% Options for packages loaded elsewhere
\PassOptionsToPackage{unicode}{hyperref}
\PassOptionsToPackage{hyphens}{url}
%
\documentclass[
]{article}
\usepackage{amsmath,amssymb}
\usepackage{iftex}
\ifPDFTeX
  \usepackage[T1]{fontenc}
  \usepackage[utf8]{inputenc}
  \usepackage{textcomp} % provide euro and other symbols
\else % if luatex or xetex
  \usepackage{unicode-math} % this also loads fontspec
  \defaultfontfeatures{Scale=MatchLowercase}
  \defaultfontfeatures[\rmfamily]{Ligatures=TeX,Scale=1}
\fi
\usepackage{lmodern}
\ifPDFTeX\else
  % xetex/luatex font selection
\fi
% Use upquote if available, for straight quotes in verbatim environments
\IfFileExists{upquote.sty}{\usepackage{upquote}}{}
\IfFileExists{microtype.sty}{% use microtype if available
  \usepackage[]{microtype}
  \UseMicrotypeSet[protrusion]{basicmath} % disable protrusion for tt fonts
}{}
\makeatletter
\@ifundefined{KOMAClassName}{% if non-KOMA class
  \IfFileExists{parskip.sty}{%
    \usepackage{parskip}
  }{% else
    \setlength{\parindent}{0pt}
    \setlength{\parskip}{6pt plus 2pt minus 1pt}}
}{% if KOMA class
  \KOMAoptions{parskip=half}}
\makeatother
\usepackage{xcolor}
\usepackage{color}
\usepackage{fancyvrb}
\newcommand{\VerbBar}{|}
\newcommand{\VERB}{\Verb[commandchars=\\\{\}]}
\DefineVerbatimEnvironment{Highlighting}{Verbatim}{commandchars=\\\{\}}
% Add ',fontsize=\small' for more characters per line
\newenvironment{Shaded}{}{}
\newcommand{\AlertTok}[1]{\textcolor[rgb]{1.00,0.00,0.00}{\textbf{#1}}}
\newcommand{\AnnotationTok}[1]{\textcolor[rgb]{0.38,0.63,0.69}{\textbf{\textit{#1}}}}
\newcommand{\AttributeTok}[1]{\textcolor[rgb]{0.49,0.56,0.16}{#1}}
\newcommand{\BaseNTok}[1]{\textcolor[rgb]{0.25,0.63,0.44}{#1}}
\newcommand{\BuiltInTok}[1]{\textcolor[rgb]{0.00,0.50,0.00}{#1}}
\newcommand{\CharTok}[1]{\textcolor[rgb]{0.25,0.44,0.63}{#1}}
\newcommand{\CommentTok}[1]{\textcolor[rgb]{0.38,0.63,0.69}{\textit{#1}}}
\newcommand{\CommentVarTok}[1]{\textcolor[rgb]{0.38,0.63,0.69}{\textbf{\textit{#1}}}}
\newcommand{\ConstantTok}[1]{\textcolor[rgb]{0.53,0.00,0.00}{#1}}
\newcommand{\ControlFlowTok}[1]{\textcolor[rgb]{0.00,0.44,0.13}{\textbf{#1}}}
\newcommand{\DataTypeTok}[1]{\textcolor[rgb]{0.56,0.13,0.00}{#1}}
\newcommand{\DecValTok}[1]{\textcolor[rgb]{0.25,0.63,0.44}{#1}}
\newcommand{\DocumentationTok}[1]{\textcolor[rgb]{0.73,0.13,0.13}{\textit{#1}}}
\newcommand{\ErrorTok}[1]{\textcolor[rgb]{1.00,0.00,0.00}{\textbf{#1}}}
\newcommand{\ExtensionTok}[1]{#1}
\newcommand{\FloatTok}[1]{\textcolor[rgb]{0.25,0.63,0.44}{#1}}
\newcommand{\FunctionTok}[1]{\textcolor[rgb]{0.02,0.16,0.49}{#1}}
\newcommand{\ImportTok}[1]{\textcolor[rgb]{0.00,0.50,0.00}{\textbf{#1}}}
\newcommand{\InformationTok}[1]{\textcolor[rgb]{0.38,0.63,0.69}{\textbf{\textit{#1}}}}
\newcommand{\KeywordTok}[1]{\textcolor[rgb]{0.00,0.44,0.13}{\textbf{#1}}}
\newcommand{\NormalTok}[1]{#1}
\newcommand{\OperatorTok}[1]{\textcolor[rgb]{0.40,0.40,0.40}{#1}}
\newcommand{\OtherTok}[1]{\textcolor[rgb]{0.00,0.44,0.13}{#1}}
\newcommand{\PreprocessorTok}[1]{\textcolor[rgb]{0.74,0.48,0.00}{#1}}
\newcommand{\RegionMarkerTok}[1]{#1}
\newcommand{\SpecialCharTok}[1]{\textcolor[rgb]{0.25,0.44,0.63}{#1}}
\newcommand{\SpecialStringTok}[1]{\textcolor[rgb]{0.73,0.40,0.53}{#1}}
\newcommand{\StringTok}[1]{\textcolor[rgb]{0.25,0.44,0.63}{#1}}
\newcommand{\VariableTok}[1]{\textcolor[rgb]{0.10,0.09,0.49}{#1}}
\newcommand{\VerbatimStringTok}[1]{\textcolor[rgb]{0.25,0.44,0.63}{#1}}
\newcommand{\WarningTok}[1]{\textcolor[rgb]{0.38,0.63,0.69}{\textbf{\textit{#1}}}}
\setlength{\emergencystretch}{3em} % prevent overfull lines
\providecommand{\tightlist}{%
  \setlength{\itemsep}{0pt}\setlength{\parskip}{0pt}}
\setcounter{secnumdepth}{-\maxdimen} % remove section numbering
\ifLuaTeX
  \usepackage{selnolig}  % disable illegal ligatures
\fi
\usepackage{bookmark}
\IfFileExists{xurl.sty}{\usepackage{xurl}}{} % add URL line breaks if available
\urlstyle{same}
\hypersetup{
  pdftitle={Implementation plan: linearized Compton IMC scheme for RadiationIMC},
  pdfauthor={OpenAI},
  hidelinks,
  pdfcreator={LaTeX via pandoc}}

\title{Implementation plan: linearized Compton IMC scheme for
RadiationIMC}
\author{OpenAI}
\date{2026-05-10}

\begin{document}
\maketitle

\section{\texorpdfstring{Implementation plan: add the linearized Compton
IMC scheme to
\texttt{RadiationIMC}}{Implementation plan: add the linearized Compton IMC scheme to RadiationIMC}}\label{implementation-plan-add-the-linearized-compton-imc-scheme-to-radiationimc}

\subsection{0. Purpose of this document}\label{purpose-of-this-document}

This document is a coding plan for modifying the current
\texttt{RadiationIMC.cpp} implementation so it can include the same
multigroup Compton linearization used by the deterministic multigroup
diffusion code, but in an IMC transport calculation.

The key implementation problem is that the Compton-modified implicit
material response produces a signed effective-scattering matrix.
Ordinary IMC can only sample nonnegative collision probabilities. The
plan below keeps the existing IMC structure for the nonnegative part,
adds physical or residual Compton transport, and represents the signed
negative part with signed residual packets, also called anti-packets.
All Compton coefficients are frozen at the beginning of the time step.
No nonlinear iterations, source iterations, or post-transport
fixed-point iterations are added.

This document is deliberately explicit. It is written so that a simple
coding agent can implement the change by following the phases and
checklists in order.

\subsection{1. Files and code areas that
matter}\label{files-and-code-areas-that-matter}

\subsubsection{\texorpdfstring{1.1 Current IMC implementation:
\texttt{RadiationIMC.cpp}}{1.1 Current IMC implementation: RadiationIMC.cpp}}\label{current-imc-implementation-radiationimc.cpp}

The uploaded current IMC implementation has these important regions:

\begin{itemize}
\tightlist
\item
  \texttt{RadiationIMC::step}, lines 40-205.

  \begin{itemize}
  \tightlist
  \item
    Currently ignores \texttt{particlesToAdd} at line 42. The Compton
    residual source will use this vector to bank newly created positive
    packets and anti-packets.
  \item
    Computes group-dependent absorption opacity at lines 65-75.
  \item
    Uses total event opacity
    \texttt{scatteringOpacity\ +\ (1\ -\ factorFleck)\ *\ absorptionOpacity}
    at lines 76-79.
  \item
    Applies continuous true absorption with
    \texttt{factorFleck\ *\ absorptionOpacity} at lines 121-145.
  \item
    Removes low-weight packets using
    \texttt{particle.weight\ \textless{}\ particle.initialWeight\ *\ 1e-3}
    at lines 146-154. This must be changed for signed packets.
  \item
    Samples elastic scattering or ordinary IMC effective scattering at
    lines 161-177.
  \end{itemize}
\item
  \texttt{RadiationIMC::postStep}, lines 207-353.

  \begin{itemize}
  \tightlist
  \item
    Normalizes time-averaged estimators at lines 210-218.
  \item
    Updates material state from \texttt{conserved} at lines 221-305.
  \item
    Reconstructs \texttt{conserved.Erad} and \texttt{conserved.Eg} from
    particle weights at lines 307-335. This will naturally work for
    signed weights, but diagnostics must be added.
  \end{itemize}
\item
  \texttt{RadiationIMC::generateParticles}, lines 355-451.

  \begin{itemize}
  \tightlist
  \item
    Creates the standard thermal IMC source energy at line 367.
  \item
    Debits material energy at lines 394-421.
  \item
    Samples thermal emission frequencies at lines 431-444.
  \item
    Sets \texttt{particle.initialWeight\ =\ particle.weight} at line
    447. This must become an absolute value for signed packets.
  \end{itemize}
\item
  \texttt{RadiationIMC::preStep}, lines 454-513.

  \begin{itemize}
  \tightlist
  \item
    Allocates \texttt{factorFleck}, \texttt{planckOpacities}, and
    estimators at lines 459-467.
  \item
    Computes the current Fleck factor at lines 472-491 using only Planck
    absorption. This must be replaced by the Compton-modified Fleck
    factor when Compton is enabled.
  \item
    Calls \texttt{generateParticles(fullDt)} at line 500.
  \end{itemize}
\item
  \texttt{generateInitialParticles}, lines 516-562.

  \begin{itemize}
  \tightlist
  \item
    Uses positive initial particle weights. No Compton source
    modification is needed here, but \texttt{initialWeight} should be
    treated as an absolute value everywhere.
  \end{itemize}
\item
  \texttt{generateSingleParticle}, lines 564-584.

  \begin{itemize}
  \tightlist
  \item
    Use this helper to create residual child packets, then override
    location, time, frequency, and signed weight.
  \end{itemize}
\end{itemize}

\subsubsection{\texorpdfstring{1.2 Deterministic Compton implementation:
\texttt{MultigroupDiffusion.cpp}}{1.2 Deterministic Compton implementation: MultigroupDiffusion.cpp}}\label{deterministic-compton-implementation-multigroupdiffusion.cpp}

The deterministic code contains the Compton coefficient logic that
should be copied or factored into shared utilities:

\begin{itemize}
\tightlist
\item
  \texttt{calculate\_fleck\_factor}, lines 1083-1135.

  \begin{itemize}
  \tightlist
  \item
    Computes \texttt{Gamma\ =\ sigma\_planck\ +\ Upsilon}.
  \item
    Computes the modified Fleck factor.
  \item
    If the Fleck factor or induced term becomes unsafe, it recomputes
    the Compton matrices with \texttt{n\ =\ 0}.
  \end{itemize}
\item
  \texttt{generate\_S\_and\_dSdUm\_matrices}, lines 1325-1451.

  \begin{itemize}
  \tightlist
  \item
    Computes the induced occupation factors \texttt{n{[}g{]}}.
  \item
    Gets \texttt{tau} and \texttt{dtau\_dUm} from
    \texttt{compton\_matrix\_gen}.
  \item
    Builds the Compton radiation operator
    \texttt{S{[}source{]}{[}target{]}} and its material-energy
    derivative \texttt{dSdUm{[}source{]}{[}target{]}}.
  \item
    Applies the same protections/cooling limiter used by the
    deterministic code.
  \end{itemize}
\item
  \texttt{calculate\_Upsilon}, lines 1453-1477.

  \begin{itemize}
  \tightlist
  \item
    Computes the contraction of \texttt{dSdUm} with the old group energy
    density.
  \end{itemize}
\item
  \texttt{get\_implicit\_compton\_contribution}, lines 1480-1519.

  \begin{itemize}
  \tightlist
  \item
    Gives the deterministic matrix contribution from direct Compton plus
    the implicit material response.
  \item
    This is the best parity check for the IMC coefficients.
  \end{itemize}
\item
  \texttt{get\_implicit\_compton\_contribution\_to\_b}, lines 1521-1550.

  \begin{itemize}
  \tightlist
  \item
    Gives the constant source correction caused by the Compton
    linearization.
  \item
    This is the best parity check for the signed source energy generated
    by the IMC code.
  \end{itemize}
\end{itemize}

\subsection{2. Mathematical target implemented by the
code}\label{mathematical-target-implemented-by-the-code}

Everything in this section is cell-local. The cell index is omitted
unless needed.

Use group indices:

\begin{itemize}
\tightlist
\item
  \texttt{h} = source group of an existing packet or existing radiation
  energy.
\item
  \texttt{g} = target group of a newly emitted, scattered, or residual
  packet.
\item
  \texttt{q} = summed group.
\end{itemize}

Old-state quantities at time \texttt{n}:

\begin{itemize}
\tightlist
\item
  \texttt{E\_h\^{}n}: old radiation energy density in group \texttt{h}.
  In the current hydro state this is usually
  \texttt{cell.Eg{[}h{]}\ *\ cell.density}.
\item
  \texttt{T\^{}n}: old material temperature.
\item
  \texttt{U\_m\^{}n\ =\ a\_r\ (T\^{}n)\^{}4}: radiation-equivalent
  material energy variable used by the Fleck linearization.
\item
  \texttt{kappa\_g}: absorption opacity in group \texttt{g}.
\item
  \texttt{b\_g}: Planck fraction in group \texttt{g}.
\item
  \texttt{kappa\_P\ =\ sum\_g\ kappa\_g\ b\_g}: Planck mean absorption
  opacity.
\item
  \texttt{S\_\{h-\textgreater{}g\}}: direct Compton radiation-energy
  operator evaluated at \texttt{T\^{}n}.
\item
  \texttt{dS\_\{h-\textgreater{}g\}/dU\_m}: derivative of the Compton
  operator with respect to the material energy variable, evaluated at
  \texttt{T\^{}n}.
\end{itemize}

The Compton contractions used by the implicit material response are

\begin{Shaded}
\begin{Highlighting}[]
\NormalTok{D\_g = \textbackslash{}sum\_h \textbackslash{}left(\textbackslash{}frac\{\textbackslash{}partial S\_\{h\textbackslash{}to g\}\}\{\textbackslash{}partial U\_m\}\textbackslash{}right)\^{}n E\_h\^{}n,}
\end{Highlighting}
\end{Shaded}

\begin{Shaded}
\begin{Highlighting}[]
\NormalTok{\textbackslash{}Upsilon = \textbackslash{}sum\_g D\_g}
\NormalTok{          = \textbackslash{}sum\_h \textbackslash{}sum\_g \textbackslash{}left(\textbackslash{}frac\{\textbackslash{}partial S\_\{h\textbackslash{}to g\}\}\{\textbackslash{}partial U\_m\}\textbackslash{}right)\^{}n E\_h\^{}n,}
\end{Highlighting}
\end{Shaded}

\begin{Shaded}
\begin{Highlighting}[]
\NormalTok{M\_g = \textbackslash{}kappa\_g b\_g + D\_g,}
\end{Highlighting}
\end{Shaded}

\begin{Shaded}
\begin{Highlighting}[]
\NormalTok{\textbackslash{}Lambda\_h = \textbackslash{}kappa\_h {-} \textbackslash{}sum\_q S\_\{h\textbackslash{}to q\}.}
\end{Highlighting}
\end{Shaded}

The Compton-modified Fleck denominator is

\begin{Shaded}
\begin{Highlighting}[]
\NormalTok{\textbackslash{}Gamma = \textbackslash{}kappa\_P + \textbackslash{}Upsilon,}
\end{Highlighting}
\end{Shaded}

\begin{Shaded}
\begin{Highlighting}[]
\NormalTok{f = \textbackslash{}frac\{1\}\{1 + \textbackslash{}beta c \textbackslash{}Delta t\textbackslash{},\textbackslash{}Gamma\},}
\end{Highlighting}
\end{Shaded}

where \texttt{beta} must be computed with the same convention already
used by the IMC Fleck factor. In the current \texttt{RadiationIMC.cpp}
this means the expression at line 479 generalizes from

\begin{Shaded}
\begin{Highlighting}[]
\FloatTok{1.0} \OperatorTok{/} \OperatorTok{(}\FloatTok{1.0} \OperatorTok{+} \DecValTok{4} \OperatorTok{*}\NormalTok{ arad }\OperatorTok{*}\NormalTok{ T}\OperatorTok{\^{}}\DecValTok{3} \OperatorTok{*}\NormalTok{ kappaP }\OperatorTok{*}\NormalTok{ c }\OperatorTok{*}\NormalTok{ dt }\OperatorTok{*}\NormalTok{ gamma }\OperatorTok{/}\NormalTok{ cv}\OperatorTok{)}
\end{Highlighting}
\end{Shaded}

to

\begin{Shaded}
\begin{Highlighting}[]
\FloatTok{1.0} \OperatorTok{/} \OperatorTok{(}\FloatTok{1.0} \OperatorTok{+} \DecValTok{4} \OperatorTok{*}\NormalTok{ arad }\OperatorTok{*}\NormalTok{ T}\OperatorTok{\^{}}\DecValTok{3} \OperatorTok{*}\NormalTok{ Gamma }\OperatorTok{*}\NormalTok{ c }\OperatorTok{*}\NormalTok{ dt }\OperatorTok{*}\NormalTok{ gamma }\OperatorTok{/}\NormalTok{ cv}\OperatorTok{)}
\end{Highlighting}
\end{Shaded}

when Compton is enabled.

The material-mediated implicit kernel is

\begin{Shaded}
\begin{Highlighting}[]
\NormalTok{K\^{}\{mat\}\_\{h\textbackslash{}to g\} = \textbackslash{}beta c\textbackslash{}Delta t\textbackslash{}, f\textbackslash{}, M\_g\textbackslash{},\textbackslash{}Lambda\_h.}
\end{Highlighting}
\end{Shaded}

This kernel is signed. It is the main source of the negative
effective-scattering problem.

The current IMC code already represents a nonnegative base absorption
effective scattering. With the Compton-modified Fleck factor this base
kernel is

\begin{Shaded}
\begin{Highlighting}[]
\NormalTok{H\^{}\{base\}\_\{h\textbackslash{}to g\} = (1{-}f)\textbackslash{},\textbackslash{}kappa\_h\textbackslash{},q\_g,}
\end{Highlighting}
\end{Shaded}

where \texttt{q\_g} must be exactly the group probability used by the
ordinary effective absorption re-emission event. For correctness use

\begin{Shaded}
\begin{Highlighting}[]
\NormalTok{q\_g = \textbackslash{}frac\{\textbackslash{}kappa\_g b\_g\}\{\textbackslash{}kappa\_P\}}
\end{Highlighting}
\end{Shaded}

when \texttt{kappa\_P\ \textgreater{}\ 0}; otherwise set all
\texttt{q\_g\ =\ 0}.

The residual kernel that is not already represented by the current IMC
effective scattering event is

\begin{Shaded}
\begin{Highlighting}[]
\NormalTok{R\_\{h\textbackslash{}to g\} = K\^{}\{mat\}\_\{h\textbackslash{}to g\} {-} H\^{}\{base\}\_\{h\textbackslash{}to g\} + S\^{}\{res\}\_\{h\textbackslash{}to g\}.}
\end{Highlighting}
\end{Shaded}

Here \texttt{Sres} is the part of direct Compton not represented by an
explicit physical Compton collision.

Recommended production implementation:

\begin{Shaded}
\begin{Highlighting}[]
\NormalTok{S\^{}\{res\}\_\{h\textbackslash{}to g\} = S\_\{h\textbackslash{}to g\} {-} S\^{}\{analog\}\_\{h\textbackslash{}to g\},}
\end{Highlighting}
\end{Shaded}

where \texttt{Sanalog} is the expectation of the explicit positive
physical Compton collision sampler.

Simplest exact implementation:

\begin{Shaded}
\begin{Highlighting}[]
\NormalTok{S\^{}\{analog\}\_\{h\textbackslash{}to g\}=0,}
\NormalTok{\textbackslash{}qquad}
\NormalTok{S\^{}\{res\}\_\{h\textbackslash{}to g\}=S\_\{h\textbackslash{}to g\}.}
\end{Highlighting}
\end{Shaded}

The simplest implementation preserves the linearized multigroup equation
exactly in expectation, but it represents direct Compton partly with
signed residual packets instead of physical group-changing collisions.
Implement the simplest version first if time is limited, then add
physical analog Compton as an optimization and physics-improving step.

Split the residual kernel into nonnegative parts:

\begin{Shaded}
\begin{Highlighting}[]
\NormalTok{R\^{}+\_\{h\textbackslash{}to g\}=\textbackslash{}max(R\_\{h\textbackslash{}to g\},0),}
\NormalTok{\textbackslash{}qquad}
\NormalTok{R\^{}{-}\_\{h\textbackslash{}to g\}=\textbackslash{}max({-}R\_\{h\textbackslash{}to g\},0).}
\end{Highlighting}
\end{Shaded}

The residual contribution is sampled by creating child packets along
each track segment:

\begin{Shaded}
\begin{Highlighting}[]
\NormalTok{R\_\{h\textbackslash{}to g\}=R\^{}+\_\{h\textbackslash{}to g\}{-}R\^{}{-}\_\{h\textbackslash{}to g\}.}
\end{Highlighting}
\end{Shaded}

A positive residual child has the same sign as its parent. A negative
residual child has the opposite sign.

\subsection{3. Source term to emit at the beginning of the
step}\label{source-term-to-emit-at-the-beginning-of-the-step}

The current code generates the standard IMC material source in
\texttt{generateParticles}. With Compton enabled, replace that source
path by a group-specific signed source.

The base source that ordinary IMC would emit is

\begin{Shaded}
\begin{Highlighting}[]
\NormalTok{B\^{}\{base\}\_g = V\textbackslash{},c\textbackslash{}Delta t\textbackslash{}, f\textbackslash{},\textbackslash{}kappa\_g b\_g\textbackslash{},U\_m\^{}n.}
\end{Highlighting}
\end{Shaded}

The deterministic Compton source correction from
\texttt{get\_implicit\_compton\_contribution\_to\_b} is

\begin{Shaded}
\begin{Highlighting}[]
\NormalTok{B\^{}\{corr\}\_g}
\NormalTok{= V\textbackslash{},c\textbackslash{}Delta t\textbackslash{},\textbackslash{}kappa\_P U\_m\^{}n}
\NormalTok{  \textbackslash{}left[}
\NormalTok{    \textbackslash{}frac\{\textbackslash{}kappa\_g b\_g\}\{\textbackslash{}kappa\_P\}}
\NormalTok{       \textbackslash{}left(1{-}(1+\textbackslash{}beta c\textbackslash{}Delta t\textbackslash{},\textbackslash{}kappa\_P)f\textbackslash{}right)}
\NormalTok{    {-} \textbackslash{}beta c\textbackslash{}Delta t\textbackslash{}, f\textbackslash{},D\_g}
\NormalTok{  \textbackslash{}right].}
\end{Highlighting}
\end{Shaded}

The total source emitted by the Compton-enabled IMC code should be

\begin{Shaded}
\begin{Highlighting}[]
\NormalTok{B\^{}\{total\}\_g = B\^{}\{base\}\_g + B\^{}\{corr\}\_g.}
\end{Highlighting}
\end{Shaded}

Equivalently, after simplification,

\begin{Shaded}
\begin{Highlighting}[]
\NormalTok{B\^{}\{total\}\_g}
\NormalTok{= V\textbackslash{},c\textbackslash{}Delta t\textbackslash{},U\_m\^{}n}
\NormalTok{  \textbackslash{}left[}
\NormalTok{      \textbackslash{}kappa\_g b\_g\textbackslash{}left(1{-}\textbackslash{}beta c\textbackslash{}Delta t\textbackslash{},\textbackslash{}kappa\_P f\textbackslash{}right)}
\NormalTok{      {-} \textbackslash{}kappa\_P\textbackslash{},\textbackslash{}beta c\textbackslash{}Delta t\textbackslash{},f\textbackslash{},D\_g}
\NormalTok{  \textbackslash{}right].}
\end{Highlighting}
\end{Shaded}

Use the unsimplified \texttt{Bbase\ +\ Bcorr} form in code because it is
easier to compare to the deterministic implementation.

Important: when Compton is enabled, do not also call the old scalar
thermal-source logic from line 367. Either generate \texttt{Btotal\_g}
directly, or keep the old source and generate only \texttt{Bcorr\_g}.
The recommended implementation is to generate \texttt{Btotal\_g}
directly in a new function \texttt{generateComptonParticles(fullDt)}.

Split the source by sign:

\begin{Shaded}
\begin{Highlighting}[]
\NormalTok{B\^{}\{total,+\}\_g = \textbackslash{}max(B\^{}\{total\}\_g,0),}
\NormalTok{\textbackslash{}qquad}
\NormalTok{B\^{}\{total,{-}\}\_g = \textbackslash{}max({-}B\^{}\{total\}\_g,0).}
\end{Highlighting}
\end{Shaded}

Emit positive source packets for \texttt{Btotal,+}. Emit anti-packets or
negative signed packets for \texttt{Btotal,-}. Material energy is
updated by the negative of the signed emitted radiation energy.
Therefore a negative radiation source packet increases the material
energy.

\subsection{4. High-level code
architecture}\label{high-level-code-architecture}

Add Compton support in small phases. Keep \texttt{withCompton=false}
behavior as close to bitwise identical as possible.

\subsubsection{Phase A: signed packet support with no
Compton}\label{phase-a-signed-packet-support-with-no-compton}

Goal: allow the transport code to safely carry negative packet weights,
but do not generate any negative packets yet.

Required changes:

\begin{enumerate}
\def\labelenumi{\arabic{enumi}.}
\tightlist
\item
  Keep using \texttt{particle.weight} as the signed packet energy.
\item
  Treat \texttt{particle.initialWeight} as an absolute reference weight.
\item
  Replace all packet-removal checks based on raw \texttt{weight} by
  checks based on \texttt{std::abs(weight)}.
\item
  Make all material and momentum updates use the signed packet weight.
  Most already do this automatically.
\item
  Make all estimators and final tallies signed. Most already do this
  automatically.
\item
  Add diagnostics for negative radiation energy after tallying, but do
  not silently clip unless the negative value is a tiny roundoff
  tolerance and the code already has an explicit conservation
  correction.
\end{enumerate}

Specific current-code edits:

\begin{itemize}
\tightlist
\item
  In \texttt{step}, replace line 146 logic:
\end{itemize}

\begin{Shaded}
\begin{Highlighting}[]
\ControlFlowTok{if}\OperatorTok{(}\NormalTok{particle}\OperatorTok{.}\NormalTok{weight }\OperatorTok{\textless{}}\NormalTok{ particle}\OperatorTok{.}\NormalTok{initialWeight }\OperatorTok{*} \FloatTok{1e{-}3}\OperatorTok{)}
\end{Highlighting}
\end{Shaded}

with

\begin{Shaded}
\begin{Highlighting}[]
\ControlFlowTok{if}\OperatorTok{(}\BuiltInTok{std::}\NormalTok{abs}\OperatorTok{(}\NormalTok{particle}\OperatorTok{.}\NormalTok{weight}\OperatorTok{)} \OperatorTok{\textless{}}\NormalTok{ particle}\OperatorTok{.}\NormalTok{initialWeight }\OperatorTok{*} \FloatTok{1e{-}3}\OperatorTok{)}
\end{Highlighting}
\end{Shaded}

and ensure \texttt{particle.initialWeight\ \textgreater{}\ 0} for every
particle.

\begin{itemize}
\tightlist
\item
  In \texttt{generateParticles}, line 447:
\end{itemize}

\begin{Shaded}
\begin{Highlighting}[]
\NormalTok{particle}\OperatorTok{.}\NormalTok{initialWeight }\OperatorTok{=} \BuiltInTok{std::}\NormalTok{abs}\OperatorTok{(}\NormalTok{particle}\OperatorTok{.}\NormalTok{weight}\OperatorTok{);}
\end{Highlighting}
\end{Shaded}

\begin{itemize}
\tightlist
\item
  In \texttt{generateInitialParticles}, line 557:
\end{itemize}

\begin{Shaded}
\begin{Highlighting}[]
\NormalTok{p}\OperatorTok{.}\NormalTok{initialWeight }\OperatorTok{=} \BuiltInTok{std::}\NormalTok{abs}\OperatorTok{(}\NormalTok{weightPerPhoton}\OperatorTok{);}
\end{Highlighting}
\end{Shaded}

\begin{itemize}
\tightlist
\item
  For any new residual child packet:
\end{itemize}

\begin{Shaded}
\begin{Highlighting}[]
\NormalTok{child}\OperatorTok{.}\NormalTok{initialWeight }\OperatorTok{=} \BuiltInTok{std::}\NormalTok{abs}\OperatorTok{(}\NormalTok{child}\OperatorTok{.}\NormalTok{weight}\OperatorTok{);}
\end{Highlighting}
\end{Shaded}

\begin{itemize}
\tightlist
\item
  In any future roulette function, roulette on
  \texttt{std::abs(weight)}, not on \texttt{weight}.
\end{itemize}

Do not add a separate sign field unless the rest of the Monte Carlo
framework forbids negative \texttt{weight}. If a separate sign field is
required, define the signed energy as
\texttt{signedWeight\ =\ sign\ *\ weightAbs}, and replace every
energy/momentum/tally expression by \texttt{signedWeight}.

\subsubsection{Phase B: parameter and data-structure
additions}\label{phase-b-parameter-and-data-structure-additions}

Add parameters to \texttt{RadiationIMCParameters} in the header. The
exact header was not uploaded, so locate \texttt{RadiationIMCParameters}
and add fields similar to these:

\begin{Shaded}
\begin{Highlighting}[]
\DataTypeTok{bool}\NormalTok{ withCompton }\OperatorTok{=} \KeywordTok{false}\OperatorTok{;}
\DataTypeTok{bool}\NormalTok{ comptonUseInduced }\OperatorTok{=} \KeywordTok{true}\OperatorTok{;}
\DataTypeTok{bool}\NormalTok{ comptonAllowNZeroFallback }\OperatorTok{=} \KeywordTok{true}\OperatorTok{;}
\DataTypeTok{bool}\NormalTok{ comptonUseAnalogDirect }\OperatorTok{=} \KeywordTok{false}\OperatorTok{;}  \CommentTok{// start false; turn true after tests}
\DataTypeTok{bool}\NormalTok{ comptonDisableRandomWalk }\OperatorTok{=} \KeywordTok{true}\OperatorTok{;}
\DataTypeTok{double}\NormalTok{ comptonOpticalDepthTurnOff }\OperatorTok{=} \BuiltInTok{std::}\NormalTok{numeric\_limits}\OperatorTok{\textless{}}\DataTypeTok{double}\OperatorTok{\textgreater{}::}\NormalTok{infinity}\OperatorTok{();}
\DataTypeTok{double}\NormalTok{ signedPacketRouletteFloor }\OperatorTok{=} \FloatTok{0.0}\OperatorTok{;}
\DataTypeTok{double}\NormalTok{ signedPacketRouletteFraction }\OperatorTok{=} \FloatTok{1e{-}4}\OperatorTok{;}
\DataTypeTok{size\_t}\NormalTok{ minSourcePacketsPerSignedGroup }\OperatorTok{=} \DecValTok{1}\OperatorTok{;}
\DataTypeTok{size\_t}\NormalTok{ maxSourcePacketsPerCellMultiplier }\OperatorTok{=} \DecValTok{40}\OperatorTok{;}
\end{Highlighting}
\end{Shaded}

Add corresponding members to \texttt{RadiationIMC}:

\begin{Shaded}
\begin{Highlighting}[]
\DataTypeTok{bool}\NormalTok{ withCompton}\OperatorTok{;}
\DataTypeTok{bool}\NormalTok{ comptonUseInduced}\OperatorTok{;}
\DataTypeTok{bool}\NormalTok{ comptonAllowNZeroFallback}\OperatorTok{;}
\DataTypeTok{bool}\NormalTok{ comptonUseAnalogDirect}\OperatorTok{;}
\DataTypeTok{bool}\NormalTok{ comptonDisableRandomWalk}\OperatorTok{;}
\DataTypeTok{double}\NormalTok{ comptonOpticalDepthTurnOff}\OperatorTok{;}
\DataTypeTok{double}\NormalTok{ signedPacketRouletteFloor}\OperatorTok{;}
\DataTypeTok{double}\NormalTok{ signedPacketRouletteFraction}\OperatorTok{;}
\DataTypeTok{size\_t}\NormalTok{ minSourcePacketsPerSignedGroup}\OperatorTok{;}
\DataTypeTok{size\_t}\NormalTok{ maxSourcePacketsPerCellMultiplier}\OperatorTok{;}
\end{Highlighting}
\end{Shaded}

Add them to the constructor initializer list and to
\texttt{operator\textless{}\textless{}}.

Create a cell-local data structure. Put it in the header as a private
nested struct or in an anonymous namespace if only used in the
\texttt{.cpp}.

\begin{Shaded}
\begin{Highlighting}[]
\KeywordTok{struct}\NormalTok{ ComptonCellData}
\OperatorTok{\{}
    \DataTypeTok{bool}\NormalTok{ active }\OperatorTok{=} \KeywordTok{false}\OperatorTok{;}
    \DataTypeTok{bool}\NormalTok{ useNZero }\OperatorTok{=} \KeywordTok{false}\OperatorTok{;}

    \DataTypeTok{double}\NormalTok{ volume }\OperatorTok{=} \FloatTok{0.0}\OperatorTok{;}
    \DataTypeTok{double}\NormalTok{ temperature }\OperatorTok{=} \FloatTok{0.0}\OperatorTok{;}
    \DataTypeTok{double}\NormalTok{ Um }\OperatorTok{=} \FloatTok{0.0}\OperatorTok{;}
    \DataTypeTok{double}\NormalTok{ cv }\OperatorTok{=} \FloatTok{0.0}\OperatorTok{;}
    \DataTypeTok{double}\NormalTok{ beta }\OperatorTok{=} \FloatTok{0.0}\OperatorTok{;}
    \DataTypeTok{double}\NormalTok{ kappaP }\OperatorTok{=} \FloatTok{0.0}\OperatorTok{;}
    \DataTypeTok{double}\NormalTok{ Upsilon }\OperatorTok{=} \FloatTok{0.0}\OperatorTok{;}
    \DataTypeTok{double}\NormalTok{ Gamma }\OperatorTok{=} \FloatTok{0.0}\OperatorTok{;}
    \DataTypeTok{double}\NormalTok{ fleck }\OperatorTok{=} \FloatTok{1.0}\OperatorTok{;}
    \DataTypeTok{double}\NormalTok{ betaCdtF }\OperatorTok{=} \FloatTok{0.0}\OperatorTok{;}      \CommentTok{// beta * c * dt * fleck, possibly with gamma factor}

    \BuiltInTok{std::}\NormalTok{array}\OperatorTok{\textless{}}\DataTypeTok{double}\OperatorTok{,}\NormalTok{ ENERGY\_GROUPS\_NUM}\OperatorTok{\textgreater{}}\NormalTok{ kappa}\OperatorTok{\{\};}
    \BuiltInTok{std::}\NormalTok{array}\OperatorTok{\textless{}}\DataTypeTok{double}\OperatorTok{,}\NormalTok{ ENERGY\_GROUPS\_NUM}\OperatorTok{\textgreater{}}\NormalTok{ b}\OperatorTok{\{\};}
    \BuiltInTok{std::}\NormalTok{array}\OperatorTok{\textless{}}\DataTypeTok{double}\OperatorTok{,}\NormalTok{ ENERGY\_GROUPS\_NUM}\OperatorTok{\textgreater{}}\NormalTok{ qBase}\OperatorTok{\{\};}
    \BuiltInTok{std::}\NormalTok{array}\OperatorTok{\textless{}}\DataTypeTok{double}\OperatorTok{,}\NormalTok{ ENERGY\_GROUPS\_NUM}\OperatorTok{\textgreater{}}\NormalTok{ Eold}\OperatorTok{\{\};}
    \BuiltInTok{std::}\NormalTok{array}\OperatorTok{\textless{}}\DataTypeTok{double}\OperatorTok{,}\NormalTok{ ENERGY\_GROUPS\_NUM}\OperatorTok{\textgreater{}}\NormalTok{ D}\OperatorTok{\{\};}
    \BuiltInTok{std::}\NormalTok{array}\OperatorTok{\textless{}}\DataTypeTok{double}\OperatorTok{,}\NormalTok{ ENERGY\_GROUPS\_NUM}\OperatorTok{\textgreater{}}\NormalTok{ M}\OperatorTok{\{\};}
    \BuiltInTok{std::}\NormalTok{array}\OperatorTok{\textless{}}\DataTypeTok{double}\OperatorTok{,}\NormalTok{ ENERGY\_GROUPS\_NUM}\OperatorTok{\textgreater{}}\NormalTok{ Lambda}\OperatorTok{\{\};}
    \BuiltInTok{std::}\NormalTok{array}\OperatorTok{\textless{}}\DataTypeTok{double}\OperatorTok{,}\NormalTok{ ENERGY\_GROUPS\_NUM}\OperatorTok{\textgreater{}}\NormalTok{ rowS}\OperatorTok{\{\};}

    \BuiltInTok{std::}\NormalTok{array}\OperatorTok{\textless{}}\DataTypeTok{double}\OperatorTok{,}\NormalTok{ ENERGY\_GROUPS\_NUM}\OperatorTok{\textgreater{}}\NormalTok{ Bbase}\OperatorTok{\{\};}
    \BuiltInTok{std::}\NormalTok{array}\OperatorTok{\textless{}}\DataTypeTok{double}\OperatorTok{,}\NormalTok{ ENERGY\_GROUPS\_NUM}\OperatorTok{\textgreater{}}\NormalTok{ Bcorr}\OperatorTok{\{\};}
    \BuiltInTok{std::}\NormalTok{array}\OperatorTok{\textless{}}\DataTypeTok{double}\OperatorTok{,}\NormalTok{ ENERGY\_GROUPS\_NUM}\OperatorTok{\textgreater{}}\NormalTok{ Btotal}\OperatorTok{\{\};}

    \BuiltInTok{std::}\NormalTok{array}\OperatorTok{\textless{}}\BuiltInTok{std::}\NormalTok{array}\OperatorTok{\textless{}}\DataTypeTok{double}\OperatorTok{,}\NormalTok{ ENERGY\_GROUPS\_NUM}\OperatorTok{\textgreater{},}\NormalTok{ ENERGY\_GROUPS\_NUM}\OperatorTok{\textgreater{}}\NormalTok{ tau}\OperatorTok{\{\};}
    \BuiltInTok{std::}\NormalTok{array}\OperatorTok{\textless{}}\BuiltInTok{std::}\NormalTok{array}\OperatorTok{\textless{}}\DataTypeTok{double}\OperatorTok{,}\NormalTok{ ENERGY\_GROUPS\_NUM}\OperatorTok{\textgreater{},}\NormalTok{ ENERGY\_GROUPS\_NUM}\OperatorTok{\textgreater{}}\NormalTok{ dtau\_dUm}\OperatorTok{\{\};}
    \BuiltInTok{std::}\NormalTok{array}\OperatorTok{\textless{}}\BuiltInTok{std::}\NormalTok{array}\OperatorTok{\textless{}}\DataTypeTok{double}\OperatorTok{,}\NormalTok{ ENERGY\_GROUPS\_NUM}\OperatorTok{\textgreater{},}\NormalTok{ ENERGY\_GROUPS\_NUM}\OperatorTok{\textgreater{}}\NormalTok{ S}\OperatorTok{\{\};}
    \BuiltInTok{std::}\NormalTok{array}\OperatorTok{\textless{}}\BuiltInTok{std::}\NormalTok{array}\OperatorTok{\textless{}}\DataTypeTok{double}\OperatorTok{,}\NormalTok{ ENERGY\_GROUPS\_NUM}\OperatorTok{\textgreater{},}\NormalTok{ ENERGY\_GROUPS\_NUM}\OperatorTok{\textgreater{}}\NormalTok{ dSdUm}\OperatorTok{\{\};}
    \BuiltInTok{std::}\NormalTok{array}\OperatorTok{\textless{}}\BuiltInTok{std::}\NormalTok{array}\OperatorTok{\textless{}}\DataTypeTok{double}\OperatorTok{,}\NormalTok{ ENERGY\_GROUPS\_NUM}\OperatorTok{\textgreater{},}\NormalTok{ ENERGY\_GROUPS\_NUM}\OperatorTok{\textgreater{}}\NormalTok{ Kmat}\OperatorTok{\{\};}
    \BuiltInTok{std::}\NormalTok{array}\OperatorTok{\textless{}}\BuiltInTok{std::}\NormalTok{array}\OperatorTok{\textless{}}\DataTypeTok{double}\OperatorTok{,}\NormalTok{ ENERGY\_GROUPS\_NUM}\OperatorTok{\textgreater{},}\NormalTok{ ENERGY\_GROUPS\_NUM}\OperatorTok{\textgreater{}}\NormalTok{ Hbase}\OperatorTok{\{\};}
    \BuiltInTok{std::}\NormalTok{array}\OperatorTok{\textless{}}\BuiltInTok{std::}\NormalTok{array}\OperatorTok{\textless{}}\DataTypeTok{double}\OperatorTok{,}\NormalTok{ ENERGY\_GROUPS\_NUM}\OperatorTok{\textgreater{},}\NormalTok{ ENERGY\_GROUPS\_NUM}\OperatorTok{\textgreater{}}\NormalTok{ Sanalog}\OperatorTok{\{\};}
    \BuiltInTok{std::}\NormalTok{array}\OperatorTok{\textless{}}\BuiltInTok{std::}\NormalTok{array}\OperatorTok{\textless{}}\DataTypeTok{double}\OperatorTok{,}\NormalTok{ ENERGY\_GROUPS\_NUM}\OperatorTok{\textgreater{},}\NormalTok{ ENERGY\_GROUPS\_NUM}\OperatorTok{\textgreater{}}\NormalTok{ Sres}\OperatorTok{\{\};}
    \BuiltInTok{std::}\NormalTok{array}\OperatorTok{\textless{}}\BuiltInTok{std::}\NormalTok{array}\OperatorTok{\textless{}}\DataTypeTok{double}\OperatorTok{,}\NormalTok{ ENERGY\_GROUPS\_NUM}\OperatorTok{\textgreater{},}\NormalTok{ ENERGY\_GROUPS\_NUM}\OperatorTok{\textgreater{}}\NormalTok{ R}\OperatorTok{\{\};}
    \BuiltInTok{std::}\NormalTok{array}\OperatorTok{\textless{}}\BuiltInTok{std::}\NormalTok{array}\OperatorTok{\textless{}}\DataTypeTok{double}\OperatorTok{,}\NormalTok{ ENERGY\_GROUPS\_NUM}\OperatorTok{\textgreater{},}\NormalTok{ ENERGY\_GROUPS\_NUM}\OperatorTok{\textgreater{}}\NormalTok{ Rpos}\OperatorTok{\{\};}
    \BuiltInTok{std::}\NormalTok{array}\OperatorTok{\textless{}}\BuiltInTok{std::}\NormalTok{array}\OperatorTok{\textless{}}\DataTypeTok{double}\OperatorTok{,}\NormalTok{ ENERGY\_GROUPS\_NUM}\OperatorTok{\textgreater{},}\NormalTok{ ENERGY\_GROUPS\_NUM}\OperatorTok{\textgreater{}}\NormalTok{ Rneg}\OperatorTok{\{\};}

    \BuiltInTok{std::}\NormalTok{array}\OperatorTok{\textless{}}\DataTypeTok{double}\OperatorTok{,}\NormalTok{ ENERGY\_GROUPS\_NUM}\OperatorTok{\textgreater{}}\NormalTok{ RposRow}\OperatorTok{\{\};}
    \BuiltInTok{std::}\NormalTok{array}\OperatorTok{\textless{}}\DataTypeTok{double}\OperatorTok{,}\NormalTok{ ENERGY\_GROUPS\_NUM}\OperatorTok{\textgreater{}}\NormalTok{ RnegRow}\OperatorTok{\{\};}
    \BuiltInTok{std::}\NormalTok{array}\OperatorTok{\textless{}}\DataTypeTok{double}\OperatorTok{,}\NormalTok{ ENERGY\_GROUPS\_NUM}\OperatorTok{\textgreater{}}\NormalTok{ directOutRate}\OperatorTok{\{\};}

    \BuiltInTok{std::}\NormalTok{array}\OperatorTok{\textless{}}\BuiltInTok{std::}\NormalTok{array}\OperatorTok{\textless{}}\DataTypeTok{double}\OperatorTok{,}\NormalTok{ ENERGY\_GROUPS\_NUM }\OperatorTok{+} \DecValTok{1}\OperatorTok{\textgreater{},}\NormalTok{ ENERGY\_GROUPS\_NUM}\OperatorTok{\textgreater{}}\NormalTok{ RposCdf}\OperatorTok{\{\};}
    \BuiltInTok{std::}\NormalTok{array}\OperatorTok{\textless{}}\BuiltInTok{std::}\NormalTok{array}\OperatorTok{\textless{}}\DataTypeTok{double}\OperatorTok{,}\NormalTok{ ENERGY\_GROUPS\_NUM }\OperatorTok{+} \DecValTok{1}\OperatorTok{\textgreater{},}\NormalTok{ ENERGY\_GROUPS\_NUM}\OperatorTok{\textgreater{}}\NormalTok{ RnegCdf}\OperatorTok{\{\};}
    \BuiltInTok{std::}\NormalTok{array}\OperatorTok{\textless{}}\BuiltInTok{std::}\NormalTok{array}\OperatorTok{\textless{}}\DataTypeTok{double}\OperatorTok{,}\NormalTok{ ENERGY\_GROUPS\_NUM }\OperatorTok{+} \DecValTok{1}\OperatorTok{\textgreater{},}\NormalTok{ ENERGY\_GROUPS\_NUM}\OperatorTok{\textgreater{}}\NormalTok{ directCdf}\OperatorTok{\{\};}
    \BuiltInTok{std::}\NormalTok{array}\OperatorTok{\textless{}}\DataTypeTok{double}\OperatorTok{,}\NormalTok{ ENERGY\_GROUPS\_NUM }\OperatorTok{+} \DecValTok{1}\OperatorTok{\textgreater{}}\NormalTok{ qBaseCdf}\OperatorTok{\{\};}
\OperatorTok{\};}
\end{Highlighting}
\end{Shaded}

Add the storage:

\begin{Shaded}
\begin{Highlighting}[]
\BuiltInTok{std::}\NormalTok{vector}\OperatorTok{\textless{}}\NormalTok{ComptonCellData}\OperatorTok{\textgreater{}}\NormalTok{ comptonData}\OperatorTok{;}
\BuiltInTok{std::}\NormalTok{array}\OperatorTok{\textless{}}\DataTypeTok{double}\OperatorTok{,}\NormalTok{ ENERGY\_GROUPS\_NUM}\OperatorTok{\textgreater{}}\NormalTok{ energyGroupCenters}\OperatorTok{;}
\BuiltInTok{std::}\NormalTok{array}\OperatorTok{\textless{}}\DataTypeTok{double}\OperatorTok{,}\NormalTok{ ENERGY\_GROUPS\_NUM}\OperatorTok{\textgreater{}}\NormalTok{ energyGroupWidths}\OperatorTok{;}
\end{Highlighting}
\end{Shaded}

If \texttt{ENERGY\_GROUPS\_NUM} is not compile-time compatible with
\texttt{std::array}, use \texttt{std::vector} with fixed sizes instead.

\subsubsection{Phase C: shared helper
functions}\label{phase-c-shared-helper-functions}

Add these helpers before \texttt{RadiationIMC::step} or as private
methods.

\paragraph{C.1 Group center and width
initialization}\label{c.1-group-center-and-width-initialization}

Do this once in the constructor:

\begin{Shaded}
\begin{Highlighting}[]
\ControlFlowTok{for}\OperatorTok{(}\DataTypeTok{size\_t}\NormalTok{ g }\OperatorTok{=} \DecValTok{0}\OperatorTok{;}\NormalTok{ g }\OperatorTok{\textless{}}\NormalTok{ ENERGY\_GROUPS\_NUM}\OperatorTok{;} \OperatorTok{++}\NormalTok{g}\OperatorTok{)}
\OperatorTok{\{}
    \AttributeTok{const} \DataTypeTok{double}\NormalTok{ e0 }\OperatorTok{=}\NormalTok{ ComputationalCell3D}\OperatorTok{::}\NormalTok{energyBoundaries}\OperatorTok{[}\NormalTok{g}\OperatorTok{];}
    \AttributeTok{const} \DataTypeTok{double}\NormalTok{ e1 }\OperatorTok{=}\NormalTok{ ComputationalCell3D}\OperatorTok{::}\NormalTok{energyBoundaries}\OperatorTok{[}\NormalTok{g }\OperatorTok{+} \DecValTok{1}\OperatorTok{];}
\NormalTok{    energyGroupCenters}\OperatorTok{[}\NormalTok{g}\OperatorTok{]} \OperatorTok{=} \BuiltInTok{std::}\NormalTok{sqrt}\OperatorTok{(}\NormalTok{e0 }\OperatorTok{*}\NormalTok{ e1}\OperatorTok{);}  \CommentTok{// logarithmic center is usually better for photon groups}
\NormalTok{    energyGroupWidths}\OperatorTok{[}\NormalTok{g}\OperatorTok{]} \OperatorTok{=}\NormalTok{ e1 }\OperatorTok{{-}}\NormalTok{ e0}\OperatorTok{;}
\OperatorTok{\}}
\end{Highlighting}
\end{Shaded}

If the deterministic code uses arithmetic centers, use those instead.
The IMC and deterministic Compton matrices must use the same centers.

\paragraph{C.2 Planck group fractions}\label{c.2-planck-group-fractions}

\begin{Shaded}
\begin{Highlighting}[]
\BuiltInTok{std::}\NormalTok{array}\OperatorTok{\textless{}}\DataTypeTok{double}\OperatorTok{,}\NormalTok{ ENERGY\_GROUPS\_NUM}\OperatorTok{\textgreater{}}\NormalTok{ ComputePlanckFractions}\OperatorTok{(}\AttributeTok{const}\NormalTok{ ComputationalCell3D}\OperatorTok{\&}\NormalTok{ cell}\OperatorTok{)}
\OperatorTok{\{}
    \BuiltInTok{std::}\NormalTok{array}\OperatorTok{\textless{}}\DataTypeTok{double}\OperatorTok{,}\NormalTok{ ENERGY\_GROUPS\_NUM}\OperatorTok{\textgreater{}}\NormalTok{ b}\OperatorTok{\{\};}
    \AttributeTok{const} \DataTypeTok{double}\NormalTok{ kT }\OperatorTok{=}\NormalTok{ units}\OperatorTok{::}\NormalTok{k\_boltz }\OperatorTok{*}\NormalTok{ cell}\OperatorTok{.}\NormalTok{temperature}\OperatorTok{;}
    \DataTypeTok{double}\NormalTok{ sum }\OperatorTok{=} \FloatTok{0.0}\OperatorTok{;}
    \ControlFlowTok{for}\OperatorTok{(}\DataTypeTok{size\_t}\NormalTok{ g }\OperatorTok{=} \DecValTok{0}\OperatorTok{;}\NormalTok{ g }\OperatorTok{\textless{}}\NormalTok{ ENERGY\_GROUPS\_NUM}\OperatorTok{;} \OperatorTok{++}\NormalTok{g}\OperatorTok{)}
    \OperatorTok{\{}
        \AttributeTok{const} \DataTypeTok{double}\NormalTok{ x0 }\OperatorTok{=}\NormalTok{ ComputationalCell3D}\OperatorTok{::}\NormalTok{energyBoundaries}\OperatorTok{[}\NormalTok{g}\OperatorTok{]} \OperatorTok{/}\NormalTok{ kT}\OperatorTok{;}
        \AttributeTok{const} \DataTypeTok{double}\NormalTok{ x1 }\OperatorTok{=}\NormalTok{ ComputationalCell3D}\OperatorTok{::}\NormalTok{energyBoundaries}\OperatorTok{[}\NormalTok{g }\OperatorTok{+} \DecValTok{1}\OperatorTok{]} \OperatorTok{/}\NormalTok{ kT}\OperatorTok{;}
\NormalTok{        b}\OperatorTok{[}\NormalTok{g}\OperatorTok{]} \OperatorTok{=}\NormalTok{ planck\_integral}\OperatorTok{::}\NormalTok{planck\_integral}\OperatorTok{(}\NormalTok{x0}\OperatorTok{,}\NormalTok{ x1}\OperatorTok{);}
\NormalTok{        sum }\OperatorTok{+=}\NormalTok{ b}\OperatorTok{[}\NormalTok{g}\OperatorTok{];}
    \OperatorTok{\}}
    \CommentTok{// Do not renormalize silently. Log if sum is far from 1 because that means the group grid misses part of the Planck spectrum.}
    \ControlFlowTok{return}\NormalTok{ b}\OperatorTok{;}
\OperatorTok{\}}
\end{Highlighting}
\end{Shaded}

Include the same \texttt{planck\_integral} header used by
\texttt{MultigroupDiffusion.cpp}.

\paragraph{C.3 Safe CDF builder}\label{c.3-safe-cdf-builder}

\begin{Shaded}
\begin{Highlighting}[]
\KeywordTok{template} \OperatorTok{\textless{}}\DataTypeTok{size\_t}\NormalTok{ N}\OperatorTok{\textgreater{}}
\BuiltInTok{std::}\NormalTok{array}\OperatorTok{\textless{}}\DataTypeTok{double}\OperatorTok{,}\NormalTok{ N }\OperatorTok{+} \DecValTok{1}\OperatorTok{\textgreater{}}\NormalTok{ BuildCdf}\OperatorTok{(}\AttributeTok{const} \BuiltInTok{std::}\NormalTok{array}\OperatorTok{\textless{}}\DataTypeTok{double}\OperatorTok{,}\NormalTok{ N}\OperatorTok{\textgreater{}\&}\NormalTok{ weights}\OperatorTok{)}
\OperatorTok{\{}
    \BuiltInTok{std::}\NormalTok{array}\OperatorTok{\textless{}}\DataTypeTok{double}\OperatorTok{,}\NormalTok{ N }\OperatorTok{+} \DecValTok{1}\OperatorTok{\textgreater{}}\NormalTok{ cdf}\OperatorTok{\{\};}
\NormalTok{    cdf}\OperatorTok{[}\DecValTok{0}\OperatorTok{]} \OperatorTok{=} \FloatTok{0.0}\OperatorTok{;}
    \DataTypeTok{double}\NormalTok{ sum }\OperatorTok{=} \FloatTok{0.0}\OperatorTok{;}
    \ControlFlowTok{for}\OperatorTok{(}\DataTypeTok{size\_t}\NormalTok{ i }\OperatorTok{=} \DecValTok{0}\OperatorTok{;}\NormalTok{ i }\OperatorTok{\textless{}}\NormalTok{ N}\OperatorTok{;} \OperatorTok{++}\NormalTok{i}\OperatorTok{)}
    \OperatorTok{\{}
\NormalTok{        sum }\OperatorTok{+=} \BuiltInTok{std::}\NormalTok{max}\OperatorTok{(}\FloatTok{0.0}\OperatorTok{,}\NormalTok{ weights}\OperatorTok{[}\NormalTok{i}\OperatorTok{]);}
\NormalTok{        cdf}\OperatorTok{[}\NormalTok{i }\OperatorTok{+} \DecValTok{1}\OperatorTok{]} \OperatorTok{=}\NormalTok{ sum}\OperatorTok{;}
    \OperatorTok{\}}
    \ControlFlowTok{if}\OperatorTok{(}\NormalTok{sum }\OperatorTok{\textgreater{}} \FloatTok{0.0}\OperatorTok{)}
    \OperatorTok{\{}
        \ControlFlowTok{for}\OperatorTok{(}\DataTypeTok{size\_t}\NormalTok{ i }\OperatorTok{=} \DecValTok{1}\OperatorTok{;}\NormalTok{ i }\OperatorTok{\textless{}=}\NormalTok{ N}\OperatorTok{;} \OperatorTok{++}\NormalTok{i}\OperatorTok{)}
\NormalTok{            cdf}\OperatorTok{[}\NormalTok{i}\OperatorTok{]} \OperatorTok{/=}\NormalTok{ sum}\OperatorTok{;}
\NormalTok{        cdf}\OperatorTok{[}\NormalTok{N}\OperatorTok{]} \OperatorTok{=} \FloatTok{1.0}\OperatorTok{;}
    \OperatorTok{\}}
    \ControlFlowTok{return}\NormalTok{ cdf}\OperatorTok{;}
\OperatorTok{\}}
\end{Highlighting}
\end{Shaded}

Add a group sampler:

\begin{Shaded}
\begin{Highlighting}[]
\KeywordTok{template} \OperatorTok{\textless{}}\DataTypeTok{size\_t}\NormalTok{ N}\OperatorTok{\textgreater{}}
\DataTypeTok{size\_t}\NormalTok{ SampleFromCdf}\OperatorTok{(}\AttributeTok{const} \BuiltInTok{std::}\NormalTok{array}\OperatorTok{\textless{}}\DataTypeTok{double}\OperatorTok{,}\NormalTok{ N }\OperatorTok{+} \DecValTok{1}\OperatorTok{\textgreater{}\&}\NormalTok{ cdf}\OperatorTok{,} \DataTypeTok{double}\NormalTok{ xi}\OperatorTok{)}
\OperatorTok{\{}
    \ControlFlowTok{if}\OperatorTok{(}\NormalTok{cdf}\OperatorTok{[}\NormalTok{N}\OperatorTok{]} \OperatorTok{\textless{}=} \FloatTok{0.0}\OperatorTok{)}
        \ControlFlowTok{return} \BuiltInTok{std::}\NormalTok{numeric\_limits}\OperatorTok{\textless{}}\DataTypeTok{size\_t}\OperatorTok{\textgreater{}::}\NormalTok{max}\OperatorTok{();}
\NormalTok{    xi }\OperatorTok{=} \BuiltInTok{std::}\NormalTok{clamp}\OperatorTok{(}\NormalTok{xi}\OperatorTok{,} \FloatTok{0.0}\OperatorTok{,} \BuiltInTok{std::}\NormalTok{nextafter}\OperatorTok{(}\FloatTok{1.0}\OperatorTok{,} \FloatTok{0.0}\OperatorTok{));}
    \KeywordTok{auto}\NormalTok{ it }\OperatorTok{=} \BuiltInTok{std::}\NormalTok{upper\_bound}\OperatorTok{(}\NormalTok{cdf}\OperatorTok{.}\NormalTok{begin}\OperatorTok{(),}\NormalTok{ cdf}\OperatorTok{.}\NormalTok{end}\OperatorTok{(),}\NormalTok{ xi}\OperatorTok{);}
    \DataTypeTok{size\_t}\NormalTok{ g }\OperatorTok{=} \KeywordTok{static\_cast}\OperatorTok{\textless{}}\DataTypeTok{size\_t}\OperatorTok{\textgreater{}(}\BuiltInTok{std::}\NormalTok{distance}\OperatorTok{(}\NormalTok{cdf}\OperatorTok{.}\NormalTok{begin}\OperatorTok{(),}\NormalTok{ it}\OperatorTok{))} \OperatorTok{{-}} \DecValTok{1}\OperatorTok{;}
    \ControlFlowTok{return} \BuiltInTok{std::}\NormalTok{min}\OperatorTok{(}\NormalTok{g}\OperatorTok{,}\NormalTok{ N }\OperatorTok{{-}} \DecValTok{1}\OperatorTok{);}
\OperatorTok{\}}
\end{Highlighting}
\end{Shaded}

\paragraph{C.4 Frequency in a specific
group}\label{c.4-frequency-in-a-specific-group}

For a first implementation, put all group-specific packets at the group
center:

\begin{Shaded}
\begin{Highlighting}[]
\DataTypeTok{double}\NormalTok{ FrequencyForGroup}\OperatorTok{(}\DataTypeTok{size\_t}\NormalTok{ g}\OperatorTok{)} \AttributeTok{const}
\OperatorTok{\{}
    \DataTypeTok{double}\NormalTok{ frequency }\OperatorTok{=}\NormalTok{ energyGroupCenters}\OperatorTok{[}\NormalTok{g}\OperatorTok{];}
\NormalTok{    ClampFrequencyToBounds}\OperatorTok{(}\NormalTok{frequency}\OperatorTok{);}
    \ControlFlowTok{return}\NormalTok{ frequency}\OperatorTok{;}
\OperatorTok{\}}
\end{Highlighting}
\end{Shaded}

Later, replace this with a Planck-weighted or opacity-weighted
within-group sampler. The multigroup equations only care that the packet
remains in the selected group, so group-center sampling is acceptable
for the first implementation.

\paragraph{C.5 Integrated signed weight over a path
segment}\label{c.5-integrated-signed-weight-over-a-path-segment}

Residual source packets should be generated from a track-length
estimator. Because the existing code applies continuous absorption over
the segment, use the attenuation-integrated packet weight, not only the
end-of-segment weight.

Add:

\begin{Shaded}
\begin{Highlighting}[]
\DataTypeTok{double}\NormalTok{ IntegratedWeightOverDt}\OperatorTok{(}\DataTypeTok{double}\NormalTok{ signedWeightAtStart}\OperatorTok{,} \DataTypeTok{double}\NormalTok{ alpha}\OperatorTok{,} \DataTypeTok{double}\NormalTok{ dt}\OperatorTok{)}
\OperatorTok{\{}
    \CommentTok{// alpha is the continuous absorption rate in 1/time.}
    \CommentTok{// For the current code alpha = factorFleck[cell] * absorptionOpacity * units::clight.}
    \AttributeTok{const} \DataTypeTok{double}\NormalTok{ x }\OperatorTok{=}\NormalTok{ alpha }\OperatorTok{*}\NormalTok{ dt}\OperatorTok{;}
    \ControlFlowTok{if}\OperatorTok{(}\BuiltInTok{std::}\NormalTok{abs}\OperatorTok{(}\NormalTok{x}\OperatorTok{)} \OperatorTok{\textless{}} \FloatTok{1e{-}10}\OperatorTok{)}
    \OperatorTok{\{}
        \ControlFlowTok{return}\NormalTok{ signedWeightAtStart }\OperatorTok{*}\NormalTok{ dt }\OperatorTok{*} \OperatorTok{(}\FloatTok{1.0} \OperatorTok{{-}} \FloatTok{0.5} \OperatorTok{*}\NormalTok{ x }\OperatorTok{+}\NormalTok{ x }\OperatorTok{*}\NormalTok{ x }\OperatorTok{/} \FloatTok{6.0}\OperatorTok{);}
    \OperatorTok{\}}
    \ControlFlowTok{return}\NormalTok{ signedWeightAtStart }\OperatorTok{*} \OperatorTok{({-}}\BuiltInTok{std::}\NormalTok{expm1}\OperatorTok{({-}}\NormalTok{x}\OperatorTok{))} \OperatorTok{/}\NormalTok{ alpha}\OperatorTok{;}
\OperatorTok{\}}
\end{Highlighting}
\end{Shaded}

This returns \texttt{integral\_0\^{}dt\ W(t)\ dt}. The residual source
energy for an opacity-like kernel \texttt{R} is then

\begin{Shaded}
\begin{Highlighting}[]
\NormalTok{sourceEnergy }\OperatorTok{=}\NormalTok{ units}\OperatorTok{::}\NormalTok{clight }\OperatorTok{*}\NormalTok{ rateFrameFactor }\OperatorTok{*}\NormalTok{ IntegratedWeightOverDt}\OperatorTok{(...)} \OperatorTok{*}\NormalTok{ Rrow}\OperatorTok{;}
\end{Highlighting}
\end{Shaded}

For non-moving material, \texttt{rateFrameFactor\ =\ 1}. For the current
hydro/Doppler convention, use the same Doppler factor that is used to
convert opacity to event distance. A conservative first implementation
can set

\begin{Shaded}
\begin{Highlighting}[]
\DataTypeTok{double}\NormalTok{ rateFrameFactor }\OperatorTok{=} \OperatorTok{(}\NormalTok{withHydro }\OperatorTok{\&\&} \OperatorTok{!}\NormalTok{MMC}\OperatorTok{)} \OperatorTok{?}\NormalTok{ dopplerShift }\OperatorTok{:} \FloatTok{1.0}\OperatorTok{;}
\end{Highlighting}
\end{Shaded}

and then verify against nonrelativistic tests first.

\paragraph{C.6 Residual source birth time on a
segment}\label{c.6-residual-source-birth-time-on-a-segment}

For exact time placement, sample the birth time with probability
proportional to the attenuated parent weight along the segment.

\begin{Shaded}
\begin{Highlighting}[]
\DataTypeTok{double}\NormalTok{ SampleBirthTimeOnSegment}\OperatorTok{(}\DataTypeTok{double}\NormalTok{ alpha}\OperatorTok{,} \DataTypeTok{double}\NormalTok{ dt}\OperatorTok{,} \DataTypeTok{double}\NormalTok{ xi}\OperatorTok{)}
\OperatorTok{\{}
\NormalTok{    xi }\OperatorTok{=} \BuiltInTok{std::}\NormalTok{clamp}\OperatorTok{(}\NormalTok{xi}\OperatorTok{,} \FloatTok{0.0}\OperatorTok{,} \BuiltInTok{std::}\NormalTok{nextafter}\OperatorTok{(}\FloatTok{1.0}\OperatorTok{,} \FloatTok{0.0}\OperatorTok{));}
    \AttributeTok{const} \DataTypeTok{double}\NormalTok{ x }\OperatorTok{=}\NormalTok{ alpha }\OperatorTok{*}\NormalTok{ dt}\OperatorTok{;}
    \ControlFlowTok{if}\OperatorTok{(}\BuiltInTok{std::}\NormalTok{abs}\OperatorTok{(}\NormalTok{x}\OperatorTok{)} \OperatorTok{\textless{}} \FloatTok{1e{-}10}\OperatorTok{)}
        \ControlFlowTok{return}\NormalTok{ dt }\OperatorTok{*}\NormalTok{ xi}\OperatorTok{;}
    \CommentTok{// CDF(t) = (1 {-} exp({-}alpha t)) / (1 {-} exp({-}alpha dt))}
    \AttributeTok{const} \DataTypeTok{double}\NormalTok{ denom }\OperatorTok{=} \OperatorTok{{-}}\BuiltInTok{std::}\NormalTok{expm1}\OperatorTok{({-}}\NormalTok{x}\OperatorTok{);}
    \ControlFlowTok{return} \OperatorTok{{-}}\BuiltInTok{std::}\NormalTok{log1p}\OperatorTok{({-}}\NormalTok{xi }\OperatorTok{*}\NormalTok{ denom}\OperatorTok{)} \OperatorTok{/}\NormalTok{ alpha}\OperatorTok{;}
\OperatorTok{\}}
\end{Highlighting}
\end{Shaded}

If this is too much for the first coding pass, place residual children
at the end of the segment. That is less accurate in time but still often
first-order consistent. The recommended production version should sample
the birth time.

\paragraph{C.7 Signed child packet
creation}\label{c.7-signed-child-packet-creation}

Add a helper:

\begin{Shaded}
\begin{Highlighting}[]
\DataTypeTok{void}\NormalTok{ RadiationIMC}\OperatorTok{::}\NormalTok{BankSignedChildPacket}\OperatorTok{(}
    \DataTypeTok{size\_t}\NormalTok{ cellIndex}\OperatorTok{,}
    \AttributeTok{const}\NormalTok{ ComputationalCell3D}\OperatorTok{\&}\NormalTok{ cell}\OperatorTok{,}
    \AttributeTok{const}\NormalTok{ Vector3D}\OperatorTok{\&}\NormalTok{ location}\OperatorTok{,}
    \AttributeTok{const}\NormalTok{ Vector3D}\OperatorTok{\&}\NormalTok{ parentVelocity}\OperatorTok{,}
    \DataTypeTok{double}\NormalTok{ timeLeft}\OperatorTok{,}
    \DataTypeTok{size\_t}\NormalTok{ targetGroup}\OperatorTok{,}
    \DataTypeTok{double}\NormalTok{ signedEnergy}\OperatorTok{,}
    \BuiltInTok{std::}\NormalTok{vector}\OperatorTok{\textless{}}\NormalTok{Particle}\OperatorTok{\textgreater{}\&}\NormalTok{ particlesToAdd}\OperatorTok{)}
\OperatorTok{\{}
    \ControlFlowTok{if}\OperatorTok{(}\NormalTok{signedEnergy }\OperatorTok{==} \FloatTok{0.0}\OperatorTok{)}
        \ControlFlowTok{return}\OperatorTok{;}

    \CommentTok{// Optional unbiased roulette for very small children.}
    \AttributeTok{const} \DataTypeTok{double}\NormalTok{ absE }\OperatorTok{=} \BuiltInTok{std::}\NormalTok{abs}\OperatorTok{(}\NormalTok{signedEnergy}\OperatorTok{);}
    \DataTypeTok{double}\NormalTok{ floor }\OperatorTok{=}\NormalTok{ signedPacketRouletteFloor}\OperatorTok{;}
    \ControlFlowTok{if}\OperatorTok{(}\NormalTok{floor }\OperatorTok{\textless{}=} \FloatTok{0.0}\OperatorTok{)}
    \OperatorTok{\{}
        \CommentTok{// A reasonable default: a fraction of the source packet scale.}
        \CommentTok{// Keep this simple initially; set to 0 to disable roulette until tests pass.}
\NormalTok{        floor }\OperatorTok{=} \FloatTok{0.0}\OperatorTok{;}
    \OperatorTok{\}}
    \ControlFlowTok{if}\OperatorTok{(}\NormalTok{floor }\OperatorTok{\textgreater{}} \FloatTok{0.0} \OperatorTok{\&\&}\NormalTok{ absE }\OperatorTok{\textless{}}\NormalTok{ floor}\OperatorTok{)}
    \OperatorTok{\{}
        \AttributeTok{const} \DataTypeTok{double}\NormalTok{ p }\OperatorTok{=}\NormalTok{ absE }\OperatorTok{/}\NormalTok{ floor}\OperatorTok{;}
        \ControlFlowTok{if}\OperatorTok{(}\KeywordTok{this}\OperatorTok{{-}\textgreater{}}\NormalTok{dist}\OperatorTok{(}\KeywordTok{this}\OperatorTok{{-}\textgreater{}}\NormalTok{re}\OperatorTok{)} \OperatorTok{\textgreater{}}\NormalTok{ p}\OperatorTok{)}
            \ControlFlowTok{return}\OperatorTok{;}
\NormalTok{        signedEnergy }\OperatorTok{=} \BuiltInTok{std::}\NormalTok{copysign}\OperatorTok{(}\NormalTok{floor}\OperatorTok{,}\NormalTok{ signedEnergy}\OperatorTok{);}
    \OperatorTok{\}}

\NormalTok{    Particle child }\OperatorTok{=} \KeywordTok{this}\OperatorTok{{-}\textgreater{}}\NormalTok{generateSingleParticle}\OperatorTok{(}\NormalTok{cellIndex}\OperatorTok{,}\NormalTok{ cell}\OperatorTok{);}
\NormalTok{    child}\OperatorTok{.}\NormalTok{location }\OperatorTok{=}\NormalTok{ location}\OperatorTok{;}
\NormalTok{    child}\OperatorTok{.}\NormalTok{cellIndex }\OperatorTok{=}\NormalTok{ cellIndex}\OperatorTok{;}
\NormalTok{    child}\OperatorTok{.}\NormalTok{cellID }\OperatorTok{=}\NormalTok{ cell}\OperatorTok{.}\NormalTok{ID}\OperatorTok{;}
\NormalTok{    child}\OperatorTok{.}\NormalTok{timeLeft }\OperatorTok{=} \BuiltInTok{std::}\NormalTok{max}\OperatorTok{(}\FloatTok{0.0}\OperatorTok{,}\NormalTok{ timeLeft}\OperatorTok{);}
\NormalTok{    child}\OperatorTok{.}\NormalTok{frequency }\OperatorTok{=}\NormalTok{ FrequencyForGroup}\OperatorTok{(}\NormalTok{targetGroup}\OperatorTok{);}
\NormalTok{    child}\OperatorTok{.}\NormalTok{weight }\OperatorTok{=}\NormalTok{ signedEnergy}\OperatorTok{;}
\NormalTok{    child}\OperatorTok{.}\NormalTok{initialWeight }\OperatorTok{=} \BuiltInTok{std::}\NormalTok{abs}\OperatorTok{(}\NormalTok{signedEnergy}\OperatorTok{);}

    \CommentTok{// Material{-}mediated residual emission is isotropic in the comoving material frame.}
\NormalTok{    child}\OperatorTok{.}\NormalTok{velocity }\OperatorTok{=} \KeywordTok{this}\OperatorTok{{-}\textgreater{}}\NormalTok{opacity}\OperatorTok{{-}\textgreater{}}\NormalTok{getRandomVelocity}\OperatorTok{(}\NormalTok{cell}\OperatorTok{);}
    \ControlFlowTok{if}\OperatorTok{(}\KeywordTok{this}\OperatorTok{{-}\textgreater{}}\NormalTok{withHydro }\OperatorTok{\&\&} \OperatorTok{!}\KeywordTok{this}\OperatorTok{{-}\textgreater{}}\NormalTok{MMC}\OperatorTok{)}
    \OperatorTok{\{}
\NormalTok{        LorentzTransformation}\OperatorTok{(}\NormalTok{child}\OperatorTok{,} \OperatorTok{{-}}\DecValTok{1} \OperatorTok{*}\NormalTok{ cell}\OperatorTok{.}\NormalTok{velocity}\OperatorTok{);}
\NormalTok{        ClampFrequencyToBounds}\OperatorTok{(}\NormalTok{child}\OperatorTok{.}\NormalTok{frequency}\OperatorTok{);}
    \OperatorTok{\}}

\NormalTok{    particlesToAdd}\OperatorTok{.}\NormalTok{push\_back}\OperatorTok{(}\NormalTok{child}\OperatorTok{);}

    \ControlFlowTok{if}\OperatorTok{(!}\KeywordTok{this}\OperatorTok{{-}\textgreater{}}\NormalTok{noHydroFeedback}\OperatorTok{)}
    \OperatorTok{\{}
        \CommentTok{// Radiation source energy came from the material if signedEnergy \textgreater{} 0.}
        \CommentTok{// If signedEnergy \textless{} 0, an anti{-}packet was emitted and the material gains energy.}
        \KeywordTok{this}\OperatorTok{{-}\textgreater{}}\NormalTok{conserved}\OperatorTok{[}\NormalTok{cellIndex}\OperatorTok{].}\NormalTok{internal\_energy }\OperatorTok{{-}=}\NormalTok{ signedEnergy}\OperatorTok{;}
        \ControlFlowTok{if}\OperatorTok{(}\KeywordTok{this}\OperatorTok{{-}\textgreater{}}\NormalTok{withHydro}\OperatorTok{)}
        \OperatorTok{\{}
            \CommentTok{// Follow the existing source convention in generateParticles.}
            \CommentTok{// For isotropic comoving emission the momentum source is the material bulk momentum, not child direction.}
            \AttributeTok{const} \DataTypeTok{double}\NormalTok{ gamma }\OperatorTok{=} \OperatorTok{(}\KeywordTok{this}\OperatorTok{{-}\textgreater{}}\NormalTok{withHydro }\OperatorTok{\&\&} \OperatorTok{!}\KeywordTok{this}\OperatorTok{{-}\textgreater{}}\NormalTok{MMC}\OperatorTok{)}
                \OperatorTok{?} \FloatTok{1.0} \OperatorTok{/} \BuiltInTok{std::}\NormalTok{sqrt}\OperatorTok{(}\FloatTok{1.0} \OperatorTok{{-}}\NormalTok{ ScalarProd}\OperatorTok{(}\NormalTok{cell}\OperatorTok{.}\NormalTok{velocity}\OperatorTok{,}\NormalTok{ cell}\OperatorTok{.}\NormalTok{velocity}\OperatorTok{)} \OperatorTok{*}\NormalTok{ units}\OperatorTok{::}\NormalTok{inv\_clight2}\OperatorTok{)}
                \OperatorTok{:} \FloatTok{1.0}\OperatorTok{;}
            \KeywordTok{this}\OperatorTok{{-}\textgreater{}}\NormalTok{conserved}\OperatorTok{[}\NormalTok{cellIndex}\OperatorTok{].}\NormalTok{energy }\OperatorTok{{-}=}\NormalTok{ signedEnergy }\OperatorTok{*}\NormalTok{ gamma}\OperatorTok{;}
            \ControlFlowTok{if}\OperatorTok{(!}\KeywordTok{this}\OperatorTok{{-}\textgreater{}}\NormalTok{diffusionPressureGradient}\OperatorTok{)}
            \OperatorTok{\{}
                \KeywordTok{this}\OperatorTok{{-}\textgreater{}}\NormalTok{conserved}\OperatorTok{[}\NormalTok{cellIndex}\OperatorTok{].}\NormalTok{momentum }\OperatorTok{{-}=}\NormalTok{ signedEnergy }\OperatorTok{*}\NormalTok{ cell}\OperatorTok{.}\NormalTok{velocity }\OperatorTok{*}\NormalTok{ units}\OperatorTok{::}\NormalTok{inv\_clight2 }\OperatorTok{*}\NormalTok{ gamma}\OperatorTok{;}
            \OperatorTok{\}}
        \OperatorTok{\}}
    \OperatorTok{\}}
\OperatorTok{\}}
\end{Highlighting}
\end{Shaded}

This helper uses signed energy. Do not reject negative
\texttt{signedEnergy}.

\subsection{5. Pre-step coefficient
construction}\label{pre-step-coefficient-construction}

Add a new private method:

\begin{Shaded}
\begin{Highlighting}[]
\DataTypeTok{void}\NormalTok{ RadiationIMC}\OperatorTok{::}\NormalTok{precomputeComptonData}\OperatorTok{(}\DataTypeTok{double}\NormalTok{ fullDt}\OperatorTok{);}
\end{Highlighting}
\end{Shaded}

Call it from \texttt{preStep} after allocating \texttt{factorFleck},
\texttt{planckOpacities}, \texttt{Erad\_time\_avg}, and
\texttt{Eg\_time\_avg}, and before \texttt{generateParticles}.

\subsubsection{\texorpdfstring{5.1 Required order in
\texttt{preStep}}{5.1 Required order in preStep}}\label{required-order-in-prestep}

Current \texttt{preStep} does:

\begin{enumerate}
\def\labelenumi{\arabic{enumi}.}
\tightlist
\item
  Allocate arrays.
\item
  Reset multigroup opacity cumulative table.
\item
  Loop cells and compute \texttt{planckOpacities} and
  \texttt{factorFleck}.
\item
  Precompute random walk data.
\item
  Generate source particles.
\item
  Generate boundary particles.
\end{enumerate}

Change it to:

\begin{Shaded}
\begin{Highlighting}[]
\ControlFlowTok{if}\OperatorTok{(}\KeywordTok{this}\OperatorTok{{-}\textgreater{}}\NormalTok{withCompton}\OperatorTok{)}
\OperatorTok{\{}
    \ControlFlowTok{if}\OperatorTok{(!}\KeywordTok{this}\OperatorTok{{-}\textgreater{}}\NormalTok{multigroupOpacity}\OperatorTok{)}
        \ControlFlowTok{throw}\NormalTok{ UniversalError}\OperatorTok{(}\StringTok{"Compton IMC requires withMultigroupOpacity"}\OperatorTok{);}

    \KeywordTok{this}\OperatorTok{{-}\textgreater{}}\NormalTok{precomputeComptonData}\OperatorTok{(}\NormalTok{fullDt}\OperatorTok{);}
\OperatorTok{\}}
\ControlFlowTok{else}
\OperatorTok{\{}
    \CommentTok{// Existing loop that computes planckOpacities and factorFleck.}
\OperatorTok{\}}

\ControlFlowTok{if}\OperatorTok{(}\KeywordTok{this}\OperatorTok{{-}\textgreater{}}\NormalTok{withRandomWalk}\OperatorTok{)}
\OperatorTok{\{}
    \ControlFlowTok{if}\OperatorTok{(}\KeywordTok{this}\OperatorTok{{-}\textgreater{}}\NormalTok{withCompton }\OperatorTok{\&\&} \KeywordTok{this}\OperatorTok{{-}\textgreater{}}\NormalTok{comptonDisableRandomWalk}\OperatorTok{)}
    \OperatorTok{\{}
        \CommentTok{// Either disable all random walk or mark Compton{-}active cells ineligible.}
    \OperatorTok{\}}
    \ControlFlowTok{else}
    \OperatorTok{\{}
        \KeywordTok{this}\OperatorTok{{-}\textgreater{}}\NormalTok{precomputeRandomWalkData}\OperatorTok{();}
    \OperatorTok{\}}
\OperatorTok{\}}

\BuiltInTok{std::}\NormalTok{vector}\OperatorTok{\textless{}}\NormalTok{Particle}\OperatorTok{\textgreater{}}\NormalTok{ newParticles }\OperatorTok{=} \KeywordTok{this}\OperatorTok{{-}\textgreater{}}\NormalTok{generateParticles}\OperatorTok{(}\NormalTok{fullDt}\OperatorTok{);}
\end{Highlighting}
\end{Shaded}

\subsubsection{\texorpdfstring{5.2 \texttt{precomputeComptonData} in
detail}{5.2 precomputeComptonData in detail}}\label{precomputecomptondata-in-detail}

Pseudo-code:

\begin{Shaded}
\begin{Highlighting}[]
\DataTypeTok{void}\NormalTok{ RadiationIMC}\OperatorTok{::}\NormalTok{precomputeComptonData}\OperatorTok{(}\DataTypeTok{double}\NormalTok{ fullDt}\OperatorTok{)}
\OperatorTok{\{}
    \AttributeTok{const} \DataTypeTok{size\_t}\NormalTok{ N }\OperatorTok{=} \KeywordTok{this}\OperatorTok{{-}\textgreater{}}\NormalTok{grid}\OperatorTok{.}\NormalTok{GetPointNo}\OperatorTok{();}
    \KeywordTok{this}\OperatorTok{{-}\textgreater{}}\NormalTok{comptonData}\OperatorTok{.}\NormalTok{assign}\OperatorTok{(}\NormalTok{N}\OperatorTok{,}\NormalTok{ ComptonCellData}\OperatorTok{\{\});}

    \ControlFlowTok{for}\OperatorTok{(}\DataTypeTok{size\_t}\NormalTok{ i }\OperatorTok{=} \DecValTok{0}\OperatorTok{;}\NormalTok{ i }\OperatorTok{\textless{}}\NormalTok{ N}\OperatorTok{;} \OperatorTok{++}\NormalTok{i}\OperatorTok{)}
    \OperatorTok{\{}
        \AttributeTok{const}\NormalTok{ ComputationalCell3D}\OperatorTok{\&}\NormalTok{ cell }\OperatorTok{=} \KeywordTok{this}\OperatorTok{{-}\textgreater{}}\NormalTok{cells}\OperatorTok{[}\NormalTok{i}\OperatorTok{];}
\NormalTok{        ComptonCellData}\OperatorTok{\&}\NormalTok{ cd }\OperatorTok{=} \KeywordTok{this}\OperatorTok{{-}\textgreater{}}\NormalTok{comptonData}\OperatorTok{[}\NormalTok{i}\OperatorTok{];}
\NormalTok{        cd}\OperatorTok{.}\NormalTok{active }\OperatorTok{=} \KeywordTok{true}\OperatorTok{;}
\NormalTok{        cd}\OperatorTok{.}\NormalTok{volume }\OperatorTok{=} \KeywordTok{this}\OperatorTok{{-}\textgreater{}}\NormalTok{grid}\OperatorTok{.}\NormalTok{GetVolume}\OperatorTok{(}\NormalTok{i}\OperatorTok{);}
\NormalTok{        cd}\OperatorTok{.}\NormalTok{temperature }\OperatorTok{=}\NormalTok{ cell}\OperatorTok{.}\NormalTok{temperature}\OperatorTok{;}
\NormalTok{        cd}\OperatorTok{.}\NormalTok{Um }\OperatorTok{=}\NormalTok{ units}\OperatorTok{::}\NormalTok{arad }\OperatorTok{*} \ExtensionTok{boost::}\NormalTok{math::pow}\OperatorTok{\textless{}}\DecValTok{4}\OperatorTok{\textgreater{}(}\NormalTok{cell}\OperatorTok{.}\NormalTok{temperature}\OperatorTok{);}

        \CommentTok{// 1. Compute group absorption opacities and Planck fractions.}
\NormalTok{        cd}\OperatorTok{.}\NormalTok{b }\OperatorTok{=}\NormalTok{ ComputePlanckFractions}\OperatorTok{(}\NormalTok{cell}\OperatorTok{);}
\NormalTok{        cd}\OperatorTok{.}\NormalTok{kappaP }\OperatorTok{=} \FloatTok{0.0}\OperatorTok{;}
        \ControlFlowTok{for}\OperatorTok{(}\DataTypeTok{size\_t}\NormalTok{ g }\OperatorTok{=} \DecValTok{0}\OperatorTok{;}\NormalTok{ g }\OperatorTok{\textless{}}\NormalTok{ ENERGY\_GROUPS\_NUM}\OperatorTok{;} \OperatorTok{++}\NormalTok{g}\OperatorTok{)}
        \OperatorTok{\{}
            \AttributeTok{const} \DataTypeTok{double}\NormalTok{ freq }\OperatorTok{=}\NormalTok{ energyGroupCenters}\OperatorTok{[}\NormalTok{g}\OperatorTok{];}
\NormalTok{            cd}\OperatorTok{.}\NormalTok{kappa}\OperatorTok{[}\NormalTok{g}\OperatorTok{]} \OperatorTok{=} \KeywordTok{this}\OperatorTok{{-}\textgreater{}}\NormalTok{opacity}\OperatorTok{{-}\textgreater{}}\NormalTok{CalcAbsorptionOpacity}\OperatorTok{(}\NormalTok{cell}\OperatorTok{,}\NormalTok{ freq}\OperatorTok{);}
\NormalTok{            cd}\OperatorTok{.}\NormalTok{kappaP }\OperatorTok{+=}\NormalTok{ cd}\OperatorTok{.}\NormalTok{kappa}\OperatorTok{[}\NormalTok{g}\OperatorTok{]} \OperatorTok{*}\NormalTok{ cd}\OperatorTok{.}\NormalTok{b}\OperatorTok{[}\NormalTok{g}\OperatorTok{];}
\NormalTok{            cd}\OperatorTok{.}\NormalTok{Eold}\OperatorTok{[}\NormalTok{g}\OperatorTok{]} \OperatorTok{=}\NormalTok{ cell}\OperatorTok{.}\NormalTok{Eg}\OperatorTok{[}\NormalTok{g}\OperatorTok{]} \OperatorTok{*}\NormalTok{ cell}\OperatorTok{.}\NormalTok{density}\OperatorTok{;}
        \OperatorTok{\}}
        \KeywordTok{this}\OperatorTok{{-}\textgreater{}}\NormalTok{planckOpacities}\OperatorTok{[}\NormalTok{i}\OperatorTok{]} \OperatorTok{=}\NormalTok{ cd}\OperatorTok{.}\NormalTok{kappaP}\OperatorTok{;}

        \CommentTok{// 2. Build qBase, the target distribution of ordinary IMC effective scattering.}
        \ControlFlowTok{if}\OperatorTok{(}\NormalTok{cd}\OperatorTok{.}\NormalTok{kappaP }\OperatorTok{\textgreater{}} \FloatTok{0.0}\OperatorTok{)}
        \OperatorTok{\{}
            \ControlFlowTok{for}\OperatorTok{(}\DataTypeTok{size\_t}\NormalTok{ g }\OperatorTok{=} \DecValTok{0}\OperatorTok{;}\NormalTok{ g }\OperatorTok{\textless{}}\NormalTok{ ENERGY\_GROUPS\_NUM}\OperatorTok{;} \OperatorTok{++}\NormalTok{g}\OperatorTok{)}
\NormalTok{                cd}\OperatorTok{.}\NormalTok{qBase}\OperatorTok{[}\NormalTok{g}\OperatorTok{]} \OperatorTok{=}\NormalTok{ cd}\OperatorTok{.}\NormalTok{kappa}\OperatorTok{[}\NormalTok{g}\OperatorTok{]} \OperatorTok{*}\NormalTok{ cd}\OperatorTok{.}\NormalTok{b}\OperatorTok{[}\NormalTok{g}\OperatorTok{]} \OperatorTok{/}\NormalTok{ cd}\OperatorTok{.}\NormalTok{kappaP}\OperatorTok{;}
        \OperatorTok{\}}
\NormalTok{        cd}\OperatorTok{.}\NormalTok{qBaseCdf }\OperatorTok{=}\NormalTok{ BuildCdf}\OperatorTok{(}\NormalTok{cd}\OperatorTok{.}\NormalTok{qBase}\OperatorTok{);}

        \CommentTok{// 3. Build S and dSdUm using the deterministic routine port.}
\NormalTok{        BuildComptonMatricesForCell}\OperatorTok{(}\NormalTok{cell}\OperatorTok{,}\NormalTok{ i}\OperatorTok{,}\NormalTok{ fullDt}\OperatorTok{,} \CommentTok{/*calculate\_n=*/}\KeywordTok{this}\OperatorTok{{-}\textgreater{}}\NormalTok{comptonUseInduced}\OperatorTok{,}\NormalTok{ cd}\OperatorTok{);}

        \CommentTok{// 4. Compute D[g], Upsilon, row sums, and Lambda.}
\NormalTok{        cd}\OperatorTok{.}\NormalTok{Upsilon }\OperatorTok{=} \FloatTok{0.0}\OperatorTok{;}
        \ControlFlowTok{for}\OperatorTok{(}\DataTypeTok{size\_t}\NormalTok{ g }\OperatorTok{=} \DecValTok{0}\OperatorTok{;}\NormalTok{ g }\OperatorTok{\textless{}}\NormalTok{ ENERGY\_GROUPS\_NUM}\OperatorTok{;} \OperatorTok{++}\NormalTok{g}\OperatorTok{)}
        \OperatorTok{\{}
\NormalTok{            cd}\OperatorTok{.}\NormalTok{D}\OperatorTok{[}\NormalTok{g}\OperatorTok{]} \OperatorTok{=} \FloatTok{0.0}\OperatorTok{;}
            \ControlFlowTok{for}\OperatorTok{(}\DataTypeTok{size\_t}\NormalTok{ h }\OperatorTok{=} \DecValTok{0}\OperatorTok{;}\NormalTok{ h }\OperatorTok{\textless{}}\NormalTok{ ENERGY\_GROUPS\_NUM}\OperatorTok{;} \OperatorTok{++}\NormalTok{h}\OperatorTok{)}
\NormalTok{                cd}\OperatorTok{.}\NormalTok{D}\OperatorTok{[}\NormalTok{g}\OperatorTok{]} \OperatorTok{+=}\NormalTok{ cd}\OperatorTok{.}\NormalTok{dSdUm}\OperatorTok{[}\NormalTok{h}\OperatorTok{][}\NormalTok{g}\OperatorTok{]} \OperatorTok{*}\NormalTok{ cd}\OperatorTok{.}\NormalTok{Eold}\OperatorTok{[}\NormalTok{h}\OperatorTok{];}
\NormalTok{            cd}\OperatorTok{.}\NormalTok{Upsilon }\OperatorTok{+=}\NormalTok{ cd}\OperatorTok{.}\NormalTok{D}\OperatorTok{[}\NormalTok{g}\OperatorTok{];}
\NormalTok{            cd}\OperatorTok{.}\NormalTok{M}\OperatorTok{[}\NormalTok{g}\OperatorTok{]} \OperatorTok{=}\NormalTok{ cd}\OperatorTok{.}\NormalTok{kappa}\OperatorTok{[}\NormalTok{g}\OperatorTok{]} \OperatorTok{*}\NormalTok{ cd}\OperatorTok{.}\NormalTok{b}\OperatorTok{[}\NormalTok{g}\OperatorTok{]} \OperatorTok{+}\NormalTok{ cd}\OperatorTok{.}\NormalTok{D}\OperatorTok{[}\NormalTok{g}\OperatorTok{];}
        \OperatorTok{\}}
        \ControlFlowTok{for}\OperatorTok{(}\DataTypeTok{size\_t}\NormalTok{ h }\OperatorTok{=} \DecValTok{0}\OperatorTok{;}\NormalTok{ h }\OperatorTok{\textless{}}\NormalTok{ ENERGY\_GROUPS\_NUM}\OperatorTok{;} \OperatorTok{++}\NormalTok{h}\OperatorTok{)}
        \OperatorTok{\{}
\NormalTok{            cd}\OperatorTok{.}\NormalTok{rowS}\OperatorTok{[}\NormalTok{h}\OperatorTok{]} \OperatorTok{=} \FloatTok{0.0}\OperatorTok{;}
            \ControlFlowTok{for}\OperatorTok{(}\DataTypeTok{size\_t}\NormalTok{ q }\OperatorTok{=} \DecValTok{0}\OperatorTok{;}\NormalTok{ q }\OperatorTok{\textless{}}\NormalTok{ ENERGY\_GROUPS\_NUM}\OperatorTok{;} \OperatorTok{++}\NormalTok{q}\OperatorTok{)}
\NormalTok{                cd}\OperatorTok{.}\NormalTok{rowS}\OperatorTok{[}\NormalTok{h}\OperatorTok{]} \OperatorTok{+=}\NormalTok{ cd}\OperatorTok{.}\NormalTok{S}\OperatorTok{[}\NormalTok{h}\OperatorTok{][}\NormalTok{q}\OperatorTok{];}
\NormalTok{            cd}\OperatorTok{.}\NormalTok{Lambda}\OperatorTok{[}\NormalTok{h}\OperatorTok{]} \OperatorTok{=}\NormalTok{ cd}\OperatorTok{.}\NormalTok{kappa}\OperatorTok{[}\NormalTok{h}\OperatorTok{]} \OperatorTok{{-}}\NormalTok{ cd}\OperatorTok{.}\NormalTok{rowS}\OperatorTok{[}\NormalTok{h}\OperatorTok{];}
        \OperatorTok{\}}

        \CommentTok{// 5. Compute modified Fleck factor.}
\NormalTok{        cd}\OperatorTok{.}\NormalTok{Gamma }\OperatorTok{=}\NormalTok{ cd}\OperatorTok{.}\NormalTok{kappaP }\OperatorTok{+}\NormalTok{ cd}\OperatorTok{.}\NormalTok{Upsilon}\OperatorTok{;}
        \AttributeTok{const} \DataTypeTok{double}\NormalTok{ gammaLorentz }\OperatorTok{=} \OperatorTok{(}\KeywordTok{this}\OperatorTok{{-}\textgreater{}}\NormalTok{withHydro }\OperatorTok{\&\&} \OperatorTok{!}\KeywordTok{this}\OperatorTok{{-}\textgreater{}}\NormalTok{MMC}\OperatorTok{)}
            \OperatorTok{?} \FloatTok{1.0} \OperatorTok{/} \BuiltInTok{std::}\NormalTok{sqrt}\OperatorTok{(}\FloatTok{1.0} \OperatorTok{{-}}\NormalTok{ ScalarProd}\OperatorTok{(}\NormalTok{cell}\OperatorTok{.}\NormalTok{velocity}\OperatorTok{,}\NormalTok{ cell}\OperatorTok{.}\NormalTok{velocity}\OperatorTok{)} \OperatorTok{*}\NormalTok{ units}\OperatorTok{::}\NormalTok{inv\_clight2}\OperatorTok{)}
            \OperatorTok{:} \FloatTok{1.0}\OperatorTok{;}
\NormalTok{        cd}\OperatorTok{.}\NormalTok{cv }\OperatorTok{=} \KeywordTok{this}\OperatorTok{{-}\textgreater{}}\NormalTok{eos}\OperatorTok{{-}\textgreater{}}\NormalTok{dT2cv}\OperatorTok{(}\NormalTok{cell}\OperatorTok{.}\NormalTok{density}\OperatorTok{,}\NormalTok{ cell}\OperatorTok{.}\NormalTok{temperature}\OperatorTok{,}\NormalTok{ cell}\OperatorTok{.}\NormalTok{tracers}\OperatorTok{,}\NormalTok{ cell}\OperatorTok{.}\NormalTok{tracerNames}\OperatorTok{);}
\NormalTok{        cd}\OperatorTok{.}\NormalTok{beta }\OperatorTok{=} \FloatTok{4.0} \OperatorTok{*}\NormalTok{ units}\OperatorTok{::}\NormalTok{arad }\OperatorTok{*} \ExtensionTok{boost::}\NormalTok{math::pow}\OperatorTok{\textless{}}\DecValTok{3}\OperatorTok{\textgreater{}(}\NormalTok{cell}\OperatorTok{.}\NormalTok{temperature}\OperatorTok{)} \OperatorTok{/}\NormalTok{ cd}\OperatorTok{.}\NormalTok{cv}\OperatorTok{;}
        \AttributeTok{const} \DataTypeTok{double}\NormalTok{ cdtEff }\OperatorTok{=}\NormalTok{ units}\OperatorTok{::}\NormalTok{clight }\OperatorTok{*}\NormalTok{ fullDt }\OperatorTok{*}\NormalTok{ gammaLorentz}\OperatorTok{;}
        \AttributeTok{const} \DataTypeTok{double}\NormalTok{ denom }\OperatorTok{=} \FloatTok{1.0} \OperatorTok{+}\NormalTok{ cd}\OperatorTok{.}\NormalTok{beta }\OperatorTok{*}\NormalTok{ cdtEff }\OperatorTok{*}\NormalTok{ cd}\OperatorTok{.}\NormalTok{Gamma}\OperatorTok{;}

        \ControlFlowTok{if}\OperatorTok{(}\NormalTok{denom }\OperatorTok{\textless{}=} \FloatTok{0.0}\OperatorTok{)}
        \OperatorTok{\{}
            \CommentTok{// Try n=0 fallback, matching deterministic code.}
            \ControlFlowTok{if}\OperatorTok{(}\KeywordTok{this}\OperatorTok{{-}\textgreater{}}\NormalTok{comptonAllowNZeroFallback}\OperatorTok{)}
            \OperatorTok{\{}
\NormalTok{                BuildComptonMatricesForCell}\OperatorTok{(}\NormalTok{cell}\OperatorTok{,}\NormalTok{ i}\OperatorTok{,}\NormalTok{ fullDt}\OperatorTok{,} \CommentTok{/*calculate\_n=*/}\KeywordTok{false}\OperatorTok{,}\NormalTok{ cd}\OperatorTok{);}
\NormalTok{                cd}\OperatorTok{.}\NormalTok{useNZero }\OperatorTok{=} \KeywordTok{true}\OperatorTok{;}
                \CommentTok{// Recompute D, Upsilon, M, rowS, Lambda, Gamma, denom.}
            \OperatorTok{\}}
            \ControlFlowTok{if}\OperatorTok{(}\NormalTok{denom }\OperatorTok{\textless{}=} \FloatTok{0.0}\OperatorTok{)}
                \ControlFlowTok{throw}\NormalTok{ UniversalError}\OperatorTok{(}\StringTok{"Compton Fleck denominator is nonpositive; reduce dt or apply limiter"}\OperatorTok{);}
        \OperatorTok{\}}

\NormalTok{        cd}\OperatorTok{.}\NormalTok{fleck }\OperatorTok{=} \FloatTok{1.0} \OperatorTok{/}\NormalTok{ denom}\OperatorTok{;}
        \KeywordTok{this}\OperatorTok{{-}\textgreater{}}\NormalTok{factorFleck}\OperatorTok{[}\NormalTok{i}\OperatorTok{]} \OperatorTok{=}\NormalTok{ cd}\OperatorTok{.}\NormalTok{fleck}\OperatorTok{;}
\NormalTok{        cd}\OperatorTok{.}\NormalTok{betaCdtF }\OperatorTok{=}\NormalTok{ cd}\OperatorTok{.}\NormalTok{beta }\OperatorTok{*}\NormalTok{ cdtEff }\OperatorTok{*}\NormalTok{ cd}\OperatorTok{.}\NormalTok{fleck}\OperatorTok{;}
        \CommentTok{// For numerical consistency with the existing IMC effective scattering, you can also compute:}
        \CommentTok{// if(abs(cd.Gamma) \textgreater{} tiny) cd.betaCdtF = (1.0 {-} cd.fleck) / cd.Gamma;}

        \CommentTok{// 6. Build Kmat, Hbase, Sanalog, Sres, R, splits, and CDFs.}
\NormalTok{        BuildComptonResidualKernels}\OperatorTok{(}\NormalTok{i}\OperatorTok{,}\NormalTok{ fullDt}\OperatorTok{,}\NormalTok{ cd}\OperatorTok{);}

        \CommentTok{// 7. Build Bbase, Bcorr, and Btotal.}
\NormalTok{        BuildComptonSources}\OperatorTok{(}\NormalTok{i}\OperatorTok{,}\NormalTok{ fullDt}\OperatorTok{,}\NormalTok{ cd}\OperatorTok{);}
    \OperatorTok{\}}
\OperatorTok{\}}
\end{Highlighting}
\end{Shaded}

When implementing the \texttt{n\ =\ 0} fallback, make a helper
\texttt{RecomputeComptonContractions(cd)} to avoid duplicated code. It
should recompute \texttt{D}, \texttt{Upsilon}, \texttt{M},
\texttt{rowS}, \texttt{Lambda}, and \texttt{Gamma} after the matrices
change.

\subsubsection{\texorpdfstring{5.3 Port
\texttt{BuildComptonMatricesForCell}}{5.3 Port BuildComptonMatricesForCell}}\label{port-buildcomptonmatricesforcell}

Copy the logic from
\texttt{MultigroupDiffusion::generate\_S\_and\_dSdUm\_matrices}, lines
1325-1451, into a method available to \texttt{RadiationIMC}.

Suggested signature:

\begin{Shaded}
\begin{Highlighting}[]
\DataTypeTok{void}\NormalTok{ RadiationIMC}\OperatorTok{::}\NormalTok{BuildComptonMatricesForCell}\OperatorTok{(}
    \AttributeTok{const}\NormalTok{ ComputationalCell3D}\OperatorTok{\&}\NormalTok{ cell}\OperatorTok{,}
    \DataTypeTok{size\_t}\NormalTok{ cellIndex}\OperatorTok{,}
    \DataTypeTok{double}\NormalTok{ fullDt}\OperatorTok{,}
    \DataTypeTok{bool}\NormalTok{ calculateN}\OperatorTok{,}
\NormalTok{    ComptonCellData}\OperatorTok{\&}\NormalTok{ cd}\OperatorTok{);}
\end{Highlighting}
\end{Shaded}

Implementation details:

\begin{enumerate}
\def\labelenumi{\arabic{enumi}.}
\tightlist
\item
  Compute occupation number \texttt{n{[}g{]}} exactly as deterministic
  code does:
\end{enumerate}

\begin{Shaded}
\begin{Highlighting}[]
\AttributeTok{const} \DataTypeTok{double}\NormalTok{ fac }\OperatorTok{=} \ExtensionTok{boost::}\NormalTok{math::pow}\OperatorTok{\textless{}}\DecValTok{3}\OperatorTok{\textgreater{}(}\NormalTok{units}\OperatorTok{::}\NormalTok{clight}\OperatorTok{)} \OperatorTok{/} \OperatorTok{(}\FloatTok{8.0} \OperatorTok{*}\NormalTok{ M\_PI }\OperatorTok{*}\NormalTok{ units}\OperatorTok{::}\NormalTok{planck\_constant}\OperatorTok{);}
\ControlFlowTok{for}\OperatorTok{(}\DataTypeTok{size\_t}\NormalTok{ g }\OperatorTok{=} \DecValTok{0}\OperatorTok{;}\NormalTok{ g }\OperatorTok{\textless{}}\NormalTok{ ENERGY\_GROUPS\_NUM}\OperatorTok{;} \OperatorTok{++}\NormalTok{g}\OperatorTok{)}
\OperatorTok{\{}
    \ControlFlowTok{if}\OperatorTok{(}\NormalTok{calculateN}\OperatorTok{)}
    \OperatorTok{\{}
        \AttributeTok{const} \DataTypeTok{double}\NormalTok{ dnu }\OperatorTok{=}\NormalTok{ energyGroupWidths}\OperatorTok{[}\NormalTok{g}\OperatorTok{]} \OperatorTok{/}\NormalTok{ units}\OperatorTok{::}\NormalTok{planck\_constant}\OperatorTok{;}
        \AttributeTok{const} \DataTypeTok{double}\NormalTok{ nu }\OperatorTok{=}\NormalTok{ energyGroupCenters}\OperatorTok{[}\NormalTok{g}\OperatorTok{]} \OperatorTok{/}\NormalTok{ units}\OperatorTok{::}\NormalTok{planck\_constant}\OperatorTok{;}
        \AttributeTok{const} \DataTypeTok{double}\NormalTok{ Eg }\OperatorTok{=}\NormalTok{ cell}\OperatorTok{.}\NormalTok{Eg}\OperatorTok{[}\NormalTok{g}\OperatorTok{]} \OperatorTok{*}\NormalTok{ cell}\OperatorTok{.}\NormalTok{density}\OperatorTok{;}
\NormalTok{        n}\OperatorTok{[}\NormalTok{g}\OperatorTok{]} \OperatorTok{=} \BuiltInTok{std::}\NormalTok{min}\OperatorTok{(}\FloatTok{100.0}\OperatorTok{,}\NormalTok{ fac }\OperatorTok{*}\NormalTok{ Eg }\OperatorTok{/} \OperatorTok{(}\ExtensionTok{boost::}\NormalTok{math::pow}\OperatorTok{\textless{}}\DecValTok{3}\OperatorTok{\textgreater{}(}\NormalTok{nu}\OperatorTok{)} \OperatorTok{*}\NormalTok{ dnu}\OperatorTok{));}
    \OperatorTok{\}}
    \ControlFlowTok{else}
    \OperatorTok{\{}
\NormalTok{        n}\OperatorTok{[}\NormalTok{g}\OperatorTok{]} \OperatorTok{=} \FloatTok{0.0}\OperatorTok{;}
    \OperatorTok{\}}
\OperatorTok{\}}
\end{Highlighting}
\end{Shaded}

\begin{enumerate}
\def\labelenumi{\arabic{enumi}.}
\setcounter{enumi}{1}
\item
  Call the same \texttt{compton\_matrix\_gen.get\_tau\_matrix} used by
  deterministic code. This means \texttt{RadiationIMC} must own a
  \texttt{compton\_matrix\_gen} member initialized the same way as the
  deterministic code. Copy the constructor setup from
  \texttt{MultigroupDiffusion.cpp} lines 95-134.
\item
  Fill \texttt{cd.S} and \texttt{cd.dSdUm} exactly as lines 1359-1387
  do:
\end{enumerate}

\begin{Shaded}
\begin{Highlighting}[]
\NormalTok{Zero}\OperatorTok{(}\NormalTok{cd}\OperatorTok{.}\NormalTok{S}\OperatorTok{);}
\NormalTok{Zero}\OperatorTok{(}\NormalTok{cd}\OperatorTok{.}\NormalTok{dSdUm}\OperatorTok{);}
\ControlFlowTok{for}\OperatorTok{(}\DataTypeTok{size\_t}\NormalTok{ g }\OperatorTok{=} \DecValTok{0}\OperatorTok{;}\NormalTok{ g }\OperatorTok{\textless{}}\NormalTok{ ENERGY\_GROUPS\_NUM}\OperatorTok{;} \OperatorTok{++}\NormalTok{g}\OperatorTok{)}
\OperatorTok{\{}
    \ControlFlowTok{for}\OperatorTok{(}\DataTypeTok{size\_t}\NormalTok{ gt }\OperatorTok{=} \DecValTok{0}\OperatorTok{;}\NormalTok{ gt }\OperatorTok{\textless{}}\NormalTok{ ENERGY\_GROUPS\_NUM}\OperatorTok{;} \OperatorTok{++}\NormalTok{gt}\OperatorTok{)}
    \OperatorTok{\{}
        \ControlFlowTok{if}\OperatorTok{(}\NormalTok{g }\OperatorTok{+} \DecValTok{1} \OperatorTok{==}\NormalTok{ ENERGY\_GROUPS\_NUM }\OperatorTok{\&\&}\NormalTok{ gt }\OperatorTok{+} \DecValTok{1} \OperatorTok{==}\NormalTok{ ENERGY\_GROUPS\_NUM}\OperatorTok{)}
        \OperatorTok{\{}
\NormalTok{            cd}\OperatorTok{.}\NormalTok{S}\OperatorTok{[}\NormalTok{g}\OperatorTok{][}\NormalTok{g}\OperatorTok{]} \OperatorTok{+=} \OperatorTok{(}\NormalTok{up\_scattering\_last }\OperatorTok{{-}}\NormalTok{ down\_scattering\_last}\OperatorTok{)} \OperatorTok{*} \OperatorTok{(}\FloatTok{1.0} \OperatorTok{+}\NormalTok{ n}\OperatorTok{[}\NormalTok{g}\OperatorTok{]);}
\NormalTok{            cd}\OperatorTok{.}\NormalTok{dSdUm}\OperatorTok{[}\NormalTok{g}\OperatorTok{][}\NormalTok{g}\OperatorTok{]} \OperatorTok{+=}\NormalTok{ cd}\OperatorTok{.}\NormalTok{dtau\_dUm}\OperatorTok{[}\NormalTok{g}\OperatorTok{][}\NormalTok{g}\OperatorTok{]} \OperatorTok{*} \OperatorTok{(}\FloatTok{1.0} \OperatorTok{+}\NormalTok{ n}\OperatorTok{[}\NormalTok{g}\OperatorTok{]);}
            \ControlFlowTok{continue}\OperatorTok{;}
        \OperatorTok{\}}

        \AttributeTok{const} \DataTypeTok{double}\NormalTok{ inFactor }\OperatorTok{=}\NormalTok{ energyGroupCenters}\OperatorTok{[}\NormalTok{g}\OperatorTok{]} \OperatorTok{/}\NormalTok{ energyGroupCenters}\OperatorTok{[}\NormalTok{gt}\OperatorTok{]} \OperatorTok{*} \OperatorTok{(}\FloatTok{1.0} \OperatorTok{+}\NormalTok{ n}\OperatorTok{[}\NormalTok{g}\OperatorTok{]);}
\NormalTok{        cd}\OperatorTok{.}\NormalTok{S}\OperatorTok{[}\NormalTok{gt}\OperatorTok{][}\NormalTok{g}\OperatorTok{]} \OperatorTok{+=}\NormalTok{ cd}\OperatorTok{.}\NormalTok{tau}\OperatorTok{[}\NormalTok{gt}\OperatorTok{][}\NormalTok{g}\OperatorTok{]} \OperatorTok{*}\NormalTok{ inFactor}\OperatorTok{;}
\NormalTok{        cd}\OperatorTok{.}\NormalTok{dSdUm}\OperatorTok{[}\NormalTok{gt}\OperatorTok{][}\NormalTok{g}\OperatorTok{]} \OperatorTok{+=}\NormalTok{ cd}\OperatorTok{.}\NormalTok{dtau\_dUm}\OperatorTok{[}\NormalTok{gt}\OperatorTok{][}\NormalTok{g}\OperatorTok{]} \OperatorTok{*}\NormalTok{ inFactor}\OperatorTok{;}

        \AttributeTok{const} \DataTypeTok{double}\NormalTok{ outFactor }\OperatorTok{=} \FloatTok{1.0} \OperatorTok{+}\NormalTok{ n}\OperatorTok{[}\NormalTok{gt}\OperatorTok{];}
\NormalTok{        cd}\OperatorTok{.}\NormalTok{S}\OperatorTok{[}\NormalTok{g}\OperatorTok{][}\NormalTok{g}\OperatorTok{]} \OperatorTok{{-}=}\NormalTok{ cd}\OperatorTok{.}\NormalTok{tau}\OperatorTok{[}\NormalTok{g}\OperatorTok{][}\NormalTok{gt}\OperatorTok{]} \OperatorTok{*}\NormalTok{ outFactor}\OperatorTok{;}
\NormalTok{        cd}\OperatorTok{.}\NormalTok{dSdUm}\OperatorTok{[}\NormalTok{g}\OperatorTok{][}\NormalTok{g}\OperatorTok{]} \OperatorTok{{-}=}\NormalTok{ cd}\OperatorTok{.}\NormalTok{dtau\_dUm}\OperatorTok{[}\NormalTok{g}\OperatorTok{][}\NormalTok{gt}\OperatorTok{]} \OperatorTok{*}\NormalTok{ outFactor}\OperatorTok{;}
    \OperatorTok{\}}
\OperatorTok{\}}
\AttributeTok{const} \DataTypeTok{double}\NormalTok{ UmFactor }\OperatorTok{=} \FloatTok{1.0} \OperatorTok{/} \OperatorTok{(}\FloatTok{4.0} \OperatorTok{*}\NormalTok{ units}\OperatorTok{::}\NormalTok{arad }\OperatorTok{*} \ExtensionTok{boost::}\NormalTok{math::pow}\OperatorTok{\textless{}}\DecValTok{3}\OperatorTok{\textgreater{}(}\NormalTok{cell}\OperatorTok{.}\NormalTok{temperature}\OperatorTok{));}
\ControlFlowTok{for}\OperatorTok{(}\DataTypeTok{size\_t}\NormalTok{ h }\OperatorTok{=} \DecValTok{0}\OperatorTok{;}\NormalTok{ h }\OperatorTok{\textless{}}\NormalTok{ ENERGY\_GROUPS\_NUM}\OperatorTok{;} \OperatorTok{++}\NormalTok{h}\OperatorTok{)}
    \ControlFlowTok{for}\OperatorTok{(}\DataTypeTok{size\_t}\NormalTok{ g }\OperatorTok{=} \DecValTok{0}\OperatorTok{;}\NormalTok{ g }\OperatorTok{\textless{}}\NormalTok{ ENERGY\_GROUPS\_NUM}\OperatorTok{;} \OperatorTok{++}\NormalTok{g}\OperatorTok{)}
\NormalTok{        cd}\OperatorTok{.}\NormalTok{dSdUm}\OperatorTok{[}\NormalTok{h}\OperatorTok{][}\NormalTok{g}\OperatorTok{]} \OperatorTok{*=}\NormalTok{ UmFactor}\OperatorTok{;}
\end{Highlighting}
\end{Shaded}

\begin{enumerate}
\def\labelenumi{\arabic{enumi}.}
\setcounter{enumi}{3}
\tightlist
\item
  Copy protections/cooling limiter logic only after the basic
  implementation passes parity tests. The first implementation can throw
  if the limiter would be needed, then add limiter support in a later
  phase.
\end{enumerate}

\subsubsection{5.4 Build residual kernels}\label{build-residual-kernels}

Add:

\begin{Shaded}
\begin{Highlighting}[]
\DataTypeTok{void}\NormalTok{ RadiationIMC}\OperatorTok{::}\NormalTok{BuildComptonResidualKernels}\OperatorTok{(}\DataTypeTok{size\_t}\NormalTok{ cellIndex}\OperatorTok{,} \DataTypeTok{double}\NormalTok{ fullDt}\OperatorTok{,}\NormalTok{ ComptonCellData}\OperatorTok{\&}\NormalTok{ cd}\OperatorTok{)}
\end{Highlighting}
\end{Shaded}

Pseudo-code:

\begin{Shaded}
\begin{Highlighting}[]
\ControlFlowTok{for}\OperatorTok{(}\DataTypeTok{size\_t}\NormalTok{ h }\OperatorTok{=} \DecValTok{0}\OperatorTok{;}\NormalTok{ h }\OperatorTok{\textless{}}\NormalTok{ ENERGY\_GROUPS\_NUM}\OperatorTok{;} \OperatorTok{++}\NormalTok{h}\OperatorTok{)}
\OperatorTok{\{}
    \ControlFlowTok{for}\OperatorTok{(}\DataTypeTok{size\_t}\NormalTok{ g }\OperatorTok{=} \DecValTok{0}\OperatorTok{;}\NormalTok{ g }\OperatorTok{\textless{}}\NormalTok{ ENERGY\_GROUPS\_NUM}\OperatorTok{;} \OperatorTok{++}\NormalTok{g}\OperatorTok{)}
    \OperatorTok{\{}
\NormalTok{        cd}\OperatorTok{.}\NormalTok{Kmat}\OperatorTok{[}\NormalTok{h}\OperatorTok{][}\NormalTok{g}\OperatorTok{]} \OperatorTok{=}\NormalTok{ cd}\OperatorTok{.}\NormalTok{betaCdtF }\OperatorTok{*}\NormalTok{ cd}\OperatorTok{.}\NormalTok{M}\OperatorTok{[}\NormalTok{g}\OperatorTok{]} \OperatorTok{*}\NormalTok{ cd}\OperatorTok{.}\NormalTok{Lambda}\OperatorTok{[}\NormalTok{h}\OperatorTok{];}
\NormalTok{        cd}\OperatorTok{.}\NormalTok{Hbase}\OperatorTok{[}\NormalTok{h}\OperatorTok{][}\NormalTok{g}\OperatorTok{]} \OperatorTok{=} \OperatorTok{(}\FloatTok{1.0} \OperatorTok{{-}}\NormalTok{ cd}\OperatorTok{.}\NormalTok{fleck}\OperatorTok{)} \OperatorTok{*}\NormalTok{ cd}\OperatorTok{.}\NormalTok{kappa}\OperatorTok{[}\NormalTok{h}\OperatorTok{]} \OperatorTok{*}\NormalTok{ cd}\OperatorTok{.}\NormalTok{qBase}\OperatorTok{[}\NormalTok{g}\OperatorTok{];}
    \OperatorTok{\}}
\OperatorTok{\}}

\ControlFlowTok{if}\OperatorTok{(}\KeywordTok{this}\OperatorTok{{-}\textgreater{}}\NormalTok{comptonUseAnalogDirect}\OperatorTok{)}
\OperatorTok{\{}
\NormalTok{    BuildAnalogDirectComptonMatrixAndCdfs}\OperatorTok{(}\NormalTok{cd}\OperatorTok{);}
\OperatorTok{\}}
\ControlFlowTok{else}
\OperatorTok{\{}
\NormalTok{    Zero}\OperatorTok{(}\NormalTok{cd}\OperatorTok{.}\NormalTok{Sanalog}\OperatorTok{);}
    \ControlFlowTok{for}\OperatorTok{(}\DataTypeTok{size\_t}\NormalTok{ h }\OperatorTok{=} \DecValTok{0}\OperatorTok{;}\NormalTok{ h }\OperatorTok{\textless{}}\NormalTok{ ENERGY\_GROUPS\_NUM}\OperatorTok{;} \OperatorTok{++}\NormalTok{h}\OperatorTok{)}
\NormalTok{        cd}\OperatorTok{.}\NormalTok{directOutRate}\OperatorTok{[}\NormalTok{h}\OperatorTok{]} \OperatorTok{=} \FloatTok{0.0}\OperatorTok{;}
\OperatorTok{\}}

\ControlFlowTok{for}\OperatorTok{(}\DataTypeTok{size\_t}\NormalTok{ h }\OperatorTok{=} \DecValTok{0}\OperatorTok{;}\NormalTok{ h }\OperatorTok{\textless{}}\NormalTok{ ENERGY\_GROUPS\_NUM}\OperatorTok{;} \OperatorTok{++}\NormalTok{h}\OperatorTok{)}
\OperatorTok{\{}
\NormalTok{    cd}\OperatorTok{.}\NormalTok{RposRow}\OperatorTok{[}\NormalTok{h}\OperatorTok{]} \OperatorTok{=} \FloatTok{0.0}\OperatorTok{;}
\NormalTok{    cd}\OperatorTok{.}\NormalTok{RnegRow}\OperatorTok{[}\NormalTok{h}\OperatorTok{]} \OperatorTok{=} \FloatTok{0.0}\OperatorTok{;}
    \ControlFlowTok{for}\OperatorTok{(}\DataTypeTok{size\_t}\NormalTok{ g }\OperatorTok{=} \DecValTok{0}\OperatorTok{;}\NormalTok{ g }\OperatorTok{\textless{}}\NormalTok{ ENERGY\_GROUPS\_NUM}\OperatorTok{;} \OperatorTok{++}\NormalTok{g}\OperatorTok{)}
    \OperatorTok{\{}
\NormalTok{        cd}\OperatorTok{.}\NormalTok{Sres}\OperatorTok{[}\NormalTok{h}\OperatorTok{][}\NormalTok{g}\OperatorTok{]} \OperatorTok{=}\NormalTok{ cd}\OperatorTok{.}\NormalTok{S}\OperatorTok{[}\NormalTok{h}\OperatorTok{][}\NormalTok{g}\OperatorTok{]} \OperatorTok{{-}}\NormalTok{ cd}\OperatorTok{.}\NormalTok{Sanalog}\OperatorTok{[}\NormalTok{h}\OperatorTok{][}\NormalTok{g}\OperatorTok{];}
\NormalTok{        cd}\OperatorTok{.}\NormalTok{R}\OperatorTok{[}\NormalTok{h}\OperatorTok{][}\NormalTok{g}\OperatorTok{]} \OperatorTok{=}\NormalTok{ cd}\OperatorTok{.}\NormalTok{Kmat}\OperatorTok{[}\NormalTok{h}\OperatorTok{][}\NormalTok{g}\OperatorTok{]} \OperatorTok{{-}}\NormalTok{ cd}\OperatorTok{.}\NormalTok{Hbase}\OperatorTok{[}\NormalTok{h}\OperatorTok{][}\NormalTok{g}\OperatorTok{]} \OperatorTok{+}\NormalTok{ cd}\OperatorTok{.}\NormalTok{Sres}\OperatorTok{[}\NormalTok{h}\OperatorTok{][}\NormalTok{g}\OperatorTok{];}
\NormalTok{        cd}\OperatorTok{.}\NormalTok{Rpos}\OperatorTok{[}\NormalTok{h}\OperatorTok{][}\NormalTok{g}\OperatorTok{]} \OperatorTok{=} \BuiltInTok{std::}\NormalTok{max}\OperatorTok{(}\NormalTok{cd}\OperatorTok{.}\NormalTok{R}\OperatorTok{[}\NormalTok{h}\OperatorTok{][}\NormalTok{g}\OperatorTok{],} \FloatTok{0.0}\OperatorTok{);}
\NormalTok{        cd}\OperatorTok{.}\NormalTok{Rneg}\OperatorTok{[}\NormalTok{h}\OperatorTok{][}\NormalTok{g}\OperatorTok{]} \OperatorTok{=} \BuiltInTok{std::}\NormalTok{max}\OperatorTok{({-}}\NormalTok{cd}\OperatorTok{.}\NormalTok{R}\OperatorTok{[}\NormalTok{h}\OperatorTok{][}\NormalTok{g}\OperatorTok{],} \FloatTok{0.0}\OperatorTok{);}
\NormalTok{        cd}\OperatorTok{.}\NormalTok{RposRow}\OperatorTok{[}\NormalTok{h}\OperatorTok{]} \OperatorTok{+=}\NormalTok{ cd}\OperatorTok{.}\NormalTok{Rpos}\OperatorTok{[}\NormalTok{h}\OperatorTok{][}\NormalTok{g}\OperatorTok{];}
\NormalTok{        cd}\OperatorTok{.}\NormalTok{RnegRow}\OperatorTok{[}\NormalTok{h}\OperatorTok{]} \OperatorTok{+=}\NormalTok{ cd}\OperatorTok{.}\NormalTok{Rneg}\OperatorTok{[}\NormalTok{h}\OperatorTok{][}\NormalTok{g}\OperatorTok{];}
    \OperatorTok{\}}
\NormalTok{    cd}\OperatorTok{.}\NormalTok{RposCdf}\OperatorTok{[}\NormalTok{h}\OperatorTok{]} \OperatorTok{=}\NormalTok{ BuildCdf}\OperatorTok{(}\NormalTok{cd}\OperatorTok{.}\NormalTok{Rpos}\OperatorTok{[}\NormalTok{h}\OperatorTok{]);}
\NormalTok{    cd}\OperatorTok{.}\NormalTok{RnegCdf}\OperatorTok{[}\NormalTok{h}\OperatorTok{]} \OperatorTok{=}\NormalTok{ BuildCdf}\OperatorTok{(}\NormalTok{cd}\OperatorTok{.}\NormalTok{Rneg}\OperatorTok{[}\NormalTok{h}\OperatorTok{]);}
\OperatorTok{\}}
\end{Highlighting}
\end{Shaded}

Important consistency check:

For every source group \texttt{h}, the existing base effective
scattering event plus the residual split must reconstruct the desired
kernel:

\begin{Shaded}
\begin{Highlighting}[]
\NormalTok{H\^{}\{base\}\_\{h\textbackslash{}to g\} + R\^{}+\_\{h\textbackslash{}to g\} {-} R\^{}{-}\_\{h\textbackslash{}to g\} + S\^{}\{analog\}\_\{h\textbackslash{}to g\}}
\NormalTok{= K\^{}\{mat\}\_\{h\textbackslash{}to g\}+S\_\{h\textbackslash{}to g\}.}
\end{Highlighting}
\end{Shaded}

Add a debug-only assertion that the absolute difference is below a
tolerance.

\subsubsection{5.5 Analog direct Compton
option}\label{analog-direct-compton-option}

This option is recommended but can be implemented after the signed
residual version works.

Build positive physical Compton event rates from \texttt{tau}:

\begin{Shaded}
\begin{Highlighting}[]
\DataTypeTok{void}\NormalTok{ RadiationIMC}\OperatorTok{::}\NormalTok{BuildAnalogDirectComptonMatrixAndCdfs}\OperatorTok{(}\NormalTok{ComptonCellData}\OperatorTok{\&}\NormalTok{ cd}\OperatorTok{)}
\OperatorTok{\{}
\NormalTok{    Zero}\OperatorTok{(}\NormalTok{cd}\OperatorTok{.}\NormalTok{Sanalog}\OperatorTok{);}
    \ControlFlowTok{for}\OperatorTok{(}\DataTypeTok{size\_t}\NormalTok{ h }\OperatorTok{=} \DecValTok{0}\OperatorTok{;}\NormalTok{ h }\OperatorTok{\textless{}}\NormalTok{ ENERGY\_GROUPS\_NUM}\OperatorTok{;} \OperatorTok{++}\NormalTok{h}\OperatorTok{)}
    \OperatorTok{\{}
        \BuiltInTok{std::}\NormalTok{array}\OperatorTok{\textless{}}\DataTypeTok{double}\OperatorTok{,}\NormalTok{ ENERGY\_GROUPS\_NUM}\OperatorTok{\textgreater{}}\NormalTok{ targetRate}\OperatorTok{\{\};}
        \DataTypeTok{double}\NormalTok{ out }\OperatorTok{=} \FloatTok{0.0}\OperatorTok{;}
        \ControlFlowTok{for}\OperatorTok{(}\DataTypeTok{size\_t}\NormalTok{ g }\OperatorTok{=} \DecValTok{0}\OperatorTok{;}\NormalTok{ g }\OperatorTok{\textless{}}\NormalTok{ ENERGY\_GROUPS\_NUM}\OperatorTok{;} \OperatorTok{++}\NormalTok{g}\OperatorTok{)}
        \OperatorTok{\{}
            \ControlFlowTok{if}\OperatorTok{(}\NormalTok{g }\OperatorTok{==}\NormalTok{ h}\OperatorTok{)}
                \ControlFlowTok{continue}\OperatorTok{;}  \CommentTok{// same{-}group transitions do not change multigroup energy in the simple analog sampler}
            \AttributeTok{const} \DataTypeTok{double}\NormalTok{ rate }\OperatorTok{=}\NormalTok{ cd}\OperatorTok{.}\NormalTok{tau}\OperatorTok{[}\NormalTok{h}\OperatorTok{][}\NormalTok{g}\OperatorTok{]} \OperatorTok{*} \OperatorTok{(}\FloatTok{1.0} \OperatorTok{+}\NormalTok{ n}\OperatorTok{[}\NormalTok{g}\OperatorTok{]);}
\NormalTok{            targetRate}\OperatorTok{[}\NormalTok{g}\OperatorTok{]} \OperatorTok{=} \BuiltInTok{std::}\NormalTok{max}\OperatorTok{(}\NormalTok{rate}\OperatorTok{,} \FloatTok{0.0}\OperatorTok{);}
\NormalTok{            out }\OperatorTok{+=}\NormalTok{ targetRate}\OperatorTok{[}\NormalTok{g}\OperatorTok{];}

            \AttributeTok{const} \DataTypeTok{double}\NormalTok{ ratio }\OperatorTok{=}\NormalTok{ energyGroupCenters}\OperatorTok{[}\NormalTok{g}\OperatorTok{]} \OperatorTok{/}\NormalTok{ energyGroupCenters}\OperatorTok{[}\NormalTok{h}\OperatorTok{];}
\NormalTok{            cd}\OperatorTok{.}\NormalTok{Sanalog}\OperatorTok{[}\NormalTok{h}\OperatorTok{][}\NormalTok{g}\OperatorTok{]} \OperatorTok{+=}\NormalTok{ targetRate}\OperatorTok{[}\NormalTok{g}\OperatorTok{]} \OperatorTok{*}\NormalTok{ ratio}\OperatorTok{;}
        \OperatorTok{\}}
\NormalTok{        cd}\OperatorTok{.}\NormalTok{Sanalog}\OperatorTok{[}\NormalTok{h}\OperatorTok{][}\NormalTok{h}\OperatorTok{]} \OperatorTok{{-}=}\NormalTok{ out}\OperatorTok{;}
\NormalTok{        cd}\OperatorTok{.}\NormalTok{directOutRate}\OperatorTok{[}\NormalTok{h}\OperatorTok{]} \OperatorTok{=}\NormalTok{ out}\OperatorTok{;}
\NormalTok{        cd}\OperatorTok{.}\NormalTok{directCdf}\OperatorTok{[}\NormalTok{h}\OperatorTok{]} \OperatorTok{=}\NormalTok{ BuildCdf}\OperatorTok{(}\NormalTok{targetRate}\OperatorTok{);}
    \OperatorTok{\}}
\OperatorTok{\}}
\end{Highlighting}
\end{Shaded}

The exact rate formula must match the same induced factors used in
\texttt{S}. The residual \texttt{S\ -\ Sanalog} will correct any
mismatch, including last-group boundary corrections. Therefore the
analog sampler can start simple and become more accurate later without
changing the expected linear operator.

If \texttt{opacity-\textgreater{}CalcScatteringOpacity(cell)} already
includes the same electron scattering represented by \texttt{tau}, do
not double count it. Add a parameter or an opacity API method to
separate elastic non-Compton scattering from Compton scattering. Until
that separation is clear, set \texttt{comptonUseAnalogDirect=false} and
let \texttt{S} be handled by signed residuals.

\subsubsection{5.6 Build source terms}\label{build-source-terms}

Add:

\begin{Shaded}
\begin{Highlighting}[]
\DataTypeTok{void}\NormalTok{ RadiationIMC}\OperatorTok{::}\NormalTok{BuildComptonSources}\OperatorTok{(}\DataTypeTok{size\_t}\NormalTok{ cellIndex}\OperatorTok{,} \DataTypeTok{double}\NormalTok{ fullDt}\OperatorTok{,}\NormalTok{ ComptonCellData}\OperatorTok{\&}\NormalTok{ cd}\OperatorTok{)}
\end{Highlighting}
\end{Shaded}

Pseudo-code:

\begin{Shaded}
\begin{Highlighting}[]
\AttributeTok{const} \DataTypeTok{double}\NormalTok{ V }\OperatorTok{=}\NormalTok{ cd}\OperatorTok{.}\NormalTok{volume}\OperatorTok{;}
\AttributeTok{const} \DataTypeTok{double}\NormalTok{ cdt }\OperatorTok{=}\NormalTok{ units}\OperatorTok{::}\NormalTok{clight }\OperatorTok{*}\NormalTok{ fullDt}\OperatorTok{;}

\ControlFlowTok{for}\OperatorTok{(}\DataTypeTok{size\_t}\NormalTok{ g }\OperatorTok{=} \DecValTok{0}\OperatorTok{;}\NormalTok{ g }\OperatorTok{\textless{}}\NormalTok{ ENERGY\_GROUPS\_NUM}\OperatorTok{;} \OperatorTok{++}\NormalTok{g}\OperatorTok{)}
\OperatorTok{\{}
    \AttributeTok{const} \DataTypeTok{double}\NormalTok{ kgbg }\OperatorTok{=}\NormalTok{ cd}\OperatorTok{.}\NormalTok{kappa}\OperatorTok{[}\NormalTok{g}\OperatorTok{]} \OperatorTok{*}\NormalTok{ cd}\OperatorTok{.}\NormalTok{b}\OperatorTok{[}\NormalTok{g}\OperatorTok{];}
\NormalTok{    cd}\OperatorTok{.}\NormalTok{Bbase}\OperatorTok{[}\NormalTok{g}\OperatorTok{]} \OperatorTok{=}\NormalTok{ V }\OperatorTok{*}\NormalTok{ cdt }\OperatorTok{*}\NormalTok{ cd}\OperatorTok{.}\NormalTok{fleck }\OperatorTok{*}\NormalTok{ kgbg }\OperatorTok{*}\NormalTok{ cd}\OperatorTok{.}\NormalTok{Um}\OperatorTok{;}

    \ControlFlowTok{if}\OperatorTok{(}\NormalTok{cd}\OperatorTok{.}\NormalTok{kappaP }\OperatorTok{\textgreater{}} \FloatTok{0.0}\OperatorTok{)}
    \OperatorTok{\{}
\NormalTok{        cd}\OperatorTok{.}\NormalTok{Bcorr}\OperatorTok{[}\NormalTok{g}\OperatorTok{]} \OperatorTok{=}\NormalTok{ V }\OperatorTok{*}\NormalTok{ cdt }\OperatorTok{*}\NormalTok{ cd}\OperatorTok{.}\NormalTok{kappaP }\OperatorTok{*}\NormalTok{ cd}\OperatorTok{.}\NormalTok{Um }\OperatorTok{*}
            \OperatorTok{(}
                \OperatorTok{(}\NormalTok{kgbg }\OperatorTok{/}\NormalTok{ cd}\OperatorTok{.}\NormalTok{kappaP}\OperatorTok{)} \OperatorTok{*} \OperatorTok{(}\FloatTok{1.0} \OperatorTok{{-}} \OperatorTok{(}\FloatTok{1.0} \OperatorTok{+}\NormalTok{ cd}\OperatorTok{.}\NormalTok{beta }\OperatorTok{*}\NormalTok{ cdt }\OperatorTok{*}\NormalTok{ cd}\OperatorTok{.}\NormalTok{kappaP}\OperatorTok{)} \OperatorTok{*}\NormalTok{ cd}\OperatorTok{.}\NormalTok{fleck}\OperatorTok{)}
                \OperatorTok{{-}}\NormalTok{ cd}\OperatorTok{.}\NormalTok{beta }\OperatorTok{*}\NormalTok{ cdt }\OperatorTok{*}\NormalTok{ cd}\OperatorTok{.}\NormalTok{fleck }\OperatorTok{*}\NormalTok{ cd}\OperatorTok{.}\NormalTok{D}\OperatorTok{[}\NormalTok{g}\OperatorTok{]}
            \OperatorTok{);}
    \OperatorTok{\}}
    \ControlFlowTok{else}
    \OperatorTok{\{}
        \CommentTok{// Match deterministic behavior. If kappaP is zero, the expression in the deterministic code would divide by kappaP.}
        \CommentTok{// Since kgbg is also zero if kappaP is a sum of nonnegative kgbg, set correction to zero.}
\NormalTok{        cd}\OperatorTok{.}\NormalTok{Bcorr}\OperatorTok{[}\NormalTok{g}\OperatorTok{]} \OperatorTok{=} \FloatTok{0.0}\OperatorTok{;}
    \OperatorTok{\}}

\NormalTok{    cd}\OperatorTok{.}\NormalTok{Btotal}\OperatorTok{[}\NormalTok{g}\OperatorTok{]} \OperatorTok{=}\NormalTok{ cd}\OperatorTok{.}\NormalTok{Bbase}\OperatorTok{[}\NormalTok{g}\OperatorTok{]} \OperatorTok{+}\NormalTok{ cd}\OperatorTok{.}\NormalTok{Bcorr}\OperatorTok{[}\NormalTok{g}\OperatorTok{];}
\OperatorTok{\}}
\end{Highlighting}
\end{Shaded}

If hydro Lorentz \texttt{gamma} was included in the Fleck factor, do not
multiply \texttt{Btotal} by gamma here. Follow the existing source
emission convention: debit material internal energy by comoving source
energy and assign particle lab weights using Doppler factors in
\texttt{generateComptonParticles}.

\subsection{6. Source-particle generation with signed group
source}\label{source-particle-generation-with-signed-group-source}

Modify \texttt{RadiationIMC::generateParticles}:

\begin{Shaded}
\begin{Highlighting}[]
\BuiltInTok{std::}\NormalTok{vector}\OperatorTok{\textless{}}\KeywordTok{typename}\NormalTok{ RadiationIMC}\OperatorTok{::}\NormalTok{Particle}\OperatorTok{\textgreater{}}\NormalTok{ RadiationIMC}\OperatorTok{::}\NormalTok{generateParticles}\OperatorTok{(}\DataTypeTok{double}\NormalTok{ fullDt}\OperatorTok{)}
\OperatorTok{\{}
    \ControlFlowTok{if}\OperatorTok{(}\KeywordTok{this}\OperatorTok{{-}\textgreater{}}\NormalTok{withCompton}\OperatorTok{)}
        \ControlFlowTok{return} \KeywordTok{this}\OperatorTok{{-}\textgreater{}}\NormalTok{generateComptonParticles}\OperatorTok{(}\NormalTok{fullDt}\OperatorTok{);}

    \CommentTok{// existing implementation unchanged}
\OperatorTok{\}}
\end{Highlighting}
\end{Shaded}

Add:

\begin{Shaded}
\begin{Highlighting}[]
\BuiltInTok{std::}\NormalTok{vector}\OperatorTok{\textless{}}\KeywordTok{typename}\NormalTok{ RadiationIMC}\OperatorTok{::}\NormalTok{Particle}\OperatorTok{\textgreater{}}\NormalTok{ RadiationIMC}\OperatorTok{::}\NormalTok{generateComptonParticles}\OperatorTok{(}\DataTypeTok{double}\NormalTok{ fullDt}\OperatorTok{);}
\end{Highlighting}
\end{Shaded}

\subsubsection{6.1 Particle count
allocation}\label{particle-count-allocation}

The old code allocates particles per cell using positive
\texttt{energyToCreateVec}. With signed group sources, allocate using
absolute emitted energy:

\begin{Shaded}
\begin{Highlighting}[]
\BuiltInTok{std::}\NormalTok{vector}\OperatorTok{\textless{}}\DataTypeTok{double}\OperatorTok{\textgreater{}}\NormalTok{ absEnergyToCreateVec}\OperatorTok{(}\NormalTok{Ncells}\OperatorTok{,} \FloatTok{0.0}\OperatorTok{);}
\DataTypeTok{double}\NormalTok{ localAbsTotalEnergy }\OperatorTok{=} \FloatTok{0.0}\OperatorTok{;}
\ControlFlowTok{for}\OperatorTok{(}\DataTypeTok{size\_t}\NormalTok{ i }\OperatorTok{=} \DecValTok{0}\OperatorTok{;}\NormalTok{ i }\OperatorTok{\textless{}}\NormalTok{ Ncells}\OperatorTok{;} \OperatorTok{++}\NormalTok{i}\OperatorTok{)}
\OperatorTok{\{}
    \AttributeTok{const} \KeywordTok{auto}\OperatorTok{\&}\NormalTok{ cd }\OperatorTok{=}\NormalTok{ comptonData}\OperatorTok{[}\NormalTok{i}\OperatorTok{];}
    \ControlFlowTok{for}\OperatorTok{(}\DataTypeTok{size\_t}\NormalTok{ g }\OperatorTok{=} \DecValTok{0}\OperatorTok{;}\NormalTok{ g }\OperatorTok{\textless{}}\NormalTok{ ENERGY\_GROUPS\_NUM}\OperatorTok{;} \OperatorTok{++}\NormalTok{g}\OperatorTok{)}
\NormalTok{        absEnergyToCreateVec}\OperatorTok{[}\NormalTok{i}\OperatorTok{]} \OperatorTok{+=} \BuiltInTok{std::}\NormalTok{abs}\OperatorTok{(}\NormalTok{cd}\OperatorTok{.}\NormalTok{Btotal}\OperatorTok{[}\NormalTok{g}\OperatorTok{]);}
\NormalTok{    localAbsTotalEnergy }\OperatorTok{+=}\NormalTok{ absEnergyToCreateVec}\OperatorTok{[}\NormalTok{i}\OperatorTok{];}
\OperatorTok{\}}
\end{Highlighting}
\end{Shaded}

Use the same MPI all-reduce pattern as lines 371-376, but reduce
absolute total energy.

For each cell, choose a total source-particle target:

\begin{Shaded}
\begin{Highlighting}[]
\DataTypeTok{size\_t}\NormalTok{ nTargetCell}\OperatorTok{;}
\ControlFlowTok{if}\OperatorTok{(}\NormalTok{globalAbsTotalEnergy }\OperatorTok{\textgreater{}} \FloatTok{0.0}\OperatorTok{)}
\OperatorTok{\{}
\NormalTok{    nTargetCell }\OperatorTok{=} \KeywordTok{static\_cast}\OperatorTok{\textless{}}\DataTypeTok{size\_t}\OperatorTok{\textgreater{}(}\NormalTok{absEnergyToCreateVec}\OperatorTok{[}\NormalTok{i}\OperatorTok{]} \OperatorTok{/}\NormalTok{ globalAbsTotalEnergy }\OperatorTok{*}\NormalTok{ totalParticles}\OperatorTok{);}
\OperatorTok{\}}
\ControlFlowTok{else}
\OperatorTok{\{}
\NormalTok{    nTargetCell }\OperatorTok{=} \DecValTok{0}\OperatorTok{;}
\OperatorTok{\}}
\NormalTok{nTargetCell }\OperatorTok{=} \BuiltInTok{std::}\NormalTok{max}\OperatorTok{\textless{}}\DataTypeTok{size\_t}\OperatorTok{\textgreater{}(}\NormalTok{newPhotonsPerCell}\OperatorTok{,} \BuiltInTok{std::}\NormalTok{min}\OperatorTok{(}\NormalTok{nTargetCell}\OperatorTok{,}\NormalTok{ newPhotonsPerCell }\OperatorTok{*}\NormalTok{ maxSourcePacketsPerCellMultiplier}\OperatorTok{));}
\end{Highlighting}
\end{Shaded}

Then distribute \texttt{nTargetCell} among signed nonzero groups
proportional to \texttt{abs(Btotal{[}g{]})}. Make at least
\texttt{minSourcePacketsPerSignedGroup} for a nonzero group if possible.

\subsubsection{6.2 Material debit/credit}\label{material-debitcredit}

For every group source energy \texttt{Eemit\ =\ cd.Btotal{[}g{]}},
update material by \texttt{-Eemit}.

\begin{Shaded}
\begin{Highlighting}[]
\ControlFlowTok{if}\OperatorTok{(!}\KeywordTok{this}\OperatorTok{{-}\textgreater{}}\NormalTok{noHydroFeedback}\OperatorTok{)}
\OperatorTok{\{}
    \KeywordTok{this}\OperatorTok{{-}\textgreater{}}\NormalTok{conserved}\OperatorTok{[}\NormalTok{i}\OperatorTok{].}\NormalTok{internal\_energy }\OperatorTok{{-}=}\NormalTok{ Eemit}\OperatorTok{;}
    \ControlFlowTok{if}\OperatorTok{(}\KeywordTok{this}\OperatorTok{{-}\textgreater{}}\NormalTok{conserved}\OperatorTok{[}\NormalTok{i}\OperatorTok{].}\NormalTok{internal\_energy }\OperatorTok{\textless{}} \FloatTok{0.0}\OperatorTok{)}
    \OperatorTok{\{}
        \CommentTok{// Only throw if Eemit is positive and caused the material to go negative.}
        \CommentTok{// A negative Eemit should increase material energy.}
    \OperatorTok{\}}

    \ControlFlowTok{if}\OperatorTok{(}\KeywordTok{this}\OperatorTok{{-}\textgreater{}}\NormalTok{withHydro}\OperatorTok{)}
    \OperatorTok{\{}
        \AttributeTok{const} \DataTypeTok{double}\NormalTok{ gamma }\OperatorTok{=} \OperatorTok{...;}
        \KeywordTok{this}\OperatorTok{{-}\textgreater{}}\NormalTok{conserved}\OperatorTok{[}\NormalTok{i}\OperatorTok{].}\NormalTok{energy }\OperatorTok{{-}=}\NormalTok{ Eemit }\OperatorTok{*}\NormalTok{ gamma}\OperatorTok{;}
        \ControlFlowTok{if}\OperatorTok{(!}\KeywordTok{this}\OperatorTok{{-}\textgreater{}}\NormalTok{diffusionPressureGradient}\OperatorTok{)}
            \KeywordTok{this}\OperatorTok{{-}\textgreater{}}\NormalTok{conserved}\OperatorTok{[}\NormalTok{i}\OperatorTok{].}\NormalTok{momentum }\OperatorTok{{-}=}\NormalTok{ Eemit }\OperatorTok{*}\NormalTok{ cell}\OperatorTok{.}\NormalTok{velocity }\OperatorTok{*}\NormalTok{ units}\OperatorTok{::}\NormalTok{inv\_clight2 }\OperatorTok{*}\NormalTok{ gamma}\OperatorTok{;}
    \OperatorTok{\}}
\OperatorTok{\}}
\end{Highlighting}
\end{Shaded}

This mirrors lines 394-421, but \texttt{Eemit} may be negative.

\subsubsection{6.3 Particle creation}\label{particle-creation}

For each group \texttt{g} with \texttt{ng} packets:

\begin{Shaded}
\begin{Highlighting}[]
\AttributeTok{const} \DataTypeTok{double}\NormalTok{ Eemit }\OperatorTok{=}\NormalTok{ cd}\OperatorTok{.}\NormalTok{Btotal}\OperatorTok{[}\NormalTok{g}\OperatorTok{];}
\AttributeTok{const} \DataTypeTok{double}\NormalTok{ EperPacket }\OperatorTok{=}\NormalTok{ Eemit }\OperatorTok{/} \KeywordTok{static\_cast}\OperatorTok{\textless{}}\DataTypeTok{double}\OperatorTok{\textgreater{}(}\NormalTok{ng}\OperatorTok{);}
\ControlFlowTok{for}\OperatorTok{(}\DataTypeTok{size\_t}\NormalTok{ j }\OperatorTok{=} \DecValTok{0}\OperatorTok{;}\NormalTok{ j }\OperatorTok{\textless{}}\NormalTok{ ng}\OperatorTok{;} \OperatorTok{++}\NormalTok{j}\OperatorTok{)}
\OperatorTok{\{}
\NormalTok{    Particle particle }\OperatorTok{=} \KeywordTok{this}\OperatorTok{{-}\textgreater{}}\NormalTok{generateSingleParticle}\OperatorTok{(}\NormalTok{i}\OperatorTok{,}\NormalTok{ cell}\OperatorTok{);}
\NormalTok{    particle}\OperatorTok{.}\NormalTok{timeLeft }\OperatorTok{=}\NormalTok{ fullDt }\OperatorTok{*} \KeywordTok{this}\OperatorTok{{-}\textgreater{}}\NormalTok{dist}\OperatorTok{(}\KeywordTok{this}\OperatorTok{{-}\textgreater{}}\NormalTok{re}\OperatorTok{);}
\NormalTok{    particle}\OperatorTok{.}\NormalTok{cellID }\OperatorTok{=}\NormalTok{ cell}\OperatorTok{.}\NormalTok{ID}\OperatorTok{;}

    \ControlFlowTok{if}\OperatorTok{(}\KeywordTok{this}\OperatorTok{{-}\textgreater{}}\NormalTok{withHydro }\OperatorTok{\&\&} \OperatorTok{!}\KeywordTok{this}\OperatorTok{{-}\textgreater{}}\NormalTok{MMC}\OperatorTok{)}
    \OperatorTok{\{}
        \AttributeTok{const} \DataTypeTok{double}\NormalTok{ D }\OperatorTok{=}\NormalTok{ DopplerShift}\OperatorTok{(}\NormalTok{particle}\OperatorTok{,}\NormalTok{ cell}\OperatorTok{.}\NormalTok{velocity}\OperatorTok{);}
        \AttributeTok{const} \DataTypeTok{double}\NormalTok{ freqCo }\OperatorTok{=}\NormalTok{ FrequencyForGroup}\OperatorTok{(}\NormalTok{g}\OperatorTok{);}
\NormalTok{        particle}\OperatorTok{.}\NormalTok{frequency }\OperatorTok{=}\NormalTok{ freqCo }\OperatorTok{/}\NormalTok{ D}\OperatorTok{;}
\NormalTok{        particle}\OperatorTok{.}\NormalTok{weight }\OperatorTok{=}\NormalTok{ EperPacket }\OperatorTok{/}\NormalTok{ D}\OperatorTok{;}
    \OperatorTok{\}}
    \ControlFlowTok{else}
    \OperatorTok{\{}
\NormalTok{        particle}\OperatorTok{.}\NormalTok{frequency }\OperatorTok{=}\NormalTok{ FrequencyForGroup}\OperatorTok{(}\NormalTok{g}\OperatorTok{);}
\NormalTok{        particle}\OperatorTok{.}\NormalTok{weight }\OperatorTok{=}\NormalTok{ EperPacket}\OperatorTok{;}
    \OperatorTok{\}}

\NormalTok{    ClampFrequencyToBounds}\OperatorTok{(}\NormalTok{particle}\OperatorTok{.}\NormalTok{frequency}\OperatorTok{);}
\NormalTok{    particle}\OperatorTok{.}\NormalTok{initialWeight }\OperatorTok{=} \BuiltInTok{std::}\NormalTok{abs}\OperatorTok{(}\NormalTok{particle}\OperatorTok{.}\NormalTok{weight}\OperatorTok{);}
    \ControlFlowTok{if}\OperatorTok{(}\NormalTok{particle}\OperatorTok{.}\NormalTok{initialWeight }\OperatorTok{\textgreater{}} \FloatTok{0.0}\OperatorTok{)}
\NormalTok{        newParticles}\OperatorTok{.}\NormalTok{push\_back}\OperatorTok{(}\NormalTok{particle}\OperatorTok{);}
\OperatorTok{\}}
\end{Highlighting}
\end{Shaded}

Do not use \texttt{multigroupOpacity-\textgreater{}GetThermalEnergy} in
the Compton source path unless it can sample a specified target group.
The source is group-signed and must emit into the selected group.

\subsection{7. Step function changes}\label{step-function-changes}

\texttt{RadiationIMC::step} is the center of the implementation.

\subsubsection{\texorpdfstring{7.1 Remove the unused marker for
\texttt{particlesToAdd}}{7.1 Remove the unused marker for particlesToAdd}}\label{remove-the-unused-marker-for-particlestoadd}

Delete or guard line 42:

\begin{Shaded}
\begin{Highlighting}[]
\OperatorTok{(}\DataTypeTok{void}\OperatorTok{)}\NormalTok{ particlesToAdd}\OperatorTok{;}
\end{Highlighting}
\end{Shaded}

Compton residuals will push child packets into this vector.

\subsubsection{7.2 Determine the current group and consistent
opacities}\label{determine-the-current-group-and-consistent-opacities}

Near lines 65-75, when Compton is enabled, use the precomputed group
opacity from \texttt{comptonData}, not a fresh opacity call. This keeps
the transport, source, and residual matrices consistent.

\begin{Shaded}
\begin{Highlighting}[]
\DataTypeTok{size\_t}\NormalTok{ group }\OperatorTok{=} \BuiltInTok{std::}\NormalTok{numeric\_limits}\OperatorTok{\textless{}}\DataTypeTok{size\_t}\OperatorTok{\textgreater{}::}\NormalTok{max}\OperatorTok{();}
\DataTypeTok{double}\NormalTok{ absorptionOpacity}\OperatorTok{;}
\ControlFlowTok{if}\OperatorTok{(}\KeywordTok{this}\OperatorTok{{-}\textgreater{}}\NormalTok{multigroupOpacity}\OperatorTok{)}
\OperatorTok{\{}
    \DataTypeTok{double}\NormalTok{ shiftedFrequency }\OperatorTok{=}\NormalTok{ particle}\OperatorTok{.}\NormalTok{frequency }\OperatorTok{*}\NormalTok{ dopplerShift}\OperatorTok{;}
\NormalTok{    ClampFrequencyToBounds}\OperatorTok{(}\NormalTok{shiftedFrequency}\OperatorTok{);}
\NormalTok{    group }\OperatorTok{=} \KeywordTok{this}\OperatorTok{{-}\textgreater{}}\NormalTok{opacity}\OperatorTok{{-}\textgreater{}}\NormalTok{findGroup}\OperatorTok{(}\NormalTok{shiftedFrequency}\OperatorTok{);}
    \ControlFlowTok{if}\OperatorTok{(}\KeywordTok{this}\OperatorTok{{-}\textgreater{}}\NormalTok{withCompton}\OperatorTok{)}
\NormalTok{        absorptionOpacity }\OperatorTok{=} \KeywordTok{this}\OperatorTok{{-}\textgreater{}}\NormalTok{comptonData}\OperatorTok{[}\NormalTok{cellIndex}\OperatorTok{].}\NormalTok{kappa}\OperatorTok{[}\NormalTok{group}\OperatorTok{];}
    \ControlFlowTok{else}
\NormalTok{        absorptionOpacity }\OperatorTok{=} \KeywordTok{this}\OperatorTok{{-}\textgreater{}}\NormalTok{opacity}\OperatorTok{{-}\textgreater{}}\NormalTok{CalcAbsorptionOpacity}\OperatorTok{(}\NormalTok{cell}\OperatorTok{,}\NormalTok{ shiftedFrequency}\OperatorTok{);}
\OperatorTok{\}}
\ControlFlowTok{else}
\OperatorTok{\{}
\NormalTok{    absorptionOpacity }\OperatorTok{=} \KeywordTok{this}\OperatorTok{{-}\textgreater{}}\NormalTok{planckOpacities}\OperatorTok{[}\NormalTok{cellIndex}\OperatorTok{];}
\OperatorTok{\}}
\end{Highlighting}
\end{Shaded}

\subsubsection{7.3 Event opacities}\label{event-opacities}

Current line 76:

\begin{Shaded}
\begin{Highlighting}[]
\DataTypeTok{double}\NormalTok{ scatteringOpacity }\OperatorTok{=} \KeywordTok{this}\OperatorTok{{-}\textgreater{}}\NormalTok{opacity}\OperatorTok{{-}\textgreater{}}\NormalTok{CalcScatteringOpacity}\OperatorTok{(}\NormalTok{cell}\OperatorTok{);}
\end{Highlighting}
\end{Shaded}

Change to separate event channels:

\begin{Shaded}
\begin{Highlighting}[]
\DataTypeTok{double}\NormalTok{ elasticScatteringOpacity }\OperatorTok{=} \KeywordTok{this}\OperatorTok{{-}\textgreater{}}\NormalTok{opacity}\OperatorTok{{-}\textgreater{}}\NormalTok{CalcScatteringOpacity}\OperatorTok{(}\NormalTok{cell}\OperatorTok{);}
\DataTypeTok{double}\NormalTok{ effectiveAbsorptionOpacity }\OperatorTok{=} \OperatorTok{(}\FloatTok{1.0} \OperatorTok{{-}} \KeywordTok{this}\OperatorTok{{-}\textgreater{}}\NormalTok{factorFleck}\OperatorTok{[}\NormalTok{cellIndex}\OperatorTok{])} \OperatorTok{*}\NormalTok{ absorptionOpacity}\OperatorTok{;}
\DataTypeTok{double}\NormalTok{ directComptonOpacity }\OperatorTok{=} \FloatTok{0.0}\OperatorTok{;}

\ControlFlowTok{if}\OperatorTok{(}\KeywordTok{this}\OperatorTok{{-}\textgreater{}}\NormalTok{withCompton }\OperatorTok{\&\&} \KeywordTok{this}\OperatorTok{{-}\textgreater{}}\NormalTok{comptonUseAnalogDirect}\OperatorTok{)}
\OperatorTok{\{}
\NormalTok{    directComptonOpacity }\OperatorTok{=} \KeywordTok{this}\OperatorTok{{-}\textgreater{}}\NormalTok{comptonData}\OperatorTok{[}\NormalTok{cellIndex}\OperatorTok{].}\NormalTok{directOutRate}\OperatorTok{[}\NormalTok{group}\OperatorTok{];}

    \CommentTok{// Avoid double counting if CalcScatteringOpacity is the same electron scattering.}
    \CommentTok{// If no separate opacity API exists, add a parameter that either leaves elasticScatteringOpacity alone}
    \CommentTok{// or sets it to zero in Compton runs.}
\OperatorTok{\}}

\DataTypeTok{double}\NormalTok{ eventOpacity }\OperatorTok{=}\NormalTok{ elasticScatteringOpacity }\OperatorTok{+}\NormalTok{ effectiveAbsorptionOpacity }\OperatorTok{+}\NormalTok{ directComptonOpacity}\OperatorTok{;}
\end{Highlighting}
\end{Shaded}

Then use \texttt{eventOpacity} in the scattering-distance calculation:

\begin{Shaded}
\begin{Highlighting}[]
\DataTypeTok{double}\NormalTok{ scatteringLength }\OperatorTok{=} \OperatorTok{(}\NormalTok{eventOpacity }\OperatorTok{\textgreater{}} \FloatTok{0.0}\OperatorTok{)} \OperatorTok{?} \FloatTok{1.0} \OperatorTok{/}\NormalTok{ eventOpacity }\OperatorTok{:} \BuiltInTok{std::}\NormalTok{numeric\_limits}\OperatorTok{\textless{}}\DataTypeTok{double}\OperatorTok{\textgreater{}::}\NormalTok{infinity}\OperatorTok{();}
\end{Highlighting}
\end{Shaded}

If \texttt{eventOpacity\ ==\ 0}, set
\texttt{timeScattering\ =\ infinity}.

\subsubsection{7.4 Save segment-start
state}\label{save-segment-start-state}

Before subtracting \texttt{particle.timeLeft} and before updating
location, save:

\begin{Shaded}
\begin{Highlighting}[]
\AttributeTok{const}\NormalTok{ Vector3D segmentStartLocation }\OperatorTok{=}\NormalTok{ particle}\OperatorTok{.}\NormalTok{location}\OperatorTok{;}
\AttributeTok{const} \DataTypeTok{double}\NormalTok{ segmentStartTimeLeft }\OperatorTok{=}\NormalTok{ particle}\OperatorTok{.}\NormalTok{timeLeft}\OperatorTok{;}
\AttributeTok{const} \DataTypeTok{double}\NormalTok{ segmentStartWeight }\OperatorTok{=}\NormalTok{ particle}\OperatorTok{.}\NormalTok{weight}\OperatorTok{;}
\AttributeTok{const} \DataTypeTok{size\_t}\NormalTok{ segmentGroup }\OperatorTok{=}\NormalTok{ group}\OperatorTok{;}
\end{Highlighting}
\end{Shaded}

Then current code can do:

\begin{Shaded}
\begin{Highlighting}[]
\NormalTok{particle}\OperatorTok{.}\NormalTok{timeLeft }\OperatorTok{{-}=}\NormalTok{ dt}\OperatorTok{;}
\end{Highlighting}
\end{Shaded}

\subsubsection{7.5 Generate residual children on every path
segment}\label{generate-residual-children-on-every-path-segment}

After \texttt{dt} is chosen and before or just after continuous
absorption is applied, call a new helper that uses the segment-start
weight and the segment length/time.

Recommended placement:

\begin{enumerate}
\def\labelenumi{\arabic{enumi}.}
\tightlist
\item
  Save segment-start state.
\item
  Compute \texttt{alpha\ =\ factorFleck\ *\ absorptionOpacity\ *\ c}.
\item
  Emit residual children using the integrated attenuated weight over the
  segment.
\item
  Apply the existing continuous absorption/material update.
\item
  Move/update the parent and handle the event.
\end{enumerate}

Add:

\begin{Shaded}
\begin{Highlighting}[]
\ControlFlowTok{if}\OperatorTok{(}\KeywordTok{this}\OperatorTok{{-}\textgreater{}}\NormalTok{withCompton }\OperatorTok{\&\&} \KeywordTok{this}\OperatorTok{{-}\textgreater{}}\NormalTok{multigroupOpacity }\OperatorTok{\&\&}\NormalTok{ dt }\OperatorTok{\textgreater{}} \FloatTok{0.0}\OperatorTok{)}
\OperatorTok{\{}
    \AttributeTok{const} \DataTypeTok{double}\NormalTok{ alpha }\OperatorTok{=}\NormalTok{ absorptionOpacity }\OperatorTok{*} \KeywordTok{this}\OperatorTok{{-}\textgreater{}}\NormalTok{factorFleck}\OperatorTok{[}\NormalTok{cellIndex}\OperatorTok{]} \OperatorTok{*}\NormalTok{ units}\OperatorTok{::}\NormalTok{clight}\OperatorTok{;}
    \KeywordTok{this}\OperatorTok{{-}\textgreater{}}\NormalTok{EmitComptonResidualChildren}\OperatorTok{(}
\NormalTok{        cellIndex}\OperatorTok{,}
\NormalTok{        cell}\OperatorTok{,}
\NormalTok{        segmentGroup}\OperatorTok{,}
\NormalTok{        segmentStartLocation}\OperatorTok{,}
\NormalTok{        particle}\OperatorTok{.}\NormalTok{velocity}\OperatorTok{,}
\NormalTok{        segmentStartTimeLeft}\OperatorTok{,}
\NormalTok{        segmentStartWeight}\OperatorTok{,}
\NormalTok{        alpha}\OperatorTok{,}
\NormalTok{        dt}\OperatorTok{,}
\NormalTok{        dopplerShift}\OperatorTok{,}
\NormalTok{        particlesToAdd}\OperatorTok{);}
\OperatorTok{\}}
\end{Highlighting}
\end{Shaded}

Implement:

\begin{Shaded}
\begin{Highlighting}[]
\DataTypeTok{void}\NormalTok{ RadiationIMC}\OperatorTok{::}\NormalTok{EmitComptonResidualChildren}\OperatorTok{(}
    \DataTypeTok{size\_t}\NormalTok{ cellIndex}\OperatorTok{,}
    \AttributeTok{const}\NormalTok{ ComputationalCell3D}\OperatorTok{\&}\NormalTok{ cell}\OperatorTok{,}
    \DataTypeTok{size\_t}\NormalTok{ sourceGroup}\OperatorTok{,}
    \AttributeTok{const}\NormalTok{ Vector3D}\OperatorTok{\&}\NormalTok{ startLocation}\OperatorTok{,}
    \AttributeTok{const}\NormalTok{ Vector3D}\OperatorTok{\&}\NormalTok{ velocity}\OperatorTok{,}
    \DataTypeTok{double}\NormalTok{ startTimeLeft}\OperatorTok{,}
    \DataTypeTok{double}\NormalTok{ signedWeightAtStart}\OperatorTok{,}
    \DataTypeTok{double}\NormalTok{ alpha}\OperatorTok{,}
    \DataTypeTok{double}\NormalTok{ dt}\OperatorTok{,}
    \DataTypeTok{double}\NormalTok{ dopplerShift}\OperatorTok{,}
    \BuiltInTok{std::}\NormalTok{vector}\OperatorTok{\textless{}}\NormalTok{Particle}\OperatorTok{\textgreater{}\&}\NormalTok{ particlesToAdd}\OperatorTok{)}
\OperatorTok{\{}
    \AttributeTok{const}\NormalTok{ ComptonCellData}\OperatorTok{\&}\NormalTok{ cd }\OperatorTok{=} \KeywordTok{this}\OperatorTok{{-}\textgreater{}}\NormalTok{comptonData}\OperatorTok{[}\NormalTok{cellIndex}\OperatorTok{];}
    \ControlFlowTok{if}\OperatorTok{(!}\NormalTok{cd}\OperatorTok{.}\NormalTok{active}\OperatorTok{)}
        \ControlFlowTok{return}\OperatorTok{;}

    \AttributeTok{const} \DataTypeTok{double}\NormalTok{ integratedWdt }\OperatorTok{=}\NormalTok{ IntegratedWeightOverDt}\OperatorTok{(}\NormalTok{signedWeightAtStart}\OperatorTok{,}\NormalTok{ alpha}\OperatorTok{,}\NormalTok{ dt}\OperatorTok{);}
    \AttributeTok{const} \DataTypeTok{double}\NormalTok{ rateFrameFactor }\OperatorTok{=} \OperatorTok{(}\KeywordTok{this}\OperatorTok{{-}\textgreater{}}\NormalTok{withHydro }\OperatorTok{\&\&} \OperatorTok{!}\KeywordTok{this}\OperatorTok{{-}\textgreater{}}\NormalTok{MMC}\OperatorTok{)} \OperatorTok{?}\NormalTok{ dopplerShift }\OperatorTok{:} \FloatTok{1.0}\OperatorTok{;}
    \AttributeTok{const} \DataTypeTok{double}\NormalTok{ pathEnergy }\OperatorTok{=}\NormalTok{ units}\OperatorTok{::}\NormalTok{clight }\OperatorTok{*}\NormalTok{ rateFrameFactor }\OperatorTok{*}\NormalTok{ integratedWdt}\OperatorTok{;}

    \AttributeTok{const} \DataTypeTok{double}\NormalTok{ rpos }\OperatorTok{=}\NormalTok{ cd}\OperatorTok{.}\NormalTok{RposRow}\OperatorTok{[}\NormalTok{sourceGroup}\OperatorTok{];}
    \ControlFlowTok{if}\OperatorTok{(}\NormalTok{rpos }\OperatorTok{\textgreater{}} \FloatTok{0.0}\OperatorTok{)}
    \OperatorTok{\{}
        \DataTypeTok{size\_t}\NormalTok{ target }\OperatorTok{=}\NormalTok{ SampleFromCdf}\OperatorTok{(}\NormalTok{cd}\OperatorTok{.}\NormalTok{RposCdf}\OperatorTok{[}\NormalTok{sourceGroup}\OperatorTok{],} \KeywordTok{this}\OperatorTok{{-}\textgreater{}}\NormalTok{dist}\OperatorTok{(}\KeywordTok{this}\OperatorTok{{-}\textgreater{}}\NormalTok{re}\OperatorTok{));}
        \AttributeTok{const} \DataTypeTok{double}\NormalTok{ tauBirth }\OperatorTok{=}\NormalTok{ SampleBirthTimeOnSegment}\OperatorTok{(}\NormalTok{alpha}\OperatorTok{,}\NormalTok{ dt}\OperatorTok{,} \KeywordTok{this}\OperatorTok{{-}\textgreater{}}\NormalTok{dist}\OperatorTok{(}\KeywordTok{this}\OperatorTok{{-}\textgreater{}}\NormalTok{re}\OperatorTok{));}
        \AttributeTok{const}\NormalTok{ Vector3D xBirth }\OperatorTok{=}\NormalTok{ startLocation }\OperatorTok{+}\NormalTok{ velocity }\OperatorTok{*}\NormalTok{ tauBirth}\OperatorTok{;}
        \AttributeTok{const} \DataTypeTok{double}\NormalTok{ childTimeLeft }\OperatorTok{=}\NormalTok{ startTimeLeft }\OperatorTok{{-}}\NormalTok{ tauBirth}\OperatorTok{;}
        \AttributeTok{const} \DataTypeTok{double}\NormalTok{ signedEnergy }\OperatorTok{=}\NormalTok{ pathEnergy }\OperatorTok{*}\NormalTok{ rpos}\OperatorTok{;}
\NormalTok{        BankSignedChildPacket}\OperatorTok{(}\NormalTok{cellIndex}\OperatorTok{,}\NormalTok{ cell}\OperatorTok{,}\NormalTok{ xBirth}\OperatorTok{,}\NormalTok{ velocity}\OperatorTok{,}\NormalTok{ childTimeLeft}\OperatorTok{,}\NormalTok{ target}\OperatorTok{,}\NormalTok{ signedEnergy}\OperatorTok{,}\NormalTok{ particlesToAdd}\OperatorTok{);}
    \OperatorTok{\}}

    \AttributeTok{const} \DataTypeTok{double}\NormalTok{ rneg }\OperatorTok{=}\NormalTok{ cd}\OperatorTok{.}\NormalTok{RnegRow}\OperatorTok{[}\NormalTok{sourceGroup}\OperatorTok{];}
    \ControlFlowTok{if}\OperatorTok{(}\NormalTok{rneg }\OperatorTok{\textgreater{}} \FloatTok{0.0}\OperatorTok{)}
    \OperatorTok{\{}
        \DataTypeTok{size\_t}\NormalTok{ target }\OperatorTok{=}\NormalTok{ SampleFromCdf}\OperatorTok{(}\NormalTok{cd}\OperatorTok{.}\NormalTok{RnegCdf}\OperatorTok{[}\NormalTok{sourceGroup}\OperatorTok{],} \KeywordTok{this}\OperatorTok{{-}\textgreater{}}\NormalTok{dist}\OperatorTok{(}\KeywordTok{this}\OperatorTok{{-}\textgreater{}}\NormalTok{re}\OperatorTok{));}
        \AttributeTok{const} \DataTypeTok{double}\NormalTok{ tauBirth }\OperatorTok{=}\NormalTok{ SampleBirthTimeOnSegment}\OperatorTok{(}\NormalTok{alpha}\OperatorTok{,}\NormalTok{ dt}\OperatorTok{,} \KeywordTok{this}\OperatorTok{{-}\textgreater{}}\NormalTok{dist}\OperatorTok{(}\KeywordTok{this}\OperatorTok{{-}\textgreater{}}\NormalTok{re}\OperatorTok{));}
        \AttributeTok{const}\NormalTok{ Vector3D xBirth }\OperatorTok{=}\NormalTok{ startLocation }\OperatorTok{+}\NormalTok{ velocity }\OperatorTok{*}\NormalTok{ tauBirth}\OperatorTok{;}
        \AttributeTok{const} \DataTypeTok{double}\NormalTok{ childTimeLeft }\OperatorTok{=}\NormalTok{ startTimeLeft }\OperatorTok{{-}}\NormalTok{ tauBirth}\OperatorTok{;}
        \AttributeTok{const} \DataTypeTok{double}\NormalTok{ signedEnergy }\OperatorTok{=} \OperatorTok{{-}}\NormalTok{pathEnergy }\OperatorTok{*}\NormalTok{ rneg}\OperatorTok{;}
\NormalTok{        BankSignedChildPacket}\OperatorTok{(}\NormalTok{cellIndex}\OperatorTok{,}\NormalTok{ cell}\OperatorTok{,}\NormalTok{ xBirth}\OperatorTok{,}\NormalTok{ velocity}\OperatorTok{,}\NormalTok{ childTimeLeft}\OperatorTok{,}\NormalTok{ target}\OperatorTok{,}\NormalTok{ signedEnergy}\OperatorTok{,}\NormalTok{ particlesToAdd}\OperatorTok{);}
    \OperatorTok{\}}
\OperatorTok{\}}
\end{Highlighting}
\end{Shaded}

Because \texttt{pathEnergy} is signed when the parent packet is signed,
this automatically gives the correct sign behavior:

\begin{itemize}
\tightlist
\item
  Positive parent through \texttt{Rpos} creates positive child.
\item
  Positive parent through \texttt{Rneg} creates negative child.
\item
  Negative parent through \texttt{Rpos} creates negative child.
\item
  Negative parent through \texttt{Rneg} creates positive child.
\end{itemize}

\subsubsection{7.6 Continuous absorption and time-average
tallies}\label{continuous-absorption-and-time-average-tallies}

Current lines 121-145 mostly work for signed packets. However, guard
division by zero in the time-average estimator:

\begin{Shaded}
\begin{Highlighting}[]
\DataTypeTok{double}\NormalTok{ tmp2 }\OperatorTok{=}\NormalTok{ absorptionOpacity }\OperatorTok{*} \KeywordTok{this}\OperatorTok{{-}\textgreater{}}\NormalTok{factorFleck}\OperatorTok{[}\NormalTok{cellIndex}\OperatorTok{]} \OperatorTok{*}\NormalTok{ units}\OperatorTok{::}\NormalTok{clight}\OperatorTok{;}
\DataTypeTok{double}\NormalTok{ tmp }\OperatorTok{=} \OperatorTok{{-}}\NormalTok{dt }\OperatorTok{*}\NormalTok{ tmp2}\OperatorTok{;}
\DataTypeTok{double}\NormalTok{ expFactor1 }\OperatorTok{=} \BuiltInTok{std::}\NormalTok{expm1}\OperatorTok{(}\NormalTok{tmp }\OperatorTok{*}\NormalTok{ dopplerShift}\OperatorTok{);}
\DataTypeTok{double}\NormalTok{ expFactor2 }\OperatorTok{=} \BuiltInTok{std::}\NormalTok{expm1}\OperatorTok{(}\NormalTok{tmp}\OperatorTok{);}

\CommentTok{// Integrated weight used for time{-}average. Avoid division by zero.}
\DataTypeTok{double}\NormalTok{ integratedForTally}\OperatorTok{;}
\ControlFlowTok{if}\OperatorTok{(}\BuiltInTok{std::}\NormalTok{abs}\OperatorTok{(}\NormalTok{tmp2 }\OperatorTok{*}\NormalTok{ dt}\OperatorTok{)} \OperatorTok{\textless{}} \FloatTok{1e{-}12}\OperatorTok{)}
\NormalTok{    integratedForTally }\OperatorTok{=}\NormalTok{ particle}\OperatorTok{.}\NormalTok{weight }\OperatorTok{*}\NormalTok{ dt}\OperatorTok{;}
\ControlFlowTok{else}
\NormalTok{    integratedForTally }\OperatorTok{=}\NormalTok{ particle}\OperatorTok{.}\NormalTok{weight }\OperatorTok{*}\NormalTok{ expFactor2 }\OperatorTok{*} \OperatorTok{({-}}\FloatTok{1.0} \OperatorTok{/}\NormalTok{ tmp2}\OperatorTok{);}
\end{Highlighting}
\end{Shaded}

Then use \texttt{integratedForTally} instead of repeating the expression
at lines 137 and 141.

For signed packets, material deposition line 128 is already correct:

\begin{Shaded}
\begin{Highlighting}[]
\KeywordTok{this}\OperatorTok{{-}\textgreater{}}\NormalTok{conserved}\OperatorTok{[}\NormalTok{cellIndex}\OperatorTok{].}\NormalTok{internal\_energy }\OperatorTok{+=} \OperatorTok{{-}}\NormalTok{expFactor2 }\OperatorTok{*}\NormalTok{ particle}\OperatorTok{.}\NormalTok{weight}\OperatorTok{;}
\end{Highlighting}
\end{Shaded}

If \texttt{particle.weight} is negative, absorption of an anti-packet
removes material energy, as it should.

\subsubsection{7.7 Low-weight removal}\label{low-weight-removal}

Change line 146 to use absolute value:

\begin{Shaded}
\begin{Highlighting}[]
\ControlFlowTok{if}\OperatorTok{(}\BuiltInTok{std::}\NormalTok{abs}\OperatorTok{(}\NormalTok{particle}\OperatorTok{.}\NormalTok{weight}\OperatorTok{)} \OperatorTok{\textless{}}\NormalTok{ particle}\OperatorTok{.}\NormalTok{initialWeight }\OperatorTok{*} \FloatTok{1e{-}3}\OperatorTok{)}
\end{Highlighting}
\end{Shaded}

When removing the remaining signed packet at lines 149-152, the existing
material update is correct:

\begin{Shaded}
\begin{Highlighting}[]
\KeywordTok{this}\OperatorTok{{-}\textgreater{}}\NormalTok{conserved}\OperatorTok{[}\NormalTok{cellIndex}\OperatorTok{].}\NormalTok{internal\_energy }\OperatorTok{+=}\NormalTok{ particle}\OperatorTok{.}\NormalTok{weight}\OperatorTok{;}
\end{Highlighting}
\end{Shaded}

For a negative packet, this subtracts material energy, which cancels the
remaining anti-radiation.

\subsubsection{7.8 Event selection}\label{event-selection}

Replace the current effective-scatter selection at lines 161-177 with
explicit channel selection.

Pseudo-code:

\begin{Shaded}
\begin{Highlighting}[]
\ControlFlowTok{else} \ControlFlowTok{if}\OperatorTok{(}\NormalTok{min}\OperatorTok{.}\NormalTok{first }\OperatorTok{==}\NormalTok{ Events}\OperatorTok{::}\NormalTok{SCATTERING}\OperatorTok{)}
\OperatorTok{\{}
\NormalTok{    Vector3D oldVelocity }\OperatorTok{=}\NormalTok{ particle}\OperatorTok{.}\NormalTok{velocity}\OperatorTok{;}
    \AttributeTok{const} \DataTypeTok{double}\NormalTok{ oldWeight }\OperatorTok{=}\NormalTok{ particle}\OperatorTok{.}\NormalTok{weight}\OperatorTok{;}
    \AttributeTok{const} \DataTypeTok{double}\NormalTok{ D\_lab\_to\_co }\OperatorTok{=}\NormalTok{ dopplerShift}\OperatorTok{;}

    \AttributeTok{const} \DataTypeTok{double}\NormalTok{ total }\OperatorTok{=}\NormalTok{ elasticScatteringOpacity }\OperatorTok{+}\NormalTok{ effectiveAbsorptionOpacity }\OperatorTok{+}\NormalTok{ directComptonOpacity}\OperatorTok{;}
    \DataTypeTok{double}\NormalTok{ r }\OperatorTok{=} \KeywordTok{this}\OperatorTok{{-}\textgreater{}}\NormalTok{dist}\OperatorTok{(}\KeywordTok{this}\OperatorTok{{-}\textgreater{}}\NormalTok{re}\OperatorTok{)} \OperatorTok{*}\NormalTok{ total}\OperatorTok{;}

    \ControlFlowTok{if}\OperatorTok{(}\NormalTok{r }\OperatorTok{\textless{}}\NormalTok{ elasticScatteringOpacity}\OperatorTok{)}
    \OperatorTok{\{}
        \CommentTok{// Existing elastic scattering.}
\NormalTok{        particle}\OperatorTok{.}\NormalTok{velocity }\OperatorTok{=}\NormalTok{ opacity}\OperatorTok{{-}\textgreater{}}\NormalTok{getNewScatterVelocity}\OperatorTok{(}\NormalTok{cell}\OperatorTok{,}\NormalTok{ particle}\OperatorTok{);}
    \OperatorTok{\}}
    \ControlFlowTok{else} \ControlFlowTok{if}\OperatorTok{((}\NormalTok{r }\OperatorTok{{-}=}\NormalTok{ elasticScatteringOpacity}\OperatorTok{)} \OperatorTok{\textless{}}\NormalTok{ effectiveAbsorptionOpacity}\OperatorTok{)}
    \OperatorTok{\{}
        \CommentTok{// Existing ordinary IMC effective absorption scattering.}
\NormalTok{        particle}\OperatorTok{.}\NormalTok{velocity }\OperatorTok{=}\NormalTok{ opacity}\OperatorTok{{-}\textgreater{}}\NormalTok{getNewScatterVelocity}\OperatorTok{(}\NormalTok{cell}\OperatorTok{,}\NormalTok{ particle}\OperatorTok{);}
        \ControlFlowTok{if}\OperatorTok{(}\KeywordTok{this}\OperatorTok{{-}\textgreater{}}\NormalTok{multigroupOpacity}\OperatorTok{)}
        \OperatorTok{\{}
\NormalTok{            particle}\OperatorTok{.}\NormalTok{frequency }\OperatorTok{*=}\NormalTok{ dopplerShift}\OperatorTok{;} \CommentTok{// lab {-}\textgreater{} comoving}
\NormalTok{            ClampFrequencyToBounds}\OperatorTok{(}\NormalTok{particle}\OperatorTok{.}\NormalTok{frequency}\OperatorTok{);}

            \DataTypeTok{size\_t}\NormalTok{ targetGroup}\OperatorTok{;}
            \ControlFlowTok{if}\OperatorTok{(}\KeywordTok{this}\OperatorTok{{-}\textgreater{}}\NormalTok{withCompton}\OperatorTok{)}
\NormalTok{                targetGroup }\OperatorTok{=}\NormalTok{ SampleFromCdf}\OperatorTok{(}\KeywordTok{this}\OperatorTok{{-}\textgreater{}}\NormalTok{comptonData}\OperatorTok{[}\NormalTok{cellIndex}\OperatorTok{].}\NormalTok{qBaseCdf}\OperatorTok{,} \KeywordTok{this}\OperatorTok{{-}\textgreater{}}\NormalTok{dist}\OperatorTok{(}\KeywordTok{this}\OperatorTok{{-}\textgreater{}}\NormalTok{re}\OperatorTok{));}
            \ControlFlowTok{else}
\NormalTok{                targetGroup }\OperatorTok{=} \BuiltInTok{std::}\NormalTok{numeric\_limits}\OperatorTok{\textless{}}\DataTypeTok{size\_t}\OperatorTok{\textgreater{}::}\NormalTok{max}\OperatorTok{();}

            \ControlFlowTok{if}\OperatorTok{(}\KeywordTok{this}\OperatorTok{{-}\textgreater{}}\NormalTok{withCompton }\OperatorTok{\&\&}\NormalTok{ targetGroup }\OperatorTok{!=} \BuiltInTok{std::}\NormalTok{numeric\_limits}\OperatorTok{\textless{}}\DataTypeTok{size\_t}\OperatorTok{\textgreater{}::}\NormalTok{max}\OperatorTok{())}
\NormalTok{                particle}\OperatorTok{.}\NormalTok{frequency }\OperatorTok{=}\NormalTok{ FrequencyForGroup}\OperatorTok{(}\NormalTok{targetGroup}\OperatorTok{);}
            \ControlFlowTok{else}
\NormalTok{                particle}\OperatorTok{.}\NormalTok{frequency }\OperatorTok{=} \KeywordTok{this}\OperatorTok{{-}\textgreater{}}\NormalTok{multigroupOpacity}\OperatorTok{{-}\textgreater{}}\NormalTok{GetThermalEnergy}\OperatorTok{(}\NormalTok{cell}\OperatorTok{,} \KeywordTok{this}\OperatorTok{{-}\textgreater{}}\NormalTok{dist}\OperatorTok{(}\KeywordTok{this}\OperatorTok{{-}\textgreater{}}\NormalTok{re}\OperatorTok{));}
        \OperatorTok{\}}
    \OperatorTok{\}}
    \ControlFlowTok{else}
    \OperatorTok{\{}
        \CommentTok{// Direct analog Compton event. Only possible if comptonUseAnalogDirect is true.}
        \KeywordTok{this}\OperatorTok{{-}\textgreater{}}\NormalTok{ApplyAnalogComptonEvent}\OperatorTok{(}\NormalTok{cellIndex}\OperatorTok{,}\NormalTok{ cell}\OperatorTok{,}\NormalTok{ group}\OperatorTok{,}\NormalTok{ oldVelocity}\OperatorTok{,}\NormalTok{ oldWeight}\OperatorTok{,}\NormalTok{ D\_lab\_to\_co}\OperatorTok{,}\NormalTok{ particle}\OperatorTok{);}
    \OperatorTok{\}}

    \CommentTok{// Existing hydro Lorentz/momentum block, adjusted if ApplyAnalogComptonEvent already updated momentum.}
\OperatorTok{\}}
\end{Highlighting}
\end{Shaded}

Be careful not to double-apply momentum/energy exchange for direct
Compton. For the first implementation with
\texttt{comptonUseAnalogDirect=false}, this branch is not used.

\subsubsection{7.9 Direct analog Compton event
implementation}\label{direct-analog-compton-event-implementation}

Add only after residual-only mode works.

\begin{Shaded}
\begin{Highlighting}[]
\DataTypeTok{void}\NormalTok{ RadiationIMC}\OperatorTok{::}\NormalTok{ApplyAnalogComptonEvent}\OperatorTok{(}
    \DataTypeTok{size\_t}\NormalTok{ cellIndex}\OperatorTok{,}
    \AttributeTok{const}\NormalTok{ ComputationalCell3D}\OperatorTok{\&}\NormalTok{ cell}\OperatorTok{,}
    \DataTypeTok{size\_t}\NormalTok{ sourceGroup}\OperatorTok{,}
    \AttributeTok{const}\NormalTok{ Vector3D}\OperatorTok{\&}\NormalTok{ oldVelocity}\OperatorTok{,}
    \DataTypeTok{double}\NormalTok{ oldWeight}\OperatorTok{,}
    \DataTypeTok{double}\NormalTok{ dopplerShift}\OperatorTok{,}
\NormalTok{    Particle}\OperatorTok{\&}\NormalTok{ particle}\OperatorTok{)}
\OperatorTok{\{}
    \AttributeTok{const}\NormalTok{ ComptonCellData}\OperatorTok{\&}\NormalTok{ cd }\OperatorTok{=} \KeywordTok{this}\OperatorTok{{-}\textgreater{}}\NormalTok{comptonData}\OperatorTok{[}\NormalTok{cellIndex}\OperatorTok{];}
    \AttributeTok{const} \DataTypeTok{size\_t}\NormalTok{ targetGroup }\OperatorTok{=}\NormalTok{ SampleFromCdf}\OperatorTok{(}\NormalTok{cd}\OperatorTok{.}\NormalTok{directCdf}\OperatorTok{[}\NormalTok{sourceGroup}\OperatorTok{],} \KeywordTok{this}\OperatorTok{{-}\textgreater{}}\NormalTok{dist}\OperatorTok{(}\KeywordTok{this}\OperatorTok{{-}\textgreater{}}\NormalTok{re}\OperatorTok{));}
    \ControlFlowTok{if}\OperatorTok{(}\NormalTok{targetGroup }\OperatorTok{==} \BuiltInTok{std::}\NormalTok{numeric\_limits}\OperatorTok{\textless{}}\DataTypeTok{size\_t}\OperatorTok{\textgreater{}::}\NormalTok{max}\OperatorTok{())}
        \ControlFlowTok{return}\OperatorTok{;}

    \CommentTok{// New direction. If a Klein{-}Nishina angular sampler is available, use it.}
    \CommentTok{// Otherwise use the existing scattering direction. The diffusion approximation does not retain the angular detail.}
\NormalTok{    particle}\OperatorTok{.}\NormalTok{velocity }\OperatorTok{=} \KeywordTok{this}\OperatorTok{{-}\textgreater{}}\NormalTok{opacity}\OperatorTok{{-}\textgreater{}}\NormalTok{getNewScatterVelocity}\OperatorTok{(}\NormalTok{cell}\OperatorTok{,}\NormalTok{ particle}\OperatorTok{);}

    \AttributeTok{const} \DataTypeTok{double}\NormalTok{ oldCenter }\OperatorTok{=}\NormalTok{ energyGroupCenters}\OperatorTok{[}\NormalTok{sourceGroup}\OperatorTok{];}
    \AttributeTok{const} \DataTypeTok{double}\NormalTok{ newCenter }\OperatorTok{=}\NormalTok{ energyGroupCenters}\OperatorTok{[}\NormalTok{targetGroup}\OperatorTok{];}
    \AttributeTok{const} \DataTypeTok{double}\NormalTok{ ratio }\OperatorTok{=}\NormalTok{ newCenter }\OperatorTok{/}\NormalTok{ oldCenter}\OperatorTok{;}

    \AttributeTok{const} \DataTypeTok{double}\NormalTok{ newWeight }\OperatorTok{=}\NormalTok{ oldWeight }\OperatorTok{*}\NormalTok{ ratio}\OperatorTok{;}
    \AttributeTok{const} \DataTypeTok{double}\NormalTok{ materialDeposit }\OperatorTok{=}\NormalTok{ oldWeight }\OperatorTok{{-}}\NormalTok{ newWeight}\OperatorTok{;}

\NormalTok{    particle}\OperatorTok{.}\NormalTok{weight }\OperatorTok{=}\NormalTok{ newWeight}\OperatorTok{;}
\NormalTok{    particle}\OperatorTok{.}\NormalTok{frequency }\OperatorTok{=}\NormalTok{ FrequencyForGroup}\OperatorTok{(}\NormalTok{targetGroup}\OperatorTok{);}

    \ControlFlowTok{if}\OperatorTok{(!}\KeywordTok{this}\OperatorTok{{-}\textgreater{}}\NormalTok{noHydroFeedback}\OperatorTok{)}
    \OperatorTok{\{}
        \KeywordTok{this}\OperatorTok{{-}\textgreater{}}\NormalTok{conserved}\OperatorTok{[}\NormalTok{cellIndex}\OperatorTok{].}\NormalTok{internal\_energy }\OperatorTok{+=}\NormalTok{ materialDeposit}\OperatorTok{;}
        \ControlFlowTok{if}\OperatorTok{(}\KeywordTok{this}\OperatorTok{{-}\textgreater{}}\NormalTok{withHydro }\OperatorTok{\&\&} \OperatorTok{!}\KeywordTok{this}\OperatorTok{{-}\textgreater{}}\NormalTok{diffusionPressureGradient}\OperatorTok{)}
        \OperatorTok{\{}
            \KeywordTok{this}\OperatorTok{{-}\textgreater{}}\NormalTok{conserved}\OperatorTok{[}\NormalTok{cellIndex}\OperatorTok{].}\NormalTok{momentum }\OperatorTok{+=}
                \OperatorTok{(}\NormalTok{oldWeight }\OperatorTok{*}\NormalTok{ oldVelocity }\OperatorTok{{-}}\NormalTok{ newWeight }\OperatorTok{*}\NormalTok{ particle}\OperatorTok{.}\NormalTok{velocity}\OperatorTok{)} \OperatorTok{*}\NormalTok{ units}\OperatorTok{::}\NormalTok{inv\_clight2}\OperatorTok{;}
        \OperatorTok{\}}
    \OperatorTok{\}}
\OperatorTok{\}}
\end{Highlighting}
\end{Shaded}

With hydro and Lorentz transformations, this needs careful testing.
Initially run Compton tests with \texttt{withHydro=false} or
\texttt{MMC=true}. Once those pass, match the existing line 178-190
transformation convention.

\subsection{8. Random walk interaction}\label{random-walk-interaction}

The current \texttt{step} tries random walk before ordinary transport at
lines 56-63. A random-walk step would skip residual source emission
unless the random-walk estimator is extended.

For the first correct implementation:

\begin{itemize}
\tightlist
\item
  If \texttt{withCompton\ \&\&\ comptonDisableRandomWalk}, skip random
  walk entirely.
\item
  Or mark a cell ineligible for random walk if any of these are nonzero:

  \begin{itemize}
  \tightlist
  \item
    \texttt{sum\_h\ RposRow{[}h{]}}
  \item
    \texttt{sum\_h\ RnegRow{[}h{]}}
  \item
    direct Compton out rate
  \item
    significant \texttt{Btotal} signed source
  \end{itemize}
\end{itemize}

Implementation in \texttt{preStep} after
\texttt{precomputeRandomWalkData()}:

\begin{Shaded}
\begin{Highlighting}[]
\ControlFlowTok{if}\OperatorTok{(}\KeywordTok{this}\OperatorTok{{-}\textgreater{}}\NormalTok{withCompton }\OperatorTok{\&\&} \KeywordTok{this}\OperatorTok{{-}\textgreater{}}\NormalTok{comptonDisableRandomWalk}\OperatorTok{)}
\OperatorTok{\{}
    \BuiltInTok{std::}\NormalTok{fill}\OperatorTok{(}\KeywordTok{this}\OperatorTok{{-}\textgreater{}}\NormalTok{rwCellEligible}\OperatorTok{.}\NormalTok{begin}\OperatorTok{(),} \KeywordTok{this}\OperatorTok{{-}\textgreater{}}\NormalTok{rwCellEligible}\OperatorTok{.}\NormalTok{end}\OperatorTok{(),} \KeywordTok{false}\OperatorTok{);}
\OperatorTok{\}}
\end{Highlighting}
\end{Shaded}

Later, if random walk is needed with Compton, add a path-length
estimator inside \texttt{tryRandomWalkStep} and emit residual children
based on the expected path length of the random walk. Do not enable
random walk with Compton before that estimator exists.

\subsection{9. Post-step and
cancellation}\label{post-step-and-cancellation}

\subsubsection{9.1 Signed tallies}\label{signed-tallies}

\texttt{postStep} already sums \texttt{particle.weight} into
\texttt{conserved.Erad} and \texttt{conserved.Eg} at lines 315-324. This
is correct for signed packets.

Add diagnostics after the tally:

\begin{Shaded}
\begin{Highlighting}[]
\ControlFlowTok{for}\OperatorTok{(}\DataTypeTok{size\_t}\NormalTok{ i }\OperatorTok{=} \DecValTok{0}\OperatorTok{;}\NormalTok{ i }\OperatorTok{\textless{}}\NormalTok{ Ncells}\OperatorTok{;} \OperatorTok{++}\NormalTok{i}\OperatorTok{)}
\OperatorTok{\{}
    \ControlFlowTok{if}\OperatorTok{(}\KeywordTok{this}\OperatorTok{{-}\textgreater{}}\NormalTok{conserved}\OperatorTok{[}\NormalTok{i}\OperatorTok{].}\NormalTok{Erad }\OperatorTok{\textless{}} \OperatorTok{{-}}\NormalTok{radiationNegativeTolerance}\OperatorTok{)}
    \OperatorTok{\{}
\NormalTok{        UniversalError eo}\OperatorTok{(}\StringTok{"Negative radiation energy after signed Compton IMC tally"}\OperatorTok{);}
\NormalTok{        eo}\OperatorTok{.}\NormalTok{addEntry}\OperatorTok{(}\StringTok{"cell"}\OperatorTok{,}\NormalTok{ i}\OperatorTok{);}
\NormalTok{        eo}\OperatorTok{.}\NormalTok{addEntry}\OperatorTok{(}\StringTok{"Erad"}\OperatorTok{,} \KeywordTok{this}\OperatorTok{{-}\textgreater{}}\NormalTok{conserved}\OperatorTok{[}\NormalTok{i}\OperatorTok{].}\NormalTok{Erad}\OperatorTok{);}
        \ControlFlowTok{throw}\NormalTok{ eo}\OperatorTok{;}
    \OperatorTok{\}}
    \ControlFlowTok{if}\OperatorTok{(}\KeywordTok{this}\OperatorTok{{-}\textgreater{}}\NormalTok{multigroupOpacity}\OperatorTok{)}
    \OperatorTok{\{}
        \ControlFlowTok{for}\OperatorTok{(}\DataTypeTok{size\_t}\NormalTok{ g }\OperatorTok{=} \DecValTok{0}\OperatorTok{;}\NormalTok{ g }\OperatorTok{\textless{}}\NormalTok{ ENERGY\_GROUPS\_NUM}\OperatorTok{;} \OperatorTok{++}\NormalTok{g}\OperatorTok{)}
        \OperatorTok{\{}
            \ControlFlowTok{if}\OperatorTok{(}\KeywordTok{this}\OperatorTok{{-}\textgreater{}}\NormalTok{conserved}\OperatorTok{[}\NormalTok{i}\OperatorTok{].}\NormalTok{Eg}\OperatorTok{[}\NormalTok{g}\OperatorTok{]} \OperatorTok{\textless{}} \OperatorTok{{-}}\NormalTok{groupRadiationNegativeTolerance}\OperatorTok{)}
            \OperatorTok{\{}
                \CommentTok{// Throw or log; do not silently clip in production tests.}
            \OperatorTok{\}}
        \OperatorTok{\}}
    \OperatorTok{\}}
\OperatorTok{\}}
\end{Highlighting}
\end{Shaded}

Use tolerances proportional to total absolute packet energy in the cell,
not a hard-coded absolute number.

\subsubsection{9.2 Packet cancellation}\label{packet-cancellation}

Signed residual packets can increase variance. Add a cancellation pass
before final tallies or at census.

Primitive cancellation:

\begin{enumerate}
\def\labelenumi{\arabic{enumi}.}
\tightlist
\item
  Group packets by \texttt{(cellIndex,\ group)}.
\item
  Sum signed weights in each bin.
\item
  Replace all packets in the bin by one packet if angular information is
  not important, or by a small fixed number of packets with
  representative directions.
\end{enumerate}

Safer cancellation:

\begin{enumerate}
\def\labelenumi{\arabic{enumi}.}
\tightlist
\item
  Separate positive and negative packets in the same
  \texttt{(cellIndex,\ group)}.
\item
  Sort by location or distance to cell center.
\item
  Cancel pairs by reducing the smaller absolute weight.
\item
  Remove zero-weight packets.
\end{enumerate}

Do not alter material energy during radiation packet cancellation.
Positive and negative radiation packets cancel each other; material was
already updated when they were created or absorbed.

Add cancellation only after the basic scheme passes conservation tests,
because cancellation can hide bugs.

\subsection{10. Conservation rules}\label{conservation-rules}

Every code path must obey these signed conservation rules.

\subsubsection{10.1 Material source
emission}\label{material-source-emission}

When creating a signed radiation source packet with energy
\texttt{Wemit}:

\begin{Shaded}
\begin{Highlighting}[]
\NormalTok{material\_internal\_energy }\OperatorTok{{-}=}\NormalTok{ Wemit}\OperatorTok{;}
\end{Highlighting}
\end{Shaded}

If \texttt{Wemit\ \textgreater{}\ 0}, material loses energy. If
\texttt{Wemit\ \textless{}\ 0}, material gains energy.

\subsubsection{10.2 Continuous absorption}\label{continuous-absorption}

Current line 128:

\begin{Shaded}
\begin{Highlighting}[]
\NormalTok{material\_internal\_energy }\OperatorTok{+=} \OperatorTok{{-}}\NormalTok{expm1}\OperatorTok{({-}}\NormalTok{alpha }\OperatorTok{*}\NormalTok{ dt}\OperatorTok{)} \OperatorTok{*}\NormalTok{ particle}\OperatorTok{.}\NormalTok{weight}\OperatorTok{;}
\end{Highlighting}
\end{Shaded}

This is correct for signed packets. Absorbing a negative packet removes
material energy.

\subsubsection{10.3 Low-weight packet
removal}\label{low-weight-packet-removal}

When a remaining signed packet is killed and deposited locally:

\begin{Shaded}
\begin{Highlighting}[]
\NormalTok{material\_internal\_energy }\OperatorTok{+=}\NormalTok{ particle}\OperatorTok{.}\NormalTok{weight}\OperatorTok{;}
\end{Highlighting}
\end{Shaded}

This is already done at line 151. It is correct for signed packets.

\subsubsection{10.4 Residual child packet
emission}\label{residual-child-packet-emission}

When \texttt{EmitComptonResidualChildren} creates a child of signed
energy \texttt{Wchild}:

\begin{Shaded}
\begin{Highlighting}[]
\NormalTok{material\_internal\_energy }\OperatorTok{{-}=}\NormalTok{ Wchild}\OperatorTok{;}
\end{Highlighting}
\end{Shaded}

Do this in \texttt{BankSignedChildPacket}.

\subsubsection{10.5 Physical analog Compton
event}\label{physical-analog-compton-event}

If an analog Compton event changes a packet from signed energy
\texttt{Wold} to \texttt{Wnew}:

\begin{Shaded}
\begin{Highlighting}[]
\NormalTok{material\_internal\_energy }\OperatorTok{+=}\NormalTok{ Wold }\OperatorTok{{-}}\NormalTok{ Wnew}\OperatorTok{;}
\end{Highlighting}
\end{Shaded}

For a positive photon that upscatters,
\texttt{Wnew\ \textgreater{}\ Wold}, so the material loses energy. For a
positive photon that downscatters, the material gains energy. For an
anti-packet, the signs reverse automatically.

\subsection{11. Unit conventions and consistency
traps}\label{unit-conventions-and-consistency-traps}

The current IMC code appears to use unscaled physical units directly,
while \texttt{MultigroupDiffusion.cpp} contains many
\texttt{mass\_scale\_}, \texttt{length\_scale\_}, and
\texttt{time\_scale\_} conversions. Do not blindly copy those scale
factors into \texttt{RadiationIMC.cpp} unless the IMC class uses the
same nondimensional convention.

Use these rules:

\begin{enumerate}
\def\labelenumi{\arabic{enumi}.}
\tightlist
\item
  \texttt{particle.weight} is total radiation energy in the code's
  current energy units.
\item
  \texttt{cell.Eg{[}g{]}\ *\ cell.density} is group radiation energy
  density in the same units used by opacity and material coupling.
\item
  \texttt{grid.GetVolume(i)} is the volume in the same length units used
  by opacities.
\item
  \texttt{kappa}, \texttt{S}, \texttt{dSdUm}, \texttt{Kmat},
  \texttt{Hbase}, and \texttt{R} must all behave like opacities:
  multiplying by \texttt{c\ *\ dt} gives a dimensionless interaction
  strength.
\item
  \texttt{Btotal{[}g{]}} must be total energy emitted in the cell over
  the time step, so it includes \texttt{V\ *\ c\ *\ dt}.
\end{enumerate}

After implementation, add debug prints for one cell comparing:

\begin{Shaded}
\begin{Highlighting}[]
\NormalTok{sum\_g Bbase}\OperatorTok{[}\NormalTok{g}\OperatorTok{]}
\NormalTok{sum\_g Bcorr}\OperatorTok{[}\NormalTok{g}\OperatorTok{]}
\NormalTok{sum\_g Btotal}\OperatorTok{[}\NormalTok{g}\OperatorTok{]}
\NormalTok{V }\OperatorTok{*}\NormalTok{ c }\OperatorTok{*}\NormalTok{ dt }\OperatorTok{*}\NormalTok{ f }\OperatorTok{*}\NormalTok{ kappaP }\OperatorTok{*}\NormalTok{ Um}
\end{Highlighting}
\end{Shaded}

When Compton is disabled or \texttt{dSdUm\ =\ 0} and \texttt{S\ =\ 0},
\texttt{Bcorr} should be zero and \texttt{sum\_g\ Bbase{[}g{]}} should
match the old source energy from line 367.

\subsection{12. Parity checks against deterministic matrix
assembly}\label{parity-checks-against-deterministic-matrix-assembly}

Add a debug-only function:

\begin{Shaded}
\begin{Highlighting}[]
\DataTypeTok{void}\NormalTok{ RadiationIMC}\OperatorTok{::}\NormalTok{CheckComptonCoefficientParity}\OperatorTok{(}\DataTypeTok{size\_t}\NormalTok{ i}\OperatorTok{,} \DataTypeTok{double}\NormalTok{ fullDt}\OperatorTok{)} \AttributeTok{const}\OperatorTok{;}
\end{Highlighting}
\end{Shaded}

For a selected cell, reconstruct the matrix coefficient represented by
the IMC kernels.

The IMC representation has base effective scattering \texttt{Hbase},
residual \texttt{R}, and analog direct \texttt{Sanalog}:

\begin{Shaded}
\begin{Highlighting}[]
\NormalTok{K\^{}\{IMC\}\_\{h\textbackslash{}to g\}=H\^{}\{base\}\_\{h\textbackslash{}to g\}+R\_\{h\textbackslash{}to g\}+S\^{}\{analog\}\_\{h\textbackslash{}to g\}.}
\end{Highlighting}
\end{Shaded}

By construction this should equal

\begin{Shaded}
\begin{Highlighting}[]
\NormalTok{K\^{}\{target\}\_\{h\textbackslash{}to g\}=K\^{}\{mat\}\_\{h\textbackslash{}to g\}+S\_\{h\textbackslash{}to g\}.}
\end{Highlighting}
\end{Shaded}

Check:

\begin{Shaded}
\begin{Highlighting}[]
\ControlFlowTok{for}\NormalTok{ h}\OperatorTok{,}\NormalTok{g}\OperatorTok{:}
\NormalTok{    assertNear}\OperatorTok{(}\NormalTok{cd}\OperatorTok{.}\NormalTok{Hbase}\OperatorTok{[}\NormalTok{h}\OperatorTok{][}\NormalTok{g}\OperatorTok{]} \OperatorTok{+}\NormalTok{ cd}\OperatorTok{.}\NormalTok{R}\OperatorTok{[}\NormalTok{h}\OperatorTok{][}\NormalTok{g}\OperatorTok{]} \OperatorTok{+}\NormalTok{ cd}\OperatorTok{.}\NormalTok{Sanalog}\OperatorTok{[}\NormalTok{h}\OperatorTok{][}\NormalTok{g}\OperatorTok{],}
\NormalTok{               cd}\OperatorTok{.}\NormalTok{Kmat}\OperatorTok{[}\NormalTok{h}\OperatorTok{][}\NormalTok{g}\OperatorTok{]} \OperatorTok{+}\NormalTok{ cd}\OperatorTok{.}\NormalTok{S}\OperatorTok{[}\NormalTok{h}\OperatorTok{][}\NormalTok{g}\OperatorTok{]);}
\end{Highlighting}
\end{Shaded}

Also compare source:

\begin{Shaded}
\begin{Highlighting}[]
\NormalTok{Btotal}\OperatorTok{[}\NormalTok{g}\OperatorTok{]} \OperatorTok{==}\NormalTok{ Bbase}\OperatorTok{[}\NormalTok{g}\OperatorTok{]} \OperatorTok{+}\NormalTok{ Bcorr}\OperatorTok{[}\NormalTok{g}\OperatorTok{]}
\end{Highlighting}
\end{Shaded}

To compare directly to deterministic
\texttt{get\_implicit\_compton\_contribution}, multiply opacities by
\texttt{-V\ *\ c\ *\ dt} for left-hand-side matrix entries. The signs
should match these identities:

\begin{Shaded}
\begin{Highlighting}[]
\NormalTok{A\^{}\{implicit\}\_\{g,h\} = {-}V c\textbackslash{}Delta t\textbackslash{}left(H\^{}\{base\}\_\{h\textbackslash{}to g\}+R\_\{h\textbackslash{}to g\}+S\^{}\{analog\}\_\{h\textbackslash{}to g\}\textbackslash{}right).}
\end{Highlighting}
\end{Shaded}

For no Compton, \texttt{R\ =\ 0}, \texttt{Sanalog\ =\ 0}, and
\texttt{Hbase} should reduce to the ordinary IMC effective scattering
already represented by the current code.

\subsection{13. Testing plan}\label{testing-plan}

Implement tests in this order. Do not proceed to the next phase until
the previous phase passes.

\subsubsection{Test 1: no-Compton
regression}\label{test-1-no-compton-regression}

Settings:

\begin{itemize}
\tightlist
\item
  \texttt{withCompton=false}.
\item
  Same random seed as current baseline.
\end{itemize}

Expected:

\begin{itemize}
\tightlist
\item
  Results statistically match the old code.
\item
  If only signed-packet absolute-value changes were made, histories with
  positive packets should be nearly identical.
\end{itemize}

\subsubsection{Test 2: signed packet smoke
test}\label{test-2-signed-packet-smoke-test}

Manually create one positive and one negative packet in a static cell
with no Compton.

Expected:

\begin{itemize}
\tightlist
\item
  Positive packet absorption increases material energy.
\item
  Negative packet absorption decreases material energy.
\item
  \texttt{postStep} tallies the signed sum.
\item
  Low-weight removal does not immediately remove a negative packet
  because of a raw \texttt{\textless{}} comparison.
\end{itemize}

\subsubsection{Test 3: Compton coefficient construction with zero
Compton}\label{test-3-compton-coefficient-construction-with-zero-compton}

Force \texttt{S\ =\ 0} and \texttt{dSdUm\ =\ 0} after matrix
construction.

Expected:

\begin{itemize}
\tightlist
\item
  \texttt{Upsilon\ =\ 0}.
\item
  \texttt{Gamma\ =\ kappaP}.
\item
  \texttt{factorFleck} equals the old Fleck factor.
\item
  \texttt{Kmat\ =\ beta*c*dt*f*kappa\_g*b\_g*kappa\_h}.
\item
  \texttt{Hbase\ =\ Kmat} if \texttt{q\_g\ =\ kappa\_g*b\_g/kappaP} and
  \texttt{1-f\ =\ beta*c*dt*f*kappaP}.
\item
  \texttt{R\ =\ 0}.
\item
  \texttt{Bcorr\ =\ 0}.
\item
  \texttt{Btotal\ =\ Bbase}.
\end{itemize}

This is the most important algebraic check.

\subsubsection{Test 4: single-cell source-only Compton
correction}\label{test-4-single-cell-source-only-compton-correction}

Use one cell, no transport leakage, no physical scattering, and no
initial particles. Enable Compton coefficient construction but do not
move particles.

Expected:

\begin{itemize}
\tightlist
\item
  \texttt{generateComptonParticles} creates signed source packets whose
  signed group sums equal \texttt{Btotal{[}g{]}}.
\item
  Material internal energy changes by \texttt{-sum\_g\ Btotal{[}g{]}}.
\item
  Radiation plus material energy is unchanged immediately after source
  creation.
\end{itemize}

\subsubsection{Test 5: residual branch
expectation}\label{test-5-residual-branch-expectation}

Create one packet in group \texttt{h} with known weight \texttt{W}, no
absorption (\texttt{alpha=0}), no boundary crossing, and a fixed segment
\texttt{dt}.

Expected mean over many histories:

\begin{Shaded}
\begin{Highlighting}[]
\NormalTok{\textbackslash{}langle W\_\{child,g\}\textbackslash{}rangle}
\NormalTok{= W c\textbackslash{}Delta t\textbackslash{}left(R\^{}+\_\{h\textbackslash{}to g\}{-}R\^{}{-}\_\{h\textbackslash{}to g\}\textbackslash{}right).}
\end{Highlighting}
\end{Shaded}

Check positive and negative parent packets.

\subsubsection{Test 6: deterministic matrix
parity}\label{test-6-deterministic-matrix-parity}

For a static single-cell problem, compute the expected one-step
linearized update from the deterministic matrix coefficients and compare
to a high-particle IMC run.

Expected:

\begin{itemize}
\tightlist
\item
  Group means agree within Monte Carlo error.
\item
  Material energy exchange agrees within Monte Carlo error.
\end{itemize}

\subsubsection{Test 7: negative effective scattering
case}\label{test-7-negative-effective-scattering-case}

Choose a hot-electron or induced-scattering case known to produce
negative entries in \texttt{Kmat} or \texttt{R}.

Expected:

\begin{itemize}
\tightlist
\item
  No negative opacity is sampled.
\item
  \texttt{RnegRow{[}h{]}\ \textgreater{}\ 0} for at least one group.
\item
  Anti-packets are created.
\item
  The signed mean matches the deterministic linearized update.
\item
  Energy is conserved in a closed cell.
\end{itemize}

\subsubsection{Test 8: analog direct Compton
event}\label{test-8-analog-direct-compton-event}

After enabling \texttt{comptonUseAnalogDirect=true}, compare
\texttt{Sanalog} to \texttt{S}.

Expected:

\begin{itemize}
\tightlist
\item
  \texttt{Sres\ =\ S\ -\ Sanalog} is small except for known
  last-group/boundary terms.
\item
  MC direct Compton events produce the correct mean group transfer and
  material energy exchange.
\end{itemize}

\subsubsection{Test 9: MPI and boundary
particles}\label{test-9-mpi-and-boundary-particles}

Expected:

\begin{itemize}
\tightlist
\item
  MPI reductions for source allocation use absolute signed source
  energy.
\item
  Boundary particle generation remains positive and unchanged.
\item
  Signed residual packets crossing MPI boundaries are communicated like
  ordinary packets. If the communication code assumes positive weights,
  fix that assumption.
\end{itemize}

\subsubsection{Test 10: random walk disabled with
Compton}\label{test-10-random-walk-disabled-with-compton}

Expected:

\begin{itemize}
\tightlist
\item
  With \texttt{withCompton=true} and
  \texttt{comptonDisableRandomWalk=true}, random walk is not used.
\item
  With \texttt{withCompton=false}, random walk behavior is unchanged.
\end{itemize}

\subsection{14. Common implementation errors to
avoid}\label{common-implementation-errors-to-avoid}

\begin{enumerate}
\def\labelenumi{\arabic{enumi}.}
\item
  \textbf{Double-counting the thermal source.} If
  \texttt{generateComptonParticles} emits \texttt{Btotal}, do not also
  run the old line-367 source logic.
\item
  \textbf{Using the wrong target distribution for base effective
  scattering.} The residual subtracts \texttt{Hbase}. \texttt{Hbase}
  must use the exact same \texttt{q\_g} that the effective scattering
  event samples. Prefer an explicit \texttt{qBaseCdf} built from
  \texttt{kappa\_g*b\_g/kappaP}.
\item
  \textbf{Raw negative-weight removal.} Any check like
  \texttt{weight\ \textless{}\ threshold} will instantly delete
  anti-packets. Use \texttt{abs(weight)}.
\item
  \textbf{Forgetting to update material for residual child emission.}
  Every signed child source has a material energy change
  \texttt{-Wchild}.
\item
  \textbf{Clipping negative residuals to zero.} This changes the
  linearized physics. Split into \texttt{Rpos} and \texttt{Rneg}; do not
  discard \texttt{Rneg}.
\item
  \textbf{Letting a nonpositive Fleck denominator through.} The signed
  split fixes negative event probabilities, not an invalid
  linearization. If \texttt{1\ +\ beta*c*dt*Gamma\ \textless{}=\ 0},
  reject the step, reduce \texttt{dt}, or apply the deterministic
  limiter/fallback.
\item
  \textbf{Using new-state temperatures.} All Compton matrices,
  derivatives, Planck fractions, and opacities in this scheme are frozen
  at the beginning of the step.
\item
  \textbf{Random walk skipping residual sources.} Disable random walk
  for Compton-active cells until a residual track-length estimator is
  implemented inside random walk.
\item
  \textbf{Double-counting electron scattering.} If
  \texttt{CalcScatteringOpacity} already represents the same electron
  scattering that \texttt{tau} represents, do not also add analog
  Compton collision opacity without subtracting/replacing the elastic
  component.
\item
  \textbf{Treating direct Compton weight changes as nonconservative.} If
  a packet changes group energy by a ratio, update material by
  \texttt{Wold\ -\ Wnew}.
\end{enumerate}

\subsection{15. Recommended implementation sequence for a coding
agent}\label{recommended-implementation-sequence-for-a-coding-agent}

Follow this exact order.

\subsubsection{Commit 1: no-op refactor and signed
safety}\label{commit-1-no-op-refactor-and-signed-safety}

\begin{itemize}
\tightlist
\item
  Add \texttt{std::abs} packet removal.
\item
  Set \texttt{initialWeight\ =\ abs(weight)} everywhere.
\item
  Add helper \texttt{FrequencyForGroup} and group center arrays.
\item
  Verify no-Compton regression.
\end{itemize}

\subsubsection{Commit 2: Compton parameters and data
skeleton}\label{commit-2-compton-parameters-and-data-skeleton}

\begin{itemize}
\tightlist
\item
  Add parameters and members.
\item
  Add \texttt{ComptonCellData} with zero initialization.
\item
  Add \texttt{withCompton} to constructor and
  \texttt{operator\textless{}\textless{}}.
\item
  Do not change runtime behavior yet.
\end{itemize}

\subsubsection{Commit 3: Planck/group opacity
precompute}\label{commit-3-planckgroup-opacity-precompute}

\begin{itemize}
\tightlist
\item
  Add \texttt{precomputeComptonData} with only \texttt{kappa},
  \texttt{b}, \texttt{kappaP}, \texttt{qBase}.
\item
  When \texttt{withCompton=true}, compute modified source still
  disabled.
\item
  Print/check \texttt{kappaP} versus \texttt{CalcPlanckOpacity}.
\end{itemize}

\subsubsection{Commit 4: Port Compton matrix
generation}\label{commit-4-port-compton-matrix-generation}

\begin{itemize}
\tightlist
\item
  Add \texttt{compton\_matrix\_gen} member.
\item
  Port \texttt{generate\_S\_and\_dSdUm\_matrices} into
  \texttt{BuildComptonMatricesForCell}.
\item
  Add \texttt{D}, \texttt{Upsilon}, \texttt{M}, \texttt{Lambda}, and
  \texttt{Gamma} computations.
\item
  Add n=0 fallback.
\item
  Add tests that print matrices for one cell and compare against
  deterministic code.
\end{itemize}

\subsubsection{Commit 5: Modified Fleck
factor}\label{commit-5-modified-fleck-factor}

\begin{itemize}
\tightlist
\item
  Replace \texttt{factorFleck{[}i{]}} by Compton-modified \texttt{f}
  when \texttt{withCompton=true}.
\item
  Keep old Fleck when \texttt{withCompton=false}.
\item
  Run zero-Compton test: \texttt{S=dSdUm=0} gives old Fleck.
\end{itemize}

\subsubsection{Commit 6: Group-signed source
generation}\label{commit-6-group-signed-source-generation}

\begin{itemize}
\tightlist
\item
  Add \texttt{BuildComptonSources}.
\item
  Add \texttt{generateComptonParticles}.
\item
  Change \texttt{generateParticles} to call it only when
  \texttt{withCompton=true}.
\item
  Verify source-only conservation.
\end{itemize}

\subsubsection{Commit 7: Residual kernels and anti-packet branch
source}\label{commit-7-residual-kernels-and-anti-packet-branch-source}

\begin{itemize}
\tightlist
\item
  Add \texttt{BuildComptonResidualKernels} with
  \texttt{comptonUseAnalogDirect=false}.
\item
  Add \texttt{EmitComptonResidualChildren} and
  \texttt{BankSignedChildPacket}.
\item
  Use \texttt{particlesToAdd} in \texttt{step}.
\item
  Disable random walk for Compton.
\item
  Run residual branch expectation test.
\end{itemize}

\subsubsection{Commit 8: Deterministic parity
test}\label{commit-8-deterministic-parity-test}

\begin{itemize}
\tightlist
\item
  Add debug parity function.
\item
  Verify \texttt{Hbase\ +\ R\ +\ Sanalog\ =\ Kmat\ +\ S}.
\item
  Verify source equals deterministic \texttt{Bbase+Bcorr}.
\end{itemize}

\subsubsection{Commit 9: Optional physical analog
Compton}\label{commit-9-optional-physical-analog-compton}

\begin{itemize}
\tightlist
\item
  Add \texttt{BuildAnalogDirectComptonMatrixAndCdfs}.
\item
  Add \texttt{directComptonOpacity} event channel.
\item
  Add \texttt{ApplyAnalogComptonEvent}.
\item
  Keep \texttt{Sres\ =\ S\ -\ Sanalog} so exactness is preserved.
\item
  Test direct Compton energy exchange.
\end{itemize}

\subsubsection{Commit 10: variance
controls}\label{commit-10-variance-controls}

\begin{itemize}
\tightlist
\item
  Add residual child roulette.
\item
  Add census cancellation.
\item
  Add diagnostics for negative group radiation from Monte Carlo noise.
\end{itemize}

\subsection{16. Minimal acceptance
criteria}\label{minimal-acceptance-criteria}

The implementation is acceptable when all of the following are true:

\begin{enumerate}
\def\labelenumi{\arabic{enumi}.}
\tightlist
\item
  \texttt{withCompton=false} reproduces the old IMC behavior.
\item
  Compton coefficients are frozen in \texttt{preStep} and not updated
  during transport.
\item
  The modified Fleck factor uses \texttt{Gamma\ =\ kappaP\ +\ Upsilon}.
\item
  Source packets represent \texttt{Btotal\_g\ =\ Bbase\_g\ +\ Bcorr\_g}
  by group and sign.
\item
  The existing base effective scattering is represented by
  \texttt{Hbase} and is subtracted from the residual, so it is not
  double-counted.
\item
  The residual matrix satisfies \texttt{R\ =\ Kmat\ -\ Hbase\ +\ Sres}.
\item
  \texttt{Rpos} and \texttt{Rneg} are never sampled as negative
  probabilities; negative parts create anti-packets or signed residual
  packets.
\item
  \texttt{particlesToAdd} is used for residual children.
\item
  Every signed packet creation or destruction updates material energy
  conservatively.
\item
  Random walk is disabled or extended so that residual sources are not
  skipped.
\item
  A debug parity test confirms that the IMC kernels reconstruct the
  deterministic linearized matrix.
\item
  No source iteration or nonlinear iteration is introduced.
\end{enumerate}

\subsection{17. Final notes on the two valid implementation
modes}\label{final-notes-on-the-two-valid-implementation-modes}

There are two valid ways to implement direct Compton in this framework.

\subsubsection{Mode 1: residual-only direct
Compton}\label{mode-1-residual-only-direct-compton}

Set \texttt{comptonUseAnalogDirect=false}.

Then

\begin{Shaded}
\begin{Highlighting}[]
\NormalTok{S\^{}\{analog\}=0,}
\NormalTok{\textbackslash{}qquad}
\NormalTok{S\^{}\{res\}=S,}
\NormalTok{\textbackslash{}qquad}
\NormalTok{R=K\^{}\{mat\}{-}H\^{}\{base\}+S.}
\end{Highlighting}
\end{Shaded}

Advantages:

\begin{itemize}
\tightlist
\item
  Fastest to code.
\item
  Algebraically exact for the linearized multigroup equation in
  expectation.
\item
  No risk of double-counting \texttt{CalcScatteringOpacity}.
\end{itemize}

Disadvantages:

\begin{itemize}
\tightlist
\item
  Direct Compton is represented through signed branch sources rather
  than physical collisions.
\item
  May have higher variance.
\end{itemize}

\subsubsection{Mode 2: physical analog direct Compton plus residual
correction}\label{mode-2-physical-analog-direct-compton-plus-residual-correction}

Set \texttt{comptonUseAnalogDirect=true}.

Then

\begin{Shaded}
\begin{Highlighting}[]
\NormalTok{S\^{}\{res\}=S{-}S\^{}\{analog\},}
\NormalTok{\textbackslash{}qquad}
\NormalTok{R=K\^{}\{mat\}{-}H\^{}\{base\}+S\^{}\{res\}.}
\end{Highlighting}
\end{Shaded}

Advantages:

\begin{itemize}
\tightlist
\item
  Physical Compton group changes are sampled as positive events.
\item
  Lower variance when \texttt{Sanalog} closely matches \texttt{S}.
\item
  The residual still guarantees exact parity with the deterministic
  linearized operator.
\end{itemize}

Disadvantages:

\begin{itemize}
\tightlist
\item
  Requires careful event-rate and material energy-exchange
  implementation.
\item
  Requires avoiding double counting with existing elastic/electron
  scattering opacity.
\end{itemize}

Recommended path: implement Mode 1 first, verify the implicit scheme,
then add Mode 2 as an optimization.

\end{document}
